\input{preamble}

% OK, start here.
%
\begin{document}

\title{Méthodes simpliciales}


\maketitle

\phantomsection
\label{section-phantom}

\tableofcontents

\section{Introduction}
\label{section-introduction}

\noindent
Voici une introduction minimale aux méthodes simpliciales.
Nous y ajoutons simplement les éléments dont nous aurons besoin par la suite.
On pourra consulter \cite{SimpHom} comme référence générale sur ce sujet.
Le mémoire de Quillen sur l'algèbre homotopique, voir \cite{Quillen},
ou celui d'Artin et Mazur sur l'homotopie étale, voir \cite{ArtinMazur},
donnent des exemples de ce que permettent ces méthodes.

\section{La catégorie des ensembles finis ordonnés}
\label{section-Delta}

\noindent
La catégorie $\Delta$ est la catégorie dont
\begin{enumerate}
\item les objets sont $[0], [1], [2], \ldots$, où
$[n] = \{0, 1, 2, \ldots, n\}$, et
\item un morphisme $[n] \to [m]$ est une application croissante au sens large
$\{0, 1, 2, \ldots, n\} \to \{0, 1, 2, \ldots, m\}$
entre les ensembles correspondants.
\end{enumerate}
Par définition, une application $\varphi : [n] \to [m]$ est
{\it croissante au sens large} si $\varphi(i) \geq \varphi(j)$ dès que
$i \geq j$. Autrement dit, $\Delta$ est une catégorie équivalente à la
« grande » catégorie des ensembles finis totalement ordonnés non vides et
des applications croissantes au sens large.
Il existe exactement $n + 1$ morphismes $[0] \to [n]$ et exactement $1$
morphisme $[n] \to [0]$. Il existe exactement $(n + 1)(n + 2)/2$
morphismes $[1] \to [n]$ et exactement $n + 2$ morphismes
$[n] \to [1]$. Et ainsi de suite.

\begin{definition}
\label{definition-face-degeneracy}
Pour tout entier $n\geq 1$ et tout $0\leq j \leq n$, on note
{\it $\delta^n_j : [n-1] \to [n]$}
l'application injective croissante qui omet $j$. Pour tout entier
$n\geq 0$ et tout $0\leq j \leq n$, on note
{\it $\sigma^n_j : [n + 1]  \to [n]$}
l'application surjective croissante telle que
$(\sigma^n_j)^{-1}(\{j\}) = \{j, j + 1\}$.
\end{definition}

\begin{lemma}
\label{lemma-face-degeneracy}
Tout morphisme de $\Delta$ s'écrit comme une composée de morphismes
$\delta^n_j$ et $\sigma^n_j$.
\end{lemma}

\begin{proof}
Soit $\varphi : [n] \to [m]$ un morphisme de $\Delta$.
Si $j \not \in \Im(\varphi)$, on peut écrire
$\varphi$ sous la forme $\delta^m_j \circ \psi$ pour un morphisme
$\psi : [n] \to [m - 1]$. Si $\varphi(j) = \varphi(j + 1)$,
on peut écrire $\varphi$ sous la forme $\psi \circ \sigma^{n - 1}_j$
pour un morphisme $\psi : [n - 1] \to [m]$.
Chaque remplacement de ce type diminue $n + m$ ; on aboutit donc à un
$\varphi$ à la fois injectif et surjectif, donc à un morphisme identité.
\end{proof}

\begin{lemma}
\label{lemma-relations-face-degeneracy}
Les morphismes $\delta^n_j$ et $\sigma^n_j$ satisfont aux relations suivantes.
\begin{enumerate}
\item Si $0 \leq i < j \leq n + 1$, alors
$\delta^{n + 1}_j \circ \delta^n_i =
\delta^{n + 1}_i \circ \delta^n_{j - 1}$.
Autrement dit, le diagramme
$$
\xymatrix{
& [n] \ar[rd]^{\delta^{n + 1}_j} & \\
[n - 1] \ar[ru]^{\delta^n_i} \ar[rd]_{\delta^n_{j - 1}} & &
[n + 1] \\
& [n] \ar[ru]_{\delta^{n + 1}_i} &
}
$$
est commutatif.
\item Si $0 \leq i < j \leq n - 1$, alors
$\sigma^{n - 1}_j \circ \delta^n_i =
\delta^{n - 1}_i \circ \sigma^{n - 2}_{j - 1}$.
Autrement dit, le diagramme
$$
\xymatrix{
& [n] \ar[rd]^{\sigma^{n - 1}_j} & \\
[n - 1] \ar[ru]^{\delta^n_i} \ar[rd]_{\sigma^{n - 2}_{j - 1}} & &
[n - 1] \\
& [n - 2] \ar[ru]_{\delta^{n - 1}_i} &
}
$$
est commutatif.
\item Si $0 \leq j \leq n - 1$, alors
$\sigma^{n - 1}_j \circ \delta^n_j = \text{id}_{[n - 1]}$
and
$\sigma^{n - 1}_j \circ \delta^n_{j + 1} = \text{id}_{[n - 1]}$.
Autrement dit, le diagramme
$$
\xymatrix{
& [n] \ar[rd]^{\sigma^{n - 1}_j} & \\
[n - 1]
\ar[ru]^{\delta^n_j}
\ar[rd]_{\delta^n_{j + 1}}
\ar[rr]^{\text{id}_{[n - 1]}} & & [n - 1] \\
& [n] \ar[ru]_{\sigma^{n - 1}_j} &
}
$$
est commutatif.
\item Si $0 < j + 1 < i \leq n$, alors
$\sigma^{n - 1}_j \circ \delta^n_i =
\delta^{n - 1}_{i - 1} \circ \sigma^{n - 2}_j$.
Autrement dit, le diagramme
$$
\xymatrix{
& [n] \ar[rd]^{\sigma^{n - 1}_j} & \\
[n - 1] \ar[ru]^{\delta^n_i} \ar[rd]_{\sigma^{n - 2}_j} & &
[n - 1] \\
& [n - 2] \ar[ru]_{\delta^{n - 1}_{i - 1}} &
}
$$
est commutatif.
\item Si $0 \leq i \leq j \leq n - 1$, alors
$\sigma^{n - 1}_j \circ \sigma^n_i =
\sigma^{n - 1}_i \circ \sigma^n_{j + 1}$.
Autrement dit, le diagramme
$$
\xymatrix{
& [n] \ar[rd]^{\sigma^{n - 1}_j} & \\
[n + 1] \ar[ru]^{\sigma^n_i} \ar[rd]_{\sigma^n_{j + 1}} & &
[n - 1] \\
& [n] \ar[ru]_{\sigma^{n - 1}_i} &
}
$$
est commutatif.
\end{enumerate}
\end{lemma}

\begin{proof}
Omis.
\end{proof}

\begin{lemma}
\label{lemma-face-degeneracy-category}
La catégorie $\Delta$ est la catégorie universelle ayant pour objets les
$[n]$, $n \geq 0$, et pour morphismes les $\delta^n_j$ et $\sigma^n_j$,
telle que (a) tout morphisme soit une composée de ces morphismes,
(b) les relations énumérées dans le lemme
\ref{lemma-relations-face-degeneracy} soient satisfaites et
(c) toute relation entre les morphismes soit une conséquence de celles-ci.
\end{lemma}

\begin{proof}
Omis.
\end{proof}







\section{Objets simpliciaux}
\label{section-simplicial-object}

\begin{definition}
\label{definition-simplicial-object}
Soit $\mathcal{C}$ une catégorie.
\begin{enumerate}
\item Un {\it objet simplicial $U$ de $\mathcal{C}$}
est un foncteur contravariant $U$ de $\Delta$ vers
$\mathcal{C}$, c'est-à-dire un foncteur
$$
U : \Delta^{opp} \longrightarrow \mathcal{C}
$$
\item Si $\mathcal{C}$ est la catégorie des ensembles, on appelle
$U$ un {\it ensemble simplicial}.
\item Si $\mathcal{C}$ est la catégorie des groupes abéliens,
on appelle $U$ un {\it groupe abélien simplicial}.
\item Un {\it morphisme d'objets simpliciaux $U \to U'$}
est une transformation de foncteurs.
\item La {\it catégorie des objets simpliciaux de $\mathcal{C}$}
est notée $\text{Simp}(\mathcal{C})$.
\end{enumerate}
\end{definition}

\noindent
Cela signifie que l'on dispose d'objets $U([0]), U([1]), U([2]), \ldots$
et, pour $\varphi$ une application croissante au sens large
$\varphi : [m] \to [n]$, d'un morphisme
$U(\varphi) : U([n]) \to U([m])$, satisfaisant
$U(\varphi \circ \psi) = U(\psi) \circ U(\varphi)$.

\medskip\noindent
En particulier, il existe un unique morphisme $U([0]) \to U([n])$ et
exactement $n + 1$ morphismes $U([n]) \to U([0])$, correspondant aux
$n + 1$ applications $[0] \to [n]$. Pour parler commodément de ces
objets simpliciaux, il nous faut davantage de notations. Nous utilisons
pour cela les morphismes distingués à la section \ref{section-Delta}.

\begin{lemma}
\label{lemma-characterize-simplicial-object}
Soit $\mathcal{C}$ une catégorie.
\begin{enumerate}
\item Étant donné un objet simplicial $U$ de $\mathcal{C}$,
on obtient une suite d'objets $U_n = U([n])$ munis des morphismes
$d^n_j = U(\delta^n_j) : U_n \to U_{n-1}$ et
$s^n_j = U(\sigma^n_j) : U_n \to U_{n + 1}$. Ces morphismes
satisfont aux relations opposées à celles du lemme
\ref{lemma-relations-face-degeneracy}, à savoir :
\begin{enumerate}
\item Si $0 \leq i < j \leq n + 1$, alors
$d^n_i \circ d^{n + 1}_j = d^n_{j - 1} \circ d^{n + 1}_i$.
\item Si $0 \leq i < j \leq n - 1$, alors
$d^n_i \circ s^{n - 1}_j = s^{n - 2}_{j - 1} \circ d^{n - 1}_i$.
\item Si $0 \leq j \leq n - 1$, alors
$\text{id} = d^n_j \circ s^{n - 1}_j = d^n_{j + 1} \circ s^{n - 1}_j$.
\item Si $0 < j + 1 < i \leq n$, alors
$d^n_i \circ s^{n - 1}_j = s^{n - 2}_j \circ d^{n - 1}_{i - 1}$.
\item Si $0 \leq i \leq j \leq n - 1$, alors
$s^n_i \circ s^{n - 1}_j = s^n_{j + 1} \circ s^{n - 1}_i$.
\end{enumerate}
\item Réciproquement, étant donnés une suite d'objets $U_n$ et des
morphismes $d^n_j$, $s^n_j$ satisfaisant à (1)(a) -- (e), il existe un
unique objet simplicial $U$ de $\mathcal{C}$ tel que $U_n = U([n])$,
$d^n_j = U(\delta^n_j)$ et $s^n_j = U(\sigma^n_j)$.
\item Un morphisme entre objets simpliciaux $U$ et $U'$ est donné par une
famille de morphismes $U_n \to U'_n$ qui commutent aux morphismes
$d^n_j$ et $s^n_j$.
\end{enumerate}
\end{lemma}

\begin{proof}
Cela résulte du lemme \ref{lemma-face-degeneracy-category}.
\end{proof}

\begin{remark}
\label{remark-relations}
Par abus de notation, nous écrivons parfois $d_i : U_n \to U_{n - 1}$
au lieu de $d^n_i$, et de même $s_i : U_n \to U_{n + 1}$.
Les relations entre les morphismes $d^n_i$ et $s^n_i$
peuvent alors s'écrire comme suit :
\begin{enumerate}
\item Si $i < j$, alors $d_i \circ d_j = d_{j - 1} \circ d_i$.
\item Si $i < j$, alors $d_i \circ s_j = s_{j - 1} \circ d_i$.
\item On a $\text{id} = d_j \circ s_j = d_{j + 1} \circ s_j$.
\item Si $i > j + 1$, alors $d_i \circ s_j = s_j \circ d_{i - 1}$.
\item Si $i \leq j$, alors $s_i \circ s_j = s_{j + 1} \circ s_i$.
\end{enumerate}
Cela signifie que, dès que les composées des deux membres sont définies,
l'égalité correspondante est satisfaite.
\end{remark}

\noindent
On obtient un unique morphisme $s^0_0 = U(\sigma^0_0) : U_0 \to U_1$ et
deux morphismes $d^1_0 = U(\delta^1_0)$ et
$d^1_1 = U(\delta^1_1)$, qui sont des morphismes $U_1 \to U_0$.
Il existe deux morphismes $s^1_0 = U(\sigma^1_0)$,
$s^1_1 = U(\sigma^1_1)$, qui sont des morphismes $U_1 \to U_2$,
et trois morphismes
$d^2_0 = U(\delta^2_0)$, $d^2_1 = U(\delta^2_1)$,
$d^2_2 = U(\delta^2_2)$, qui sont des morphismes $U_2 \to U_1$.
Et ainsi de suite.

\medskip\noindent
On se représente $U$ par le diagramme suivant :
$$
\xymatrix{
U_2
\ar@<2ex>[r]
\ar@<0ex>[r]
\ar@<-2ex>[r]
&
U_1
\ar@<1ex>[r]
\ar@<-1ex>[r]
\ar@<1ex>[l]
\ar@<-1ex>[l]
&
U_0
\ar@<0ex>[l]
}
$$
Les morphismes $d$ sont les flèches dirigées vers la droite et les
morphismes $s$ celles dirigées vers la gauche.

\begin{example}
\label{example-constant-simplicial-object}
L'exemple le plus simple est l'objet simplicial {\it constant} de valeur
$X \in \Ob(\mathcal{C})$. Autrement dit, $U_n = X$ et toutes les
applications sont $\text{id}_X$.
\end{example}

\begin{example}
\label{example-fibre-products-simplicial-object}
Supposons que $Y\to X$ soit un morphisme de $\mathcal{C}$ tel que tous les
produits fibrés $Y \times_X Y \times_X \ldots \times_X Y$ existent.
On prend alors pour $U_n$ le produit fibré à $(n + 1)$ facteurs et l'on
associe à $\varphi : [n] \to [m]$ l'application (en « coordonnées »)
$(y_0, \ldots, y_m) \mapsto (y_{\varphi(0)}, \ldots, y_{\varphi(n)})$.
Autrement dit, l'application $U_0 = Y \to U_1 = Y \times_X Y$ est
l'application diagonale, tandis que les deux applications
$U_1 = Y \times_X Y \to U_0 = Y$ sont les projections.
\end{example}

\noindent
L'exemple \ref{example-fibre-products-simplicial-object} est important
du point de vue géométrique. Il invite à considérer les applications
$d^n_j : U_n \to U_{n - 1}$ comme des projections (qui oublient la
$j$-ième composante) et les applications $s^n_j : U_n \to U_{n + 1}$
comme des applications diagonales (qui répètent la $j$-ième coordonnée).
Nous y reviendrons dans les sections suivantes.

\begin{lemma}
\label{lemma-si-injective}
Soit $\mathcal{C}$ une catégorie et soit $U$ un objet simplicial de
$\mathcal{C}$. Chacun des morphismes $s^n_i : U_n \to U_{n + 1}$
admet un inverse à gauche. En particulier, $s^n_i$ est un monomorphisme.
\end{lemma}

\begin{proof}
En effet, $d_i^{n + 1} \circ s^n_i = \text{id}_{U_n}$.
\end{proof}

\section{Objets simpliciaux comme préfaisceaux}
\label{section-simplicial-presheaves}

\noindent
On peut encore considérer un objet simplicial de $\mathcal{C}$ comme un
préfaisceau à valeurs dans $\mathcal{C}$ sur $\Delta$. Voir
Sites, définition \ref{sites-definition-presheaf}.
En effet, si $U$ et $U'$ sont des objets simpliciaux de $\mathcal{C}$,
alors
\begin{equation}
\label{equation-simplicial-set-presheaf}
\Mor(U, U') = \Mor_{\textit{PSh}(\Delta)}(U, U').
\end{equation}
Une partie de ce qui suit pourrait être remplacée par les constructions
plus générales du chapitre sur les sites. Les arguments propres aux objets
simpliciaux donnent toutefois une présentation plus transparente.
















\section{Objets cosimpliciaux}
\label{section-cosimplicial-object}

\noindent
On pourrait définir un objet cosimplicial d'une catégorie $\mathcal{C}$
comme un objet simplicial de la catégorie opposée $\mathcal{C}^{opp}$.
Cette façon de penser étant peu naturelle, nous les introduisons séparément
et en relevons quelques propriétés élémentaires.

\begin{definition}
\label{definition-cosimplicial-object}
Soit $\mathcal{C}$ une catégorie.
\begin{enumerate}
\item Un {\it objet cosimplicial $U$ de $\mathcal{C}$}
est un foncteur covariant $U$ de $\Delta$ vers
$\mathcal{C}$, c'est-à-dire un foncteur
$$
U : \Delta \longrightarrow \mathcal{C}
$$
\item Si $\mathcal{C}$ est la catégorie des ensembles, on appelle
$U$ un {\it ensemble cosimplicial}.
\item Si $\mathcal{C}$ est la catégorie des groupes abéliens,
on appelle $U$ un {\it groupe abélien cosimplicial}.
\item Un {\it morphisme d'objets cosimpliciaux $U \to U'$}
est une transformation de foncteurs.
\item La {\it catégorie des objets cosimpliciaux de $\mathcal{C}$}
est notée $\text{CoSimp}(\mathcal{C})$.
\end{enumerate}
\end{definition}

\noindent
Cela signifie que l'on dispose d'objets $U([0]), U([1]), U([2]), \ldots$
et, pour $\varphi$ une application croissante au sens large
$\varphi : [m] \to [n]$, d'un morphisme
$U(\varphi) : U([m]) \to U([n])$, satisfaisant
$U(\varphi \circ \psi) = U(\varphi) \circ U(\psi)$.

\medskip\noindent
En particulier, il existe un unique morphisme $U([n]) \to U([0])$ et
exactement $n + 1$ morphismes $U([0]) \to U([n])$, correspondant aux
$n + 1$ applications $[0] \to [n]$. Pour parler commodément de ces
objets cosimpliciaux, nous utilisons les morphismes distingués à la
section \ref{section-Delta}.

\begin{lemma}
\label{lemma-characterize-cosimplicial-object}
Soit $\mathcal{C}$ une catégorie.
\begin{enumerate}
\item Étant donné un objet cosimplicial $U$ de $\mathcal{C}$, on obtient
une suite d'objets $U_n = U([n])$ munis des morphismes
$\delta^n_j = U(\delta^n_j) : U_{n - 1} \to U_n$ et
$\sigma^n_j = U(\sigma^n_j) : U_{n + 1} \to U_n$. Ces morphismes
satisfont aux relations du lemme \ref{lemma-relations-face-degeneracy}.
\item Réciproquement, étant donnés une suite d'objets $U_n$ et des
morphismes $\delta^n_j$, $\sigma^n_j$ satisfaisant à ces relations, il
existe un unique objet cosimplicial $U$ de $\mathcal{C}$ tel que
$U_n = U([n])$, $\delta^n_j = U(\delta^n_j)$ et
$\sigma^n_j = U(\sigma^n_j)$.
\item Un morphisme entre objets cosimpliciaux $U$ et $U'$ est donné par
une famille de morphismes $U_n \to U'_n$ qui commutent aux morphismes
$\delta^n_j$ et $\sigma^n_j$.
\end{enumerate}
\end{lemma}

\begin{proof}
Cela résulte du lemme \ref{lemma-face-degeneracy-category}.
\end{proof}

\begin{remark}
\label{remark-relations-cosimplicial}
Par abus de notation, nous écrivons parfois
$\delta_i : U_{n - 1} \to U_n$ au lieu de $\delta^n_i$, et de même
$\sigma_i : U_{n + 1} \to U_n$. Les relations entre les morphismes
$\delta^n_i$ et $\sigma^n_i$ peuvent alors s'écrire comme suit :
\begin{enumerate}
\item Si $i < j$, alors
$\delta_j \circ \delta_i = \delta_i \circ \delta_{j - 1}$.
\item Si $i < j$, alors
$\sigma_j \circ \delta_i = \delta_i \circ \sigma_{j - 1}$.
\item On a
$\text{id} = \sigma_j \circ \delta_j = \sigma_j \circ \delta_{j + 1}$.
\item Si $i > j + 1$, alors
$\sigma_j \circ \delta_i = \delta_{i - 1} \circ \sigma_j$.
\item Si $i \leq j$, alors
$\sigma_j \circ \sigma_i = \sigma_i \circ \sigma_{j + 1}$.
\end{enumerate}
Cela signifie que, dès que les composées des deux membres sont définies,
l'égalité correspondante est satisfaite.
\end{remark}

\noindent
On obtient un unique morphisme
$\sigma^0_0 = U(\sigma^0_0) : U_1 \to U_0$ et deux morphismes
$\delta^1_0 = U(\delta^1_0)$ et $\delta^1_1 = U(\delta^1_1)$,
qui sont des morphismes $U_0 \to U_1$. Il existe deux morphismes
$\sigma^1_0 = U(\sigma^1_0)$, $\sigma^1_1 = U(\sigma^1_1)$,
qui sont des morphismes $U_2 \to U_1$, et trois morphismes
$\delta^2_0 = U(\delta^2_0)$, $\delta^2_1 = U(\delta^2_1)$,
$\delta^2_2 = U(\delta^2_2)$, qui sont des morphismes $U_2 \to U_3$.
Et ainsi de suite.

\medskip\noindent
On se représente $U$ par le diagramme suivant :
$$
\xymatrix{
U_0
\ar@<1ex>[r]
\ar@<-1ex>[r]
&
U_1
\ar@<0ex>[l]
\ar@<2ex>[r]
\ar@<0ex>[r]
\ar@<-2ex>[r]
&
U_2
\ar@<1ex>[l]
\ar@<-1ex>[l]
}
$$
Les morphismes $\delta$ sont les flèches dirigées vers la droite et les
morphismes $\sigma$ celles dirigées vers la gauche.

\begin{example}
\label{example-constant-cosimplicial-object}
L'exemple le plus simple est l'objet cosimplicial {\it constant} de valeur
$X \in \Ob(\mathcal{C})$. Autrement dit, $U_n = X$ et toutes les
applications sont $\text{id}_X$.
\end{example}

\begin{example}
\label{example-push-outs-simplicial-object}
Supposons que $X\to Y$ soit un morphisme de $\mathcal{C}$ tel que toutes les
sommes amalgamées $Y\amalg_X Y \amalg_X \ldots \amalg_X Y$ existent.
On prend alors pour $U_n$ la somme amalgamée à $(n + 1)$ facteurs et
l'on associe à $\varphi : [n] \to [m]$ l'application
$$
(y \text{ dans la }i\text{-ième composante})
\mapsto
(y \text{ dans la }\varphi(i)\text{-ième composante})
$$
en « coordonnées ». Autrement dit, l'application
$U_1 = Y \amalg_X Y \to U_0 = Y$ est l'identité sur chaque composante.
Les deux applications $U_0 = Y \to U_1 = Y \amalg_X Y$ sont les deux
coprojections.
\end{example}

\begin{example}
\label{example-simplex-cosimplicial-set}
Pour tout $n \geq 0$, on note $C[n]$ l'ensemble cosimplicial
$$
\Delta \longrightarrow \textit{Sets},\quad
[k] \longmapsto \Mor_{\Delta}([n], [k])
$$
Cet exemple est dual de l'exemple \ref{example-simplex-simplicial-set}.
\end{example}

\begin{lemma}
\label{lemma-di-injective}
Soit $\mathcal{C}$ une catégorie et soit $U$ un objet cosimplicial de
$\mathcal{C}$. Chacun des morphismes $\delta^n_i : U_{n - 1} \to U_n$
admet un inverse à gauche. En particulier, $\delta^n_i$ est un
monomorphisme.
\end{lemma}

\begin{proof}
En effet,
$\sigma_i^{n - 1} \circ \delta^n_i = \text{id}_{U_{n - 1}}$
pour $i < n$ ; lorsque $i = n$, on utilise
$\sigma_{n - 1}^{n - 1} \circ \delta^n_n = \text{id}_{U_{n - 1}}$.
\end{proof}


























\section{Produits d'objets simpliciaux}
\label{section-products}

\noindent
Il serait naturel de définir le produit d'objets simpliciaux comme leur
produit dans la catégorie des objets simpliciaux. Il pourrait alors arriver,
ce qui prêterait à confusion, que ce produit existe sans être décrit comme
ci-dessous. Nous le définissons donc directement.

\begin{definition}
\label{definition-product}
Soit $\mathcal{C}$ une catégorie. Soient $U$ et $V$ des objets simpliciaux
de $\mathcal{C}$. Supposons que les produits $U_n \times V_n$ existent
dans $\mathcal{C}$. Le {\it produit de $U$ et $V$} est l'objet simplicial
$U \times V$ défini comme suit :
\begin{enumerate}
\item $(U \times V)_n = U_n \times V_n$,
\item $d^n_i = (d^n_i, d^n_i)$, and
\item $s^n_i = (s^n_i, s^n_i)$.
\end{enumerate}
Autrement dit, $U \times V$ est le produit des préfaisceaux $U$ et $V$
sur $\Delta$.
\end{definition}

\begin{lemma}
\label{lemma-product}
Si $U$ et $V$ sont des objets simpliciaux de la catégorie $\mathcal{C}$
et si $U \times V$ existe, alors
$$
\Mor(W, U \times V) =
\Mor(W, U) \times
\Mor(W, V)
$$
pour tout troisième objet simplicial $W$ de $\mathcal{C}$.
\end{lemma}

\begin{proof}
Omis.
\end{proof}


\section{Produits fibrés d'objets simpliciaux}
\label{section-fibre-products}

\noindent
Il serait naturel de définir le produit fibré d'objets simpliciaux comme
leur produit fibré dans la catégorie des objets simpliciaux. Il pourrait
alors arriver, ce qui prêterait à confusion, que ce produit fibré existe
sans être décrit comme ci-dessous. Nous le définissons donc directement.

\begin{definition}
\label{definition-fibre-product}
Soit $\mathcal{C}$ une catégorie. Soient $U, V, W$ des objets simpliciaux
de $\mathcal{C}$ et soient $a : V \to U$, $b : W \to U$ des morphismes.
Supposons que les produits fibrés $V_n \times_{U_n} W_n$ existent dans
$\mathcal{C}$. Le {\it produit fibré de $V$ et $W$ au-dessus de $U$}
est l'objet simplicial $V \times_U W$ défini comme suit :
\begin{enumerate}
\item $(V \times_U W)_n = V_n \times_{U_n} W_n$,
\item $d^n_i = (d^n_i, d^n_i)$, and
\item $s^n_i = (s^n_i, s^n_i)$.
\end{enumerate}
Autrement dit, $V \times_U W$ est le produit fibré des préfaisceaux
$V$ et $W$ au-dessus du préfaisceau $U$ sur $\Delta$.
\end{definition}

\begin{lemma}
\label{lemma-fibre-product}
Si $U, V, W$ sont des objets simpliciaux de la catégorie $\mathcal{C}$,
si $a : V \to U$, $b : W \to U$ sont des morphismes et si
$V \times_U W$ existe, alors
$$
\Mor(T, V \times_U W) =
\Mor(T, V) \times_{\Mor(T, U)}
\Mor(T, W)
$$
pour tout quatrième objet simplicial $T$ de $\mathcal{C}$.
\end{lemma}

\begin{proof}
Omis.
\end{proof}

\section{Sommes amalgamées d'objets simpliciaux}
\label{section-push-outs}

\noindent
Il serait naturel de définir la somme amalgamée d'objets simpliciaux comme
leur somme amalgamée dans la catégorie des objets simpliciaux. Il pourrait
alors arriver, ce qui prêterait à confusion, que ces sommes existent sans
être décrites comme ci-dessous. Nous les définissons donc directement.

\begin{definition}
\label{definition-push-out}
Soit $\mathcal{C}$ une catégorie. Soient $U, V, W$ des objets simpliciaux
de $\mathcal{C}$ et soient $a : U \to V$, $b : U \to W$ des morphismes.
Supposons que les sommes amalgamées $V_n \amalg_{U_n} W_n$ existent dans
$\mathcal{C}$. La {\it somme amalgamée de $V$ et $W$ sous $U$} est
l'objet simplicial $V\amalg_U W$ défini comme suit :
\begin{enumerate}
\item $(V \amalg_U W)_n = V_n \amalg_{U_n} W_n$,
\item $d^n_i = (d^n_i, d^n_i)$, and
\item $s^n_i = (s^n_i, s^n_i)$.
\end{enumerate}
Autrement dit, $V\amalg_U W$ est la somme amalgamée des préfaisceaux
$V$ et $W$ sous le préfaisceau $U$ sur $\Delta$.
\end{definition}

\begin{lemma}
\label{lemma-push-out}
Si $U, V, W$ sont des objets simpliciaux de la catégorie $\mathcal{C}$,
si $a : U \to V$, $b : U \to W$ sont des morphismes et si
$V\amalg_U W$ existe, alors
$$
\Mor(V\amalg_U W, T) =
\Mor(V, T) \times_{\Mor(U, T)}
\Mor(W, T)
$$
pour tout quatrième objet simplicial $T$ de $\mathcal{C}$.
\end{lemma}

\begin{proof}
Omis.
\end{proof}











\section{Produits d'objets cosimpliciaux}
\label{section-products-cosimplicial}

\noindent
Il serait naturel de définir le produit d'objets cosimpliciaux comme leur
produit dans la catégorie des objets cosimpliciaux. Il pourrait alors
arriver, ce qui prêterait à confusion, que ce produit existe sans être
décrit comme ci-dessous. Nous le définissons donc directement.

\begin{definition}
\label{definition-product-cosimplicial-objects}
Soit $\mathcal{C}$ une catégorie. Soient $U$ et $V$ des objets
cosimpliciaux de $\mathcal{C}$. Supposons que les produits
$U_n \times V_n$ existent dans $\mathcal{C}$. Le {\it produit de $U$ et
$V$} est l'objet cosimplicial $U \times V$ défini comme suit :
\begin{enumerate}
\item $(U \times V)_n = U_n \times V_n$,
\item pour tout $\varphi : [n] \to [m]$, l'application
$(U \times V)(\varphi) : U_n \times V_n \to U_m \times V_m$
est le produit $U(\varphi) \times V(\varphi)$.
\end{enumerate}
\end{definition}

\begin{lemma}
\label{lemma-product-cosimplicial-objects}
Si $U$ et $V$ sont des objets cosimpliciaux de la catégorie $\mathcal{C}$
et si $U \times V$ existe, alors
$$
\Mor(W, U \times V) =
\Mor(W, U) \times
\Mor(W, V)
$$
pour tout troisième objet cosimplicial $W$ de $\mathcal{C}$.
\end{lemma}

\begin{proof}
Omis.
\end{proof}

\section{Produits fibrés d'objets cosimpliciaux}
\label{section-fibre-products-cosimplicial}

\noindent
Il serait naturel de définir le produit fibré d'objets cosimpliciaux comme
leur produit fibré dans la catégorie des objets cosimpliciaux. Il pourrait
alors arriver, ce qui prêterait à confusion, que ce produit existe sans être
décrit comme ci-dessous. Nous le définissons donc directement.

\begin{definition}
\label{definition-fibre-product-cosimplicial-objects}
Soit $\mathcal{C}$ une catégorie. Soient $U, V, W$ des objets
cosimpliciaux de $\mathcal{C}$ et soient $a : V \to U$ et $b : W \to U$
des morphismes. Supposons que les produits fibrés
$V_n \times_{U_n} W_n$ existent dans $\mathcal{C}$. Le {\it produit
fibré de $V$ et $W$ au-dessus de $U$} est l'objet cosimplicial
$V \times_U W$ défini comme suit :
\begin{enumerate}
\item $(V \times_U W)_n = V_n \times_{U_n} W_n$,
\item pour tout $\varphi : [n] \to [m]$, l'application
$(V \times_U W)(\varphi) : V_n \times_{U_n} W_n \to V_m \times_{U_m} W_m$
est le produit $V(\varphi) \times_{U(\varphi)} W(\varphi)$.
\end{enumerate}
\end{definition}

\begin{lemma}
\label{lemma-fibre-product-cosimplicial-objects}
Si $U, V, W$ sont des objets cosimpliciaux de la catégorie $\mathcal{C}$,
si $a : V \to U$, $b : W \to U$ sont des morphismes et si
$V \times_U W$ existe, alors
$$
\Mor(T, V \times_U W) =
\Mor(T, V) \times_{\Mor(T, U)}
\Mor(T, W)
$$
pour tout quatrième objet cosimplicial $T$ de $\mathcal{C}$.
\end{lemma}

\begin{proof}
Omis.
\end{proof}














\section{Ensembles simpliciaux}
\label{section-simplicial-set}

\noindent
Soit $U$ un ensemble simplicial. Il est utile de considérer $U_0$ comme
l'ensemble des {\it $0$-simplexes}, $U_1$ comme l'ensemble des
{\it $1$-simplexes}, $U_2$ comme l'ensemble des {\it $2$-simplexes},
et ainsi de suite.

\medskip\noindent
On considère $s^n_j : U_n \to U_{n + 1}$ comme l'application qui associe
à un $n$-simplexe $A$ le $(n + 1)$-simplexe dégénéré $B$ dont l'arête
$(j, j + 1)$ est contractée sur le sommet $j$ de $A$. On considère
$d^n_j : U_n \to U_{n - 1}$ comme l'application qui associe à un
$n$-simplexe $A$ l'une de ses faces, celle qui omet le sommet $j$.
On peut ainsi visualiser géométriquement les relations entre les
applications $s^n_j$ et $d^n_j$.

\begin{definition}
\label{definition-terminology-simplicial-sets}
Soit $U$ un ensemble simplicial. On dit que $x$ est un
{\it $n$-simplexe de $U$} si $x$ est un élément de $U_n$. On dit que
$y$ est la $j$-ième {\it face de $x$} si $d^n_jx = y$. On dit que
$z$ est la $j$-ième {\it dégénérescence de $x$} si $z = s^n_jx$.
Un simplexe est dit {\it dégénéré} s'il est la dégénérescence d'un autre
simplexe.
\end{definition}

\noindent
Voici quelques exemples fondamentaux.

\begin{example}
\label{example-simplex-simplicial-set}
Pour tout $n \geq 0$, on note $\Delta[n]$ l'ensemble simplicial
$$
\Delta^{opp} \longrightarrow \textit{Sets},\quad
[k] \longmapsto \Mor_{\Delta}([k], [n])
$$
Nous laissons au lecteur le soin de vérifier les assertions suivantes.
Tout $m$-simplexe de $\Delta[n]$ avec $m > n$ est dégénéré.
Il existe un unique $n$-simplexe non dégénéré de $\Delta[n]$, à savoir
$\text{id}_{[n]}$.
\end{example}

\begin{lemma}
\label{lemma-simplex-map}
Soit $U$ un ensemble simplicial et soit $n \geq 0$ un entier.
Il existe une bijection canonique
$$
\Mor(\Delta[n], U)
\longrightarrow
U_n
$$
qui à un morphisme $\varphi$ associe la valeur de $\varphi$ sur l'unique
$n$-simplexe non dégénéré de $\Delta[n]$.
\end{lemma}

\begin{proof}
Omis.
\end{proof}

\begin{example}
\label{example-simplex-category}
Considérons la catégorie $\Delta/[n]$ des objets au-dessus de $[n]$ dans
$\Delta$, voir Categories, exemple \ref{categories-example-category-over-X}.
Il existe un foncteur $p : \Delta/[n] \to \Delta$.
La catégorie fibre de $p$ au-dessus de $[k]$, voir Categories,
section \ref{categories-section-fibred-groupoids}, a pour objets l'ensemble
$\Delta[n]_k$ des $k$-simplexes de $\Delta[n]$ et pour seuls morphismes
les identités. Pour tout morphisme $\varphi : [k] \to [l]$ de $\Delta$
et tout objet $\psi : [l] \to [n]$ de la catégorie fibre au-dessus de
$[l]$, il existe un unique objet au-dessus de $[k]$ muni d'un morphisme
au-dessus de $\varphi$, à savoir $\psi \circ \varphi : [k] \to [n]$.
Ainsi $\Delta/[n]$ est fibrée en ensembles au-dessus de $\Delta$.
Autrement dit, on peut voir $\Delta/[n]$ comme un préfaisceau d'ensembles
sur $\Delta$. Voir aussi Categories, exemple
\ref{categories-example-fibred-category-from-functor-of-points}.
Ce préfaisceau d'ensembles coïncide avec l'ensemble simplicial $\Delta[n]$.
En particulier, l'équation (\ref{equation-simplicial-set-presheaf}) et le
lemme \ref{lemma-simplex-map} donnent la formule
$$
\Mor_{\textit{PSh}(\Delta)}(\Delta/[n], U) = U_n
$$
pour tout ensemble simplicial $U$.
\end{example}

\begin{lemma}
\label{lemma-product-degenerate}
Soient $U$ et $V$ des ensembles simpliciaux et soient $a, b \geq 0$
des entiers. Supposons que tout $n$-simplexe de $U$ soit dégénéré si
$n > a$, et que tout $n$-simplexe de $V$ soit dégénéré si $n > b$.
Alors tout $n$-simplexe de $U \times V$ est dégénéré si $n > a + b$.
\end{lemma}

\begin{proof}
Supposons $n > a + b$ et soit
$(u, v) \in (U \times V)_n = U_n \times V_n$.
Par hypothèse, il existe $\alpha : [n] \to [a]$, $u' \in U_a$,
$\beta : [n] \to [b]$ et $v' \in V_b$ tels que
$u = U(\alpha)(u')$ et $v = V(\beta)(v')$. Puisque $n > a + b$,
il existe $0 \leq i \leq a + b$ tel que
$\alpha(i) = \alpha(i + 1)$ et $\beta(i) = \beta(i + 1)$.
Il s'ensuit immédiatement que $(u, v)$ appartient à l'image de
$s^{n - 1}_i$.
\end{proof}



\section{Objets simpliciaux tronqués et foncteurs squelette}
\label{section-skeleton}

\noindent
Notons $\Delta_{\leq n}$ la sous-catégorie pleine de $\Delta$ ayant pour
objets $[0], [1], [2], \ldots, [n]$, et soit $\mathcal{C}$ une catégorie.

\begin{definition}
\label{definition-truncated-simplicial-object}
Un {\it objet simplicial $n$-tronqué de $\mathcal{C}$} est un foncteur
contravariant de $\Delta_{\leq n}$ vers $\mathcal{C}$. Un
{\it morphisme d'objets simpliciaux $n$-tronqués} est une transformation
de foncteurs. On note la catégorie des objets simpliciaux $n$-tronqués de
$\mathcal{C}$ par le symbole $\text{Simp}_n(\mathcal{C})$.
\end{definition}

\noindent
Étant donné un objet simplicial $U$ de $\mathcal{C}$, sa troncature
$\text{sk}_n U$ est la restriction de $U$ à la sous-catégorie
$\Delta_{\leq n}$. Cela définit un {\it foncteur squelette}
$$
\text{sk}_n :
\text{Simp}(\mathcal{C}) \longrightarrow \text{Simp}_n(\mathcal{C})
$$
de la catégorie des objets simpliciaux de $\mathcal{C}$ vers celle des
objets simpliciaux $n$-tronqués de $\mathcal{C}$. Voir la remarque
\ref{remark-sk-literature} afin d'éviter toute confusion avec d'autres
foncteurs de la littérature.




\section{Produits avec des ensembles simpliciaux}
\label{section-product-with-simplicial-sets}

\noindent
Soit $\mathcal{C}$ une catégorie, soit $U$ un ensemble simplicial et soit
$V$ un objet simplicial de $\mathcal{C}$. On peut considérer le foncteur
covariant qui associe à tout objet simplicial $W$ de $\mathcal{C}$
l'ensemble
\begin{equation}
\label{equation-functor-product-with-simplicial-set}
\left\{
(f_{n, u} : V_n \to W_n)_{n \geq 0, u \in U_n}
\text{ tels que }
\begin{matrix}
\forall \varphi : [m] \to [n] \\
f_{m, U(\varphi)(u)} \circ V(\varphi) = W(\varphi) \circ f_{n, u}
\end{matrix}
\right\}
\end{equation}
Si ce foncteur est de la forme $\Mor_{\text{Simp}(\mathcal{C})}(Q, -)$,
on peut considérer $Q$ comme le produit de $U$ par $V$. Plutôt que de
formaliser cette idée, nous définissons directement le produit comme suit.

\begin{definition}
\label{definition-product-with-simplicial-set}
Soit $\mathcal{C}$ une catégorie dans laquelle la somme de deux objets
quelconques de $\mathcal{C}$ existe. Soit $U$ un ensemble simplicial et soit $V$ un objet
simplicial de $\mathcal{C}$. Supposons chaque $U_n$ fini et non vide.
On définit alors le {\it produit $U \times V$ de $U$ et $V$} comme
l'objet simplicial de $\mathcal{C}$ dont le terme de degré $n$ est
$$
(U \times V)_n = \coprod\nolimits_{u\in U_n} V_n
$$
et dont les applications associées à $\varphi : [m] \to [n]$ sont données
par le morphisme
$$
\coprod\nolimits_{u\in U_n} V_n
\longrightarrow
\coprod\nolimits_{u'\in U_m} V_m
$$
qui envoie la composante $V_n$ correspondant à $u$ sur la composante
$V_m$ correspondant à $u' = U(\varphi)(u)$ par le morphisme
$V(\varphi)$.
Plus généralement, si toutes les sommes affichées ci-dessus existent
(sans hypothèse globale sur $\mathcal{C}$), nous dirons que le
{\it produit $U \times V$ existe}.
\end{definition}

\begin{lemma}
\label{lemma-check-product-with-simplicial-set}
Soit $\mathcal{C}$ une catégorie dans laquelle la somme de deux objets
quelconques de $\mathcal{C}$ existe. Soit $U$ un ensemble simplicial et soit $V$ un objet
simplicial de $\mathcal{C}$. Supposons chaque $U_n$ fini et non vide.
Le foncteur
$W \mapsto \Mor_{\text{Simp}(\mathcal{C})}(U \times V, W)$
est canoniquement isomorphe au foncteur qui associe à $W$ l'ensemble de
l'équation (\ref{equation-functor-product-with-simplicial-set}).
\end{lemma}

\begin{proof}
Omis.
\end{proof}

\begin{lemma}
\label{lemma-back-to-U}
Soit $\mathcal{C}$ une catégorie dans laquelle la somme de deux objets
quelconques de $\mathcal{C}$ existe. Notons provisoirement $\textit{FSSets}$ la catégorie
des ensembles simpliciaux dont toutes les composantes sont finies et non
vides.
\begin{enumerate}
\item La règle $(U, V) \mapsto U \times V$ définit un foncteur
$\textit{FSSets} \times \text{Simp}(\mathcal{C})
\to \text{Simp}(\mathcal{C})$.
\item Pour tous $U$, $V$ comme ci-dessus, il existe un morphisme canonique
d'objets simpliciaux
$$
U \times V \longrightarrow V
$$
défini par l'identité sur chaque composante de
$(U \times V)_n = \coprod_u V_n$.
\end{enumerate}
\end{lemma}

\begin{proof}
Omis.
\end{proof}

\noindent
Étudions brièvement un cas particulier de la construction précédente.
Soit $\mathcal{C}$ une catégorie, soit $X$ un objet de $\mathcal{C}$ et
soit $k \geq 0$ un entier. Si toutes les sommes
$X \amalg \ldots \amalg X$ existent, alors, d'après la définition
ci-dessus, le produit
$$
X \times \Delta[k]
$$
existe, où $X$ est considéré comme l'objet simplicial constant
correspondant.

\begin{lemma}
\label{lemma-morphism-from-coproduct}
Avec $X$ et $k$ comme ci-dessus, on dispose, pour tout objet simplicial
$V$ de $\mathcal{C}$, de la bijection canonique
$$
\Mor_{\text{Simp}(\mathcal{C})}(X \times \Delta[k], V)
\longrightarrow
\Mor_\mathcal{C}(X, V_k).
$$
qui à $\gamma$ associe la restriction du morphisme $\gamma_k$ à la
composante correspondant à $\text{id}_{[k]}$. De même, pour tout
$n \geq k$, si $W$ est un objet simplicial $n$-tronqué de $\mathcal{C}$,
alors
$$
\Mor_{\text{Simp}_n(\mathcal{C})}(\text{sk}_n(X \times \Delta[k]), W)
=
\Mor_\mathcal{C}(X, W_k).
$$
\end{lemma}

\begin{proof}
Un morphisme $\gamma : X \times \Delta[k] \to V$ est donné par une
famille de morphismes $\gamma_\alpha : X \to V_n$, où
$\alpha : [n] \to [k]$. Ces morphismes doivent vérifier que, pour tout
$\varphi : [m] \to [n]$, les diagrammes
$$
\xymatrix{
X \ar[r]^{\gamma_\alpha} \ar[d]^{\text{id}_X} & V_n \ar[d]^{V(\varphi)} \\
X \ar[r]^{\gamma_{\alpha \circ \varphi}} & V_m
}
$$
soient commutatifs. En prenant $\alpha = \text{id}_{[k]}$, on voit que,
pour tout $\varphi : [m] \to [k]$, on a $\gamma_\varphi =
V(\varphi) \circ \gamma_{\text{id}_{[k]}}$. Le morphisme $\gamma$ est
donc déterminé par la valeur de $\gamma$ sur la composante correspondant à
$\text{id}_{[k]}$. Réciproquement, un morphisme $f : X \to V_k$ de ce
type définit aisément un morphisme $\gamma$ en posant
$\gamma_\alpha = V(\alpha) \circ f$.

\medskip\noindent
Le cas tronqué est analogue et laissé au lecteur.
\end{proof}

\noindent
Le cas $k = 0$ en fournit un exemple particulier. La formule du lemme
signifie alors simplement que
$$
\Mor_\mathcal{C}(X, V_0)
=
\Mor_{\text{Simp}(\mathcal{C})}(X, V)
$$
où, dans le membre de droite, $X$ désigne l'objet simplicial constant de
valeur $X$. Nous utiliserons désormais cette formule sans autre mention.




\section{Hom des ensembles simpliciaux vers les objets cosimpliciaux}
\label{section-hom-from-simplicial-sets-into-cosimplicial}

\noindent
Soit $\mathcal{C}$ une catégorie, soit $U$ un objet simplicial de
$\mathcal{C}$ et soit $V$ un objet cosimplicial de $\mathcal{C}$.
On obtient alors un ensemble cosimplicial $\Hom_\mathcal{C}(U, V)$
comme suit :
\begin{enumerate}
\item on pose
$\Hom_\mathcal{C}(U, V)_n = \Mor_\mathcal{C}(U_n, V_n)$, and
\item à $\varphi : [m] \to [n]$ on associe l'application
$\Hom_\mathcal{C}(U, V)_m \to \Hom_\mathcal{C}(U, V)_n$
donnée par $f \mapsto V(\varphi) \circ f \circ U(\varphi)$.
\end{enumerate}
C'est ce qui motive la définition suivante.

\begin{definition}
\label{definition-hom-deltak-cosimplicial}
Soit $\mathcal{C}$ une catégorie admettant les produits finis. Soit $V$
un objet cosimplicial de $\mathcal{C}$ et soit $U$ un ensemble simplicial
tel que chaque $U_n$ soit fini et non vide. On définit
{\it $\Hom(U, V)$} comme l'objet cosimplicial de $\mathcal{C}$ suivant :
\begin{enumerate}
\item on pose $\Hom(U, V)_n = \prod_{u \in U_n} V_n$, autrement dit
l'unique objet de $\mathcal{C}$ dont les points à valeurs dans $X$
satisfont
$$
\Mor_\mathcal{C}(X, \Hom(U, V)_n)
=
\text{Map}(U_n, \Mor_\mathcal{C}(X, V_n))
$$
et
\item à $\varphi : [m] \to [n]$ on associe l'application
$\Hom(U, V)_m \to \Hom(U, V)_n$
donnée par $f \mapsto V(\varphi) \circ f \circ U(\varphi)$ sur les
points à valeurs dans $X$ ci-dessus.
\end{enumerate}
\end{definition}

\noindent
Nous laissons au lecteur le soin d'expliciter cette définition en termes
d'applications entre produits. Remarquons également que la construction est
fonctorielle en $U$ (contravariamment) et en $V$ (covariamment), exactement
comme au lemme \ref{lemma-back-to-U} pour les produits d'ensembles
simpliciaux avec des objets simpliciaux.




\section{Hom des ensembles cosimpliciaux vers les objets simpliciaux}
\label{section-hom-from-cosimplicial-sets-into-simplicial}

\noindent
Soit $\mathcal{C}$ une catégorie, soit $U$ un objet cosimplicial de
$\mathcal{C}$ et soit $V$ un objet simplicial de $\mathcal{C}$.
On obtient alors un ensemble simplicial $\Hom_\mathcal{C}(U, V)$
comme suit :
\begin{enumerate}
\item on pose
$\Hom_\mathcal{C}(U, V)_n = \Mor_\mathcal{C}(U_n, V_n)$, and
\item à $\varphi : [m] \to [n]$ on associe l'application
$\Hom_\mathcal{C}(U, V)_n \to \Hom_\mathcal{C}(U, V)_m$
donnée par $f \mapsto V(\varphi) \circ f \circ U(\varphi)$.
\end{enumerate}
C'est ce qui motive la définition suivante.

\begin{definition}
\label{definition-hom-deltak-simplicial}
Soit $\mathcal{C}$ une catégorie admettant les produits finis. Soit $V$
un objet simplicial de $\mathcal{C}$ et soit $U$ un ensemble cosimplicial
tel que chaque $U_n$ soit fini et non vide. On définit
{\it $\Hom(U, V)$} comme l'objet simplicial de $\mathcal{C}$ suivant :
\begin{enumerate}
\item on pose $\Hom(U, V)_n = \prod_{u \in U_n} V_n$, autrement dit
l'unique objet de $\mathcal{C}$ dont les points à valeurs dans $X$
satisfont
$$
\Mor_\mathcal{C}(X, \Hom(U, V)_n)
=
\text{Map}(U_n, \Mor_\mathcal{C}(X, V_n))
$$
et
\item à $\varphi : [m] \to [n]$ on associe l'application
$\Hom(U, V)_n \to \Hom(U, V)_m$
donnée par $f \mapsto V(\varphi) \circ f \circ U(\varphi)$ sur les
points à valeurs dans $X$ ci-dessus.
\end{enumerate}
\end{definition}

\noindent
Nous laissons au lecteur le soin d'expliciter cette définition en termes
d'applications entre produits. La construction est fonctorielle en $U$
(contravariamment) et en $V$ (covariamment), exactement comme au lemme
\ref{lemma-back-to-U} pour les produits d'ensembles simpliciaux avec des
objets simpliciaux.

\medskip\noindent
Explicitons la construction précédente dans un cas particulier. Soit $X$
un objet d'une catégorie $\mathcal{C}$. Supposons que les produits
$X \times \ldots \times X$ existent et soit $k$ un entier. Considérons
l'objet simplicial $U$ de termes
$$
U_n = \prod\nolimits_{\alpha \in \Mor([k], [n])} X
$$
et dont les applications associées à $\varphi : [m] \to [n]$ sont
$$
U(\varphi) :
\prod\nolimits_{\alpha \in \Mor([k], [n])} X
\longrightarrow
\prod\nolimits_{\alpha' \in \Mor([k], [m])} X, \quad
(f_{\alpha})_{\alpha} \longmapsto
(f_{\varphi \circ \alpha'})_{\alpha'}
$$
En « coordonnées », l'élément $(x_\alpha)_\alpha$ est envoyé sur
$(x_{\varphi \circ \alpha'})_{\alpha'}$. Nous affirmons que cet objet
est égal à $\Hom(C[k], X)$, où $X$ est considéré comme l'objet simplicial
constant de valeur $X$ et où $C[k]$ est l'ensemble cosimplicial de
l'exemple \ref{example-simplex-cosimplicial-set}.

\begin{lemma}
\label{lemma-morphism-into-product}
Avec $X$, $k$ et $U$ comme ci-dessus :
\begin{enumerate}
\item pour tout objet simplicial $V$ de $\mathcal{C}$, on dispose de la
bijection canonique
$$
\Mor_{\text{Simp}(\mathcal{C})}(V, U)
\longrightarrow
\Mor_\mathcal{C}(V_k, X).
$$
qui à $\gamma$ associe le morphisme $\gamma_k$ composé avec la projection
sur le facteur correspondant à $\text{id}_{[k]}$ ;
\item de même, si $W$ est un objet simplicial $k$-tronqué de
$\mathcal{C}$, alors
$$
\Mor_{\text{Simp}_k(\mathcal{C})}(W, \text{sk}_k U)
=
\Mor_\mathcal{C}(W_k, X).
$$
\item l'objet $U$ construit ci-dessus est une réalisation de
$\Hom(C[k], X)$, où $C[k]$ est l'ensemble cosimplicial de l'exemple
\ref{example-simplex-cosimplicial-set}.
\end{enumerate}
\end{lemma}

\begin{proof}
Montrons d'abord (1). Supposons que $\gamma : V \to U$ soit un morphisme.
Il est donné par une famille de morphismes
$\gamma_{\alpha} : V_n \to X$ pour $\alpha : [k] \to [n]$.
Ces morphismes doivent vérifier que, pour tout
$\varphi : [m] \to [n]$, les diagrammes
$$
\xymatrix{
X \ar[d]^{\text{id}_X} &
V_n \ar[d]^{V(\varphi)}
\ar[l]^{\gamma_{\varphi \circ \alpha'}} \\
X &
V_m \ar[l]_{\gamma_{\alpha'}}
}
$$
soient commutatifs pour tout $\alpha' : [k] \to [m]$.
En prenant $\alpha' = \text{id}_{[k]}$, on voit que, pour tout
$\varphi : [k] \to [n]$, on a $\gamma_\varphi =
\gamma_{\text{id}_{[k]}} \circ V(\varphi)$. Le morphisme $\gamma$ est
donc déterminé par la composante de $\gamma_k$ correspondant à
$\text{id}_{[k]}$. Réciproquement, un morphisme $f : V_k \to X$ de ce
type définit aisément un morphisme $\gamma$ en posant
$\gamma_\alpha = f \circ V(\alpha)$.

\medskip\noindent
Le cas tronqué est analogue et laissé au lecteur.

\medskip\noindent
L'assertion (3) résulte immédiatement de la construction de $U$ et de
l'égalité $C[k]_n = \Mor([k], [n])$, dont les deux membres sont les
ensembles d'indices utilisés dans la construction de $U_n$.
\end{proof}




\section{Hom interne}
\label{section-internal-hom}

\noindent
Soit $\mathcal{C}$ une catégorie admettant les produits finis non vides.
Soient $U$, $V$ des objets simpliciaux de $\mathcal{C}$. Dans certains
cas, le foncteur
$$
\text{Simp}(\mathcal{C})^{opp} \longrightarrow \textit{Sets}, \quad
W \longmapsto \Mor_{\text{Simp}(\mathcal{C})}(W \times V, U)
$$
est représentable. On note alors $\SheafHom(V, U)$ l'objet simplicial
de $\mathcal{C}$ ainsi obtenu et l'on dit que le
{\it Hom interne de $V$ vers $U$ existe}. En outre, pour tout $X$ de
$\mathcal{C}$, on aurait
\begin{align*}
\Mor_\mathcal{C}(X, \SheafHom(V, U)_n)
& =
\Mor_{\text{Simp}(\mathcal{C})}(X \times \Delta[n], \SheafHom(V, U)) \\
& =
\Mor_{\text{Simp}(\mathcal{C})}(X \times \Delta[n]\times V, U) \\
& =
\Mor_{\text{Simp}(\mathcal{C})}(X, \SheafHom(\Delta[n] \times V, U)) \\
& =
\Mor_\mathcal{C}(X, \SheafHom(\Delta[n] \times V, U)_0)
\end{align*}
pourvu que $\SheafHom(\Delta[n] \times V, U)$ existe également.
La première et la dernière égalité résultent du lemme
\ref{lemma-morphism-from-coproduct}.

\medskip\noindent
Nous en retenons qu'étant donnés $U$ et $V$, si l'on veut construire le
Hom interne, il convient de construire les objets
$$
\SheafHom(\Delta[n] \times V, U)_0
$$
car ils doivent constituer le terme de degré $n$ de $\SheafHom(V, U)$.
Dans la section suivante, nous étudions une construction d'objets
simpliciaux « $\Hom(\Delta[n], U)$ ».
\section{Hom des ensembles simpliciaux vers les objets simpliciaux}
\label{section-hom-from-simplicial-sets}

\noindent
À la suite de la discussion sur le Hom interne, nous définissons l'objet
simplicial qui doit classifier les morphismes d'un ensemble simplicial
vers un objet simplicial donné de la catégorie $\mathcal{C}$.

\begin{definition}
\label{definition-hom-from-simplicial-set}
Soit $\mathcal{C}$ une catégorie dans laquelle la somme de deux objets
quelconques existe. Soit $U$ un ensemble simplicial tel que $U_n$ soit
fini et non vide pour tout $n \geq 0$. Soit $V$ un objet simplicial de
$\mathcal{C}$. On note {\it $\Hom(U, V)$} tout objet simplicial de
$\mathcal{C}$ tel que
$$
\Mor_{\text{Simp}(\mathcal{C})}(W, \Hom(U, V))
=
\Mor_{\text{Simp}(\mathcal{C})}(W \times U, V)
$$
fonctoriellement en l'objet simplicial $W$ de $\mathcal{C}$.
\end{definition}

\noindent
Bien entendu, $\Hom(U, V)$ n'existe pas nécessairement. D'après la
discussion de la section \ref{section-internal-hom}, on s'attend en outre,
lorsqu'il existe, à ce que
$\Hom(U, V)_n = \Hom(U \times \Delta[n], V)_0$.
Nous n'employons pas la notation italique pour ces objets Hom, puisque
$\Hom(U, V)$ n'est pas un Hom interne.

\begin{lemma}
\label{lemma-exists-hom-0-from-simplicial-set}
Supposons que la catégorie $\mathcal{C}$ admette les sommes de deux objets
et les limites dénombrables. Soit $U$ un ensemble simplicial tel que $U_n$
soit fini et non vide pour tout $n \geq 0$. Soit $V$ un objet simplicial
de $\mathcal{C}$. Alors le foncteur
\begin{eqnarray*}
\mathcal{C}^{opp} & \longrightarrow & \textit{Sets} \\
X
& \longmapsto &
\Mor_{\text{Simp}(\mathcal{C})}(X \times U, V)
\end{eqnarray*}
est représentable.
\end{lemma}

\begin{proof}
Un morphisme de $X \times U$ vers $V$ est donné par une famille de
morphismes $f_u : X \to V_n$, où $n \geq 0$ et $u \in U_n$.
Une telle famille définit effectivement un morphisme si et seulement si,
pour tout $\varphi : [m] \to [n]$, tous les diagrammes
$$
\xymatrix{
X \ar[r]^{f_u} \ar[d]_{\text{id}_X} & V_n \ar[d]^{V(\varphi)} \\
X \ar[r]^{f_{U(\varphi)(u)}} & V_m
}
$$
commutent. Il est donc naturel d'introduire une catégorie $\mathcal{U}$
et un foncteur
$\mathcal{V} : \mathcal{U}^{opp} \to \mathcal{C}$
comme suit :
\begin{enumerate}
\item l'ensemble des objets de $\mathcal{U}$ est
$\coprod_{n \geq 0} U_n$,
\item un morphisme de $u' \in U_m$ vers $u \in U_n$ est un
$\varphi : [m] \to [n]$ tel que $U(\varphi)(u) = u'$,
\item pour $u \in U_n$, on pose $\mathcal{V}(u) = V_n$, et
\item pour $\varphi : [m] \to [n]$ tel que $U(\varphi)(u) = u'$,
on pose $\mathcal{V}(\varphi) = V(\varphi) : V_n \to V_m$.
\end{enumerate}
Il est alors clair que notre foncteur n'est autre que le foncteur qui
définit
$$
\lim_{\mathcal{U}^{opp}} \mathcal{V}
$$
Si $\mathcal{C}$ admet les limites dénombrables, cette limite existe donc,
et avec elle un objet représentant le foncteur du lemme.
\end{proof}

\begin{lemma}
\label{lemma-exists-hom-0-from-simplicial-set-finite}
Supposons que la catégorie $\mathcal{C}$ admette les sommes de deux objets
et les limites finies. Soit $U$ un ensemble simplicial tel que $U_n$ soit
fini et non vide pour tout $n \geq 0$. Supposons que tous les
$n$-simplexes de $U$ soient dégénérés pour tout $n \gg 0$.
Soit $V$ un objet simplicial de $\mathcal{C}$. Alors le foncteur
\begin{eqnarray*}
\mathcal{C}^{opp} & \longrightarrow & \textit{Sets} \\
X
& \longmapsto &
\Mor_{\text{Simp}(\mathcal{C})}(X \times U, V)
\end{eqnarray*}
est représentable.
\end{lemma}

\begin{proof}
Il faut montrer que la catégorie $\mathcal{U}$ décrite dans la preuve du
lemme \ref{lemma-exists-hom-0-from-simplicial-set} possède une
sous-catégorie finie $\mathcal{U}'$ telle que la limite de $\mathcal{V}$
sur $\mathcal{U}'$ soit la même que la limite de $\mathcal{V}$ sur
$\mathcal{U}$.
Nous utiliserons Categories, lemme \ref{categories-lemma-initial}.
Pour $m > 0$, notons $\mathcal{U}_{\leq m}$ la sous-catégorie pleine ayant
pour objets $\coprod_{0 \leq n \leq m} U_m$. Soit $m_0$ un entier tel que
tout $n$-simplexe de l'ensemble simplicial $U$ soit dégénéré si $n > m_0$.
Pour tout $m \geq m_0$ assez grand, la sous-catégorie
$\mathcal{U}_{\leq m}$ satisfait à la propriété (1) de Categories,
définition \ref{categories-definition-initial}.

\medskip\noindent
Supposons que $u \in U_n$ et
$u' \in U_{n'}$, avec $n, n' \leq m_0$, et que
$\varphi : [k] \to [n]$, $\varphi' : [k] \to [n']$
soient des morphismes tels que $U(\varphi)(u) = U(\varphi')(u')$.
Un argument combinatoire élémentaire montre que, si $k > 2m_0$,
il existe un indice $0 \leq i \leq 2m_0$ tel que
$\varphi(i) =\varphi(i + 1)$ et $\varphi'(i) = \varphi'(i + 1)$.
(Le principe des tiroirs montre déjà que cela fonctionne si
$k > m_0^2$, ce qui suffit de toute manière à l'argument ci-dessous.)
Ainsi, si $k > 2m_0$, on peut écrire
$\varphi = \psi \circ \sigma^{k - 1}_i$ et
$\varphi' = \psi' \circ \sigma^{k - 1}_i$, pour certains
$\psi : [k - 1] \to [n]$ et $\psi' : [k - 1] \to [n']$.
Comme $s^{k - 1}_i : U_{k - 1} \to U_k$ est injective,
voir le lemme \ref{lemma-si-injective}, on en déduit également que
$U(\psi)(u) = U(\psi')(u')$. En poursuivant de la sorte,
on conclut que, étant donnés des morphismes
$u \to z$ et $u' \to z$ de $\mathcal{U}$,
avec $u, u' \in \mathcal{U}_{\leq m_0}$, il existe
un diagramme commutatif
$$
\xymatrix{
u \ar[rd] \ar[rrd] & & \\
& a \ar[r] & z \\
u' \ar[ru] \ar[rru]
}
$$
où $a \in \mathcal{U}_{\leq 2m_0}$.

\medskip\noindent
On en déduit aisément que la sous-catégorie finie
$\mathcal{U}_{\leq 2m_0}$ convient. En effet, soient
$x' \in U_n$ et $x'' \in U_{n'}$, avec $n, n' \leq 2m_0$, ainsi que
des morphismes $x' \to x$ et $x'' \to x$ de $\mathcal{U}$
ayant même but. Par le choix de $m_0$, on peut trouver des objets
$u, u'$ de $\mathcal{U}_{\leq m_0}$ et des morphismes
$u \to x'$, $u' \to x''$.
D'après ce qui précède, on peut trouver $a \in \mathcal{U}_{\leq 2m_0}$
et des morphismes $u \to a$, $u' \to a$ tels que
$$
\xymatrix{
u \ar[rd] \ar[rrd] \ar[r] & x' \ar[rd] & \\
& a \ar[r] & x \\
u' \ar[ru] \ar[rru] \ar[r] & x'' \ar[ru] &
}
$$
soit commutatif. En faisant tourner ce diagramme de 90 degrés dans le
sens des aiguilles d'une montre, on obtient le diagramme voulu comme
dans Categories, définition \ref{categories-definition-initial}, (2).
\end{proof}

\begin{lemma}
\label{lemma-exists-hom-from-simplicial-set-finite}
Supposons que la catégorie $\mathcal{C}$ possède les sommes de deux
objets quelconques et les limites finies. Soit $U$ un ensemble simplicial,
avec $U_n$ fini et non vide pour tout $n \geq 0$. Supposons que tous les
$n$-simplexes de $U$ soient dégénérés pour tout $n \gg 0$.
Soit $V$ un objet simplicial de $\mathcal{C}$.
Alors $\Hom(U, V)$ existe et l'on a les égalités attendues
$$
\Hom(U, V)_n = \Hom(U \times \Delta[n], V)_0.
$$
\end{lemma}

\begin{proof}
Construisons cet objet simplicial comme suit.
Pour $n \geq 0$, notons $\Hom(U, V)_n$ l'objet de
$\mathcal{C}$ qui représente le foncteur
$$
X
\longmapsto
\Mor_{\text{Simp}(\mathcal{C})}(X \times U \times \Delta[n], V)
$$
Celui-ci existe par le lemme
\ref{lemma-exists-hom-0-from-simplicial-set-finite}, car
$U \times \Delta[n]$ est un ensemble simplicial dont les ensembles de
simplexes sont finis et qui n'a aucun simplexe non dégénéré en degré
assez grand, voir le lemme \ref{lemma-product-degenerate}.
À partir de $\varphi : [m] \to [n]$, on obtient une application induite
d'ensembles simpliciaux $\varphi : \Delta[m] \to \Delta[n]$. On obtient
donc un morphisme
$X \times U \times \Delta[m] \to X \times U \times \Delta[n]$
fonctoriel en $X$, puis une transformation de foncteurs, laquelle donne
$$
\Hom(U, V)(\varphi) :
\Hom(U, V)_n
\longrightarrow
\Hom(U, V)_m.
$$
Cela définit manifestement un foncteur contravariant
$\Hom(U, V)$ de $\Delta$ dans la catégorie $\mathcal{C}$.
Autrement dit, on obtient un objet simplicial de $\mathcal{C}$.

\medskip\noindent
Il reste à montrer que $\Hom(U, V)$ satisfait à la propriété
universelle voulue
$$
\Mor_{\text{Simp}(\mathcal{C})}(W, \Hom(U, V))
=
\Mor_{\text{Simp}(\mathcal{C})}(W \times U, V)
$$
Pour le voir, soit $f : W \to \Hom(U, V)$. Construisons l'élément
$f' : W \times U \to V$ du membre de droite.
Par construction, chaque $f_n : W_n \to \Hom(U, V)_n$
correspond à un morphisme
$f_n : W_n \times U \times \Delta[n] \to V$. De plus, pour tout
morphisme $\varphi : [m] \to [n]$, le diagramme
$$
\xymatrix{
W_n \times U \times \Delta[m]
\ar[rr]_{W(\varphi)\times \text{id} \times \text{id}}
\ar[d]_{\text{id} \times \text{id} \times \varphi} & &
W_m \times U \times \Delta[m] \ar[d]^{f_m} \\
W_n \times U \times \Delta[n] \ar[rr]^{f_n} & & V
}
$$
est commutatif. Pour $\psi : [n] \to [k]$ appartenant à
$(\Delta[n])_k$, notons
$(f_n)_{k, \psi} : W_n \times U_k \to V_k$ la composante de
$(f_n)_k$ correspondant à l'élément $\psi$. Définissons
$f'_n : W_n \times U_n \to V_n$ par
$f'_n = (f_n)_{n, \text{id}}$, c'est-à-dire comme la restriction de
$(f_n)_n : W_n \times U_n \times (\Delta[n])_n \to V_n$
à $W_n \times U_n \times \text{id}_{[n]}$.
Pour voir que la famille $(f'_n)$ définit un morphisme d'objets
simpliciaux, il faut montrer, pour tout $\varphi : [m] \to [n]$, que
$V(\varphi) \circ f'_n =
f'_m \circ W(\varphi) \times U(\varphi)$.
Le diagramme commutatif ci-dessus affirme que
$(f_n)_{m, \varphi} : W_n \times U_m \to V_m$
est égal à
$(f_m)_{m, \text{id}} \circ W(\varphi) :
W_n \times U_m \to V_m$.
Or, puisque $f_n$ est un morphisme d'objets simpliciaux, le diagramme
$$
\xymatrix{
W_n \times U_n \times (\Delta[n])_n
\ar[r]_-{(f_n)_n}
\ar[d]_{\text{id} \times U(\varphi) \times \varphi}
& V_n \ar[d]^{V(\varphi)} \\
W_n \times U_m \times (\Delta[n])_m \ar[r]^-{(f_n)_m} & V_m
}
$$
est commutatif. Il en résulte que
$(f_n)_{m, \varphi} \circ U(\varphi)$ est égal à
$V(\varphi) \circ (f_n)_{n, \text{id}}$. Au total, on obtient
$
V(\varphi) \circ (f_n)_{n, \text{id}}
=
(f_n)_{m, \varphi} \circ U(\varphi)
=
(f_m)_{m, \text{id}} \circ W(\varphi)\circ U(\varphi)
=
(f_m)_{m, \text{id}} \circ W(\varphi)\times U(\varphi)
$
comme voulu.

\medskip\noindent
Réciproquement, étant donné un morphisme
$f' : W \times U \to V$, définissons un morphisme
$f : W \to \Hom(U, V)$ comme suit. Par le lemme
\ref{lemma-morphism-from-coproduct}, le morphisme
$\text{id} : W_n \to W_n$ correspond à un unique morphisme
$c_n : W_n \times \Delta[n] \to W$. On peut donc considérer le composé
$$
W_n \times \Delta[n] \times U
\xrightarrow{c_n}
W \times U
\xrightarrow{f'}
V.
$$
Par construction, celui-ci correspond à un unique morphisme
$f_n : W_n \to \Hom(U, V)_n$. Nous laissons au lecteur le soin de
vérifier que ces morphismes définissent le morphisme d'objets simpliciaux
voulu.

\medskip\noindent
Nous laissons également au lecteur le soin de vérifier que les opérations
$f \mapsto f'$ et $f' \mapsto f$ sont inverses l'une de l'autre.
\end{proof}

\begin{lemma}
\label{lemma-hom-from-coprod}
Supposons que la catégorie $\mathcal{C}$ possède les sommes de deux
objets quelconques et les limites finies. Soient
$a : U \to V$, $b : U \to W$ des morphismes d'ensembles simpliciaux.
Supposons que $U_n, V_n, W_n$ soient finis et non vides pour tout
$n \geq 0$, et que tous les $n$-simplexes de $U, V, W$ soient dégénérés
pour tout $n \gg 0$.
Soit $T$ un objet simplicial de $\mathcal{C}$. Alors
$$
\Hom(V, T) \times_{\Hom(U, T)} \Hom(W, T)
=
\Hom(V \amalg_U W, T)
$$
Autrement dit, le produit fibré du membre de gauche est représenté par
l'objet Hom du membre de droite.
\end{lemma}

\begin{proof}
Par le lemme \ref{lemma-exists-hom-from-simplicial-set-finite}, tous les
objets $\Hom$ requis existent et possèdent les propriétés fonctorielles
attendues. On peut alors identifier le terme de degré $n$ du membre de
gauche à l'objet qui représente le foncteur associant à $X$ le premier
ensemble de la suite d'égalités fonctorielles suivante :
\begin{align*}
&
\Mor(X \times \Delta[n],
\Hom(V, T) \times_{\Hom(U, T)} \Hom(W, T)) \\
& =
\Mor(X \times \Delta[n], \Hom(V, T))
\times_{\Mor(X \times \Delta[n], \Hom(U, T))}
\Mor(X \times \Delta[n], \Hom(W, T)) \\
& =
\Mor(X \times \Delta[n] \times V, T)
\times_{\Mor(X \times \Delta[n] \times U, T)}
\Mor(X \times \Delta[n] \times W, T) \\
& =
\Mor(X \times \Delta[n] \times (V \amalg_U W), T)
\end{align*}
On a utilisé ici le fait que
$$
(X \times \Delta[n] \times V)
\amalg_{X \times \Delta[n] \times U}
(X \times \Delta[n] \times W)
=
X \times \Delta[n] \times (V \amalg_U W)
$$
ce qui se vérifie aisément degré par degré. Le lemme en résulte, puisque
le dernier terme de la suite d'égalités affichée correspond à
$\Hom(V \amalg_U W, T)_n$.
\end{proof}













\section{Scindement des objets simpliciaux}
\label{section-splitting}

\noindent
Un sous-objet $N$ d'un objet $X$ de la catégorie $\mathcal{C}$ est un
objet $N$ de $\mathcal{C}$ muni d'un monomorphisme $N \to X$.
Bien entendu, on dit (par abus de notation) que les sous-objets
$N$, $N'$ sont égaux s'il existe un isomorphisme $N \to N'$ compatible
avec les morphismes vers $X$. L'ensemble des sous-objets est ordonné
partiellement. (Cela tient à nos conventions sur les catégories ; ce ne
serait pas vrai, par exemple, pour la catégorie des espaces à homotopie près.)

\begin{definition}
\label{definition-split}
Soit $\mathcal{C}$ une catégorie admettant les sommes finies non vides.
On dit qu'un objet simplicial $U$ de $\mathcal{C}$ est {\it scindé}
s'il existe des sous-objets $N(U_m)$ de $U_m$, $m \geq 0$, tels que
\begin{equation}
\label{equation-splitting}
\coprod\nolimits_{\varphi : [n] \to [m]\text{ surjective}}
N(U_m)
\longrightarrow
U_n
\end{equation}
soit un isomorphisme pour tout $n \geq 0$. Si $U$ est un objet
simplicial $r$-tronqué de $\mathcal{C}$, on dit que $U$ est {\it scindé}
s'il existe des sous-objets $N(U_m)$ de $U_m$, $r \geq m \geq 0$,
tels que (\ref{equation-splitting}) soit un isomorphisme pour
$r \geq n \geq 0$.
\end{definition}

\noindent
Dans ce cas, $N(U_0) = U_0$. Ensuite,
$U_1 = U_0 \amalg N(U_1)$. Puis
$$
U_2 = U_0 \amalg N(U_1) \amalg N(U_1) \amalg N(U_2).
$$
Dans de nombreuses catégories $\mathcal{C}$, tout objet simplicial est scindé.

\begin{lemma}
\label{lemma-splitting-simplicial-sets}
Soit $U$ un ensemble simplicial. Alors $U$ possède un unique scindement,
pour lequel $N(U_m)$ est l'ensemble des $m$-simplexes non dégénérés.
\end{lemma}

\begin{proof}
Il résulte immédiatement de la définition que, s'il existe un scindement,
$N(U_m)$ est nécessairement l'ensemble des simplexes non dégénérés.
Soit $x \in U_n$. Supposons qu'il existe des surjections
$\varphi : [n] \to [k]$ et $\psi : [n] \to [l]$, ainsi que des simplexes
non dégénérés $y \in U_k$, $z \in U_l$, tels que
$x = U(\varphi)(y)$ et $x = U(\psi)(z)$. Choisissons une section
$\xi : [l] \to [n]$ de $\psi$, c'est-à-dire
$\psi \circ \xi = \text{id}_{[l]}$. Alors $z = U(\xi)(x)$, d'où
$z = U(\xi)(x) = U(\varphi \circ \xi)(y)$. Comme $z$ est non dégénéré,
on en déduit que $\varphi \circ \xi : [l] \to [k]$ est surjective, et
donc que $l \geq k$. De même, $k \geq l$. Ainsi,
$\varphi \circ \xi : [l] \to [k]$ est nécessairement l'identité pour
tout choix d'une section $\xi$ de $\psi$. Il en résulte aisément que
$\psi = \varphi$.
\end{proof}

\noindent
Il peut naturellement arriver qu'une application d'ensembles simpliciaux
envoie un $n$-simplexe non dégénéré sur un $n$-simplexe dégénéré.
Le scindement du lemme \ref{lemma-splitting-simplicial-sets} n'est donc
pas fonctoriel. Voici un cas où il l'est.

\begin{lemma}
\label{lemma-injective-map-simplicial-sets}
Soit $f : U \to V$ un morphisme d'ensembles simpliciaux. Supposons que
(a) l'image de tout simplexe non dégénéré de $U$ soit un simplexe non
dégénéré de $V$ et que (b) la restriction de $f$, comme application de
l'ensemble des simplexes non dégénérés de $U$ vers celui des simplexes
non dégénérés de $V$, soit injective. Alors $f_n$ est injective pour tout
$n$. Le même énoncé vaut en remplaçant « injective » par « surjective »
ou « bijective ».
\end{lemma}

\begin{proof}
Sous l'hypothèse (a), l'application $f$ respecte les décompositions en
réunions disjointes du scindement du lemme
\ref{lemma-splitting-simplicial-sets} ; autrement dit, on obtient des
diagrammes commutatifs
$$
\xymatrix{
\coprod\nolimits_{\varphi : [n] \to [m]\text{ surjective}}
N(U_m)
\ar[r] \ar[d] &
U_n \ar[d] \\
\coprod\nolimits_{\varphi : [n] \to [m]\text{ surjective}}
N(V_m)
\ar[r] &
V_n.
}
$$
L'hypothèse (b) montre alors que la flèche verticale de gauche est
injective (resp.\ surjective, resp.\ bijective).
\end{proof}

\begin{lemma}
\label{lemma-simplicial-set-n-skel-sub}
Soit $U$ un ensemble simplicial et soit $n \geq 0$ un entier. La règle
$$
U'_m = \bigcup\nolimits_{\varphi : [m] \to [i], \ i\leq n} \Im(U(\varphi))
$$
définit un sous-ensemble simplicial $U' \subset U$ tel que
$U'_i = U_i$ pour $i \leq n$. De plus, tous les $m$-simplexes de $U'$
sont dégénérés pour tout $m > n$.
\end{lemma}

\begin{proof}
Si $x \in U_m$ et $x = U(\varphi)(y)$ pour un certain $y \in U_i$,
$i \leq n$, et un certain $\varphi : [m] \to [i]$, alors l'image
$U(\psi)(x)$, pour tout $\psi : [m'] \to [m]$, est égale à
$U(\varphi \circ \psi)(y)$, où
$\varphi \circ \psi : [m'] \to [i]$. Ainsi $U'$ est un ensemble
simplicial. Par construction, tous les simplexes de dimension au moins
$n + 1$ sont dégénérés.
\end{proof}

\begin{lemma}
\label{lemma-splitting-simplicial-groups}
Soit $U$ un groupe abélien simplicial. Alors $U$ possède un scindement
obtenu en prenant $N(U_0) = U_0$ et, pour $m \geq 1$,
$$
N(U_m) = \bigcap\nolimits_{i = 0}^{m - 1} \Ker(d^m_i).
$$
De plus, ce scindement est fonctoriel sur la catégorie des groupes
abéliens simpliciaux.
\end{lemma}

\begin{proof}
Montrons par récurrence sur $n$ que le choix de $N(U_m)$ indiqué dans
le lemme assure que (\ref{equation-splitting}) est un isomorphisme pour
$m \leq n$. C'est clair pour $n = 0$. Dans la suite de la démonstration,
nous omettrons les exposants des applications $d_i$ et $s_i$ afin
d'alléger les notations. Nous utiliserons aussi à plusieurs reprises les
relations de la remarque \ref{remark-relations}.

\medskip\noindent
Commençons par une remarque générale. Pour $0 \leq i \leq m$ et
$z \in U_m$, on a $d_i(s_i(z)) = z$. Tout $x \in U_{m + 1}$ s'écrit
donc de manière unique sous la forme $x = x' + x''$, avec
$d_i(x') = 0$ et $x'' \in \Im(s_i)$, en prenant
$x' = (x - s_i(d_i(x)))$ et $x'' = s_i(d_i(x))$. De plus, l'élément
$z \in U_m$ tel que $x'' = s_i(z)$ est unique puisque $s_i$ est injective.

\medskip\noindent
Voici un procédé pour décomposer un élément quelconque
$x \in U_{n + 1}$. Écrivons d'abord
$x = x_0 + s_0(z_0)$ avec $d_0(x_0) = 0$, puis
$x_0 = x_1 + s_1(z_1)$ avec $d_1(x_1) = 0$. En continuant ainsi, on obtient
\begin{eqnarray*}
x & = & x_0 + s_0(z_0), \\
x_0 & = & x_1 + s_1(z_1), \\
x_1 & = & x_2 + s_2(z_2), \\
\ldots & \ldots & \ldots \\
x_{n - 1} & = & x_n + s_n(z_n)
\end{eqnarray*}
où $d_i(x_i) = 0$ pour tout $i = n, \ldots, 0$.
D'après la remarque générale ci-dessus, tous les $x_i$ et $z_i$ sont
déterminés de manière unique par $x$. Nous affirmons que
$x_i \in
\Ker(d_0) \cap
\Ker(d_1) \cap
\ldots \cap
\Ker(d_i)$
et
$z_i \in
\Ker(d_0) \cap
\ldots \cap
\Ker(d_{i - 1})$
pour $i = n, \ldots, 0$. Ici comme dans la suite, une intersection vide
de noyaux désigne l'espace tout entier ; ainsi, la notation
$z_0 \in \Ker(d_0) \cap
\ldots \cap
\Ker(d_{i - 1})$
lorsque $i = 0$ signifie simplement $z_0 \in U_n$, sans condition.

\medskip\noindent
Procédons par récurrence croissante sur $i$. Pour $i = 0$, cela résulte
de la construction de $x_0$ et $z_0$. Prouvons-le pour
$0 < i \leq n$, en supposant le résultat établi pour $i - 1$. Tout d'abord,
$d_i(x_i) = 0$ par construction. Fixons donc $j$ avec $0 \leq j < i$.
Par récurrence, $d_j(x_{i - 1}) = 0$. Ainsi
$$
0 = d_j(x_{i - 1})
= d_j(x_i) + d_j(s_i(z_i))
= d_j(x_i) + s_{i - 1}(d_j(z_i)).
$$
La dernière égalité résulte des relations de la remarque
\ref{remark-relations}. Ces relations impliquent également que
$d_{i - 1}(d_j(x_i)) = d_j(d_i(x_i)) = 0$
car $d_i(x_i)= 0$ par construction. L'unicité dans la remarque générale
ci-dessus montre alors que l'égalité
$0 = x' + x'' = d_j(x_i) + s_{i - 1}(d_j(z_i))$
ne peut avoir lieu que si les deux termes sont nuls. On en conclut que
$d_j(x_i) = 0$ et, par injectivité de $s_{i - 1}$, que
$d_j(z_i) = 0$. Cela démontre l'affirmation.

\medskip\noindent
L'affirmation montre que l'on peut écrire de manière unique
$$
x = s_0(z_0) + s_1(z_1) + \ldots + s_n(z_n) + x_n
$$
avec $x_n \in N(U_{n + 1})$ et
$z_i \in \Ker(d_0) \cap \ldots \cap \Ker(d_{i - 1})$.
Autrement dit, on a obtenu une décomposition en somme directe
$$
U_{n + 1}
=
N(U_{n + 1})
\oplus
\bigoplus\nolimits_{i = 0}^{i = n}
s_i\Big(\Ker(d_0) \cap \ldots \cap \Ker(d_{i - 1})\Big)
$$
ayant la propriété suivante :
$$
\Ker(d_0) \cap \ldots \cap \Ker(d_j)
=
N(U_{n + 1}) \oplus
\bigoplus\nolimits_{i = j + 1}^{i = n}
s_i\Big(\Ker(d_0) \cap \ldots \cap \Ker(d_{i - 1})\Big)
$$
pour $j = 0, \ldots, n$. Le résultat s'en déduit comme suit. Chacun des
$z_i$ apparaissant dans l'expression de $x$ s'écrit de manière unique
$$
z_i = s_i(z'_{i, i}) + \ldots + s_{n - 1}(z'_{i, n - 1}) + z_{i, 0}
$$
avec $z_{i, 0} \in N(U_n)$ et
$z'_{i, j} \in \Ker(d_0) \cap \ldots \cap \Ker(d_{j - 1})$.
Les premières étapes de la décomposition de $z_i$ sont nulles puisque
$z_i$ appartient déjà aux noyaux de $d_0, \ldots, d_i$. On en déduit de
manière unique
$$
x = x_n + s_0(z_{0, 0}) + s_1(z_{1, 0}) + \ldots + s_n(z_{n, 0}) +
\sum\nolimits_{0 \leq i \leq j \leq n - 1} s_i(s_j(z'_{i, j})).
$$
En poursuivant ainsi, on obtient finalement une décomposition de $x$ en
somme de termes de la forme
$$
s_{i_1} s_{i_2} \ldots s_{i_k} (z)
$$
avec $0 \leq i_1 \leq i_2 \leq \ldots \leq i_k \leq n - k + 1$ et
$z \in N(U_{n + 1 - k})$. C'est précisément la décomposition requise,
car toute application surjective $[n + 1] \to [n + 1 - k]$ s'écrit de
manière unique sous la forme
$$
\sigma^{n - k}_{i_k} \ldots \sigma^{n - 1}_{i_2} \sigma^n_{i_1}
$$
avec $0 \leq i_1 \leq i_2 \leq \ldots \leq i_k \leq n - k + 1$.
\end{proof}

\begin{lemma}
\label{lemma-splitting-abelian-category}
Soit $\mathcal{A}$ une catégorie abélienne et soit $U$ un objet
simplicial de $\mathcal{A}$. Alors $U$ possède un scindement obtenu en
prenant $N(U_0) = U_0$ et, pour $m \geq 1$,
$$
N(U_m) = \bigcap\nolimits_{i = 0}^{m - 1} \Ker(d^m_i).
$$
De plus, ce scindement est fonctoriel sur la catégorie des objets
simpliciaux de $\mathcal{A}$.
\end{lemma}

\begin{proof}
Pour tout objet $A$ de $\mathcal{A}$, on obtient un groupe abélien
simplicial $\Mor_\mathcal{A}(A, U)$. Chacun d'eux est canoniquement
scindé par le lemme \ref{lemma-splitting-simplicial-groups}. De plus,
$$
N(\Mor_\mathcal{A}(A, U_m)) =
\bigcap\nolimits_{i = 0}^{m - 1} \Ker(d^m_i) =
\Mor_\mathcal{A}(A, N(U_m)).
$$
Ainsi, le morphisme (\ref{equation-splitting}) devient un isomorphisme
après application du foncteur $\Mor_\mathcal{A}(A, -)$, pour tout objet
de $\mathcal{A}$. C'est donc un isomorphisme par le lemme de Yoneda.
\end{proof}

\begin{lemma}
\label{lemma-injective-map-simplicial-abelian}
\begin{slogan}
Le foncteur de normalisation de Dold--Kan reflète
l'injectivité, la surjectivité et les isomorphismes.
\end{slogan}
Soit $\mathcal{A}$ une catégorie abélienne et soit $f : U \to V$ un
morphisme d'objets simpliciaux de $\mathcal{A}$. Si les morphismes
induits $N(f)_i : N(U)_i \to N(V)_i$ sont injectifs pour tout $i$,
alors $f_i$ est injectif pour tout $i$. Le même énoncé vaut en remplaçant
« injectif » par « surjectif » ou par « isomorphisme ».
\end{lemma}

\begin{proof}
Cela résulte immédiatement du lemme
\ref{lemma-splitting-abelian-category} et de la définition d'un scindement.
\end{proof}


\begin{lemma}
\label{lemma-N-d-in-N}
Soit $\mathcal{A}$ une catégorie abélienne, soit $U$ un objet simplicial
de $\mathcal{A}$ et soient $N(U_m)$ comme dans le lemme
\ref{lemma-splitting-abelian-category}. Alors
$d^m_m(N(U_m)) \subset N(U_{m - 1})$.
\end{lemma}

\begin{proof}
Pour $j = 0, \ldots, m - 2$, on a
$d^{m - 1}_j d^m_m = d^{m - 1}_{m - 1} d^m_j$
d'après les relations de la remarque \ref{remark-relations}.
Le résultat en découle.
\end{proof}

\begin{lemma}
\label{lemma-simplicial-abelian-n-skel-sub}
Soit $\mathcal{A}$ une catégorie abélienne, soit $U$ un objet simplicial
de $\mathcal{A}$ et soit $n \geq 0$ un entier. La règle
$$
U'_m = \sum\nolimits_{\varphi : [m] \to [i], \ i\leq n} \Im(U(\varphi))
$$
définit un sous-objet simplicial $U' \subset U$ tel que $U'_i = U_i$
pour $i \leq n$. De plus, $N(U'_m) = 0$ pour tout $m > n$.
\end{lemma}

\begin{proof}
Fixons $m$, $i \leq n$ et un morphisme $\varphi : [m] \to [i]$.
L'image par $U(\psi)$ de $\Im(U(\varphi))$, pour tout
$\psi : [m'] \to [m]$, est égale à l'image de
$U(\varphi \circ \psi)$, où
$\varphi \circ \psi : [m'] \to [i]$. Ainsi $U'$ est un objet
simplicial. Fixons $m > n$. Il faut montrer que $N(U'_m) = 0$.
Par définition de $N(U_m)$ et $N(U'_m)$, on a
$N(U'_m) = U'_m \cap N(U_m)$ (intersection de sous-objets).
Comme $U$ est scindé par le lemme
\ref{lemma-splitting-abelian-category}, il suffit de montrer que $U'_m$
est contenu dans la somme
$$
\sum\nolimits_{\varphi : [m] \to [m']\text{ surjective}, \ m' < m}
\Im(U(\varphi)|_{N(U_{m'})}).
$$
Par le scindement, chaque $U_{m'}$ est la somme des images des
$N(U_{m''})$ par $U(\psi)$, pour les applications surjectives
$\psi : [m'] \to [m'']$. La somme affichée ci-dessus est donc égale à
$$
\sum\nolimits_{\varphi : [m] \to [m']\text{ surjective}, \ m' < m}
\Im(U(\varphi)).
$$
Manifestement, $U'_m$ y est contenu : toute application
$\varphi : [m] \to [i]$, $i \leq n$, intervenant dans la définition de
$U'_m$ se factorise sous la forme $[m] \to [m'] \to [i]$, où
$[m] \to [m']$ est surjective et $m' < m$, comme dans la dernière somme
affichée ci-dessus.
\end{proof}




\section{Foncteurs cosquelette}
\label{section-coskeleton}

\noindent
Soit $\mathcal{C}$ une catégorie. Le {\it foncteur cosquelette}
(s'il existe) est un foncteur
$$
\text{cosk}_n :
\text{Simp}_n(\mathcal{C}) \longrightarrow \text{Simp}(\mathcal{C})
$$
adjoint à droite du foncteur squelette. En formule,
\begin{equation}
\label{equation-cosk}
\Mor_{\text{Simp}(\mathcal{C})}(U, \text{cosk}_n V)
=
\Mor_{\text{Simp}_n(\mathcal{C})}(\text{sk}_n U, V)
\end{equation}
Étant donné un objet simplicial $n$-tronqué $V$, on dit que
{\it $\text{cosk}_nV$ existe} s'il existe un objet
$\text{cosk}_nV \in \Ob(\text{Simp}(\mathcal{C}))$ et un morphisme
$\text{sk}_n \text{cosk}_n V \to V$ tels que la formule affichée soit
vérifiée, autrement dit si le foncteur
$U \mapsto \Mor_{\text{Simp}_n(\mathcal{C})}(\text{sk}_n U, V)$
est représentable. S'il existe, il est unique à isomorphisme unique près
par le lemme de Yoneda. Voir Categories, section
\ref{categories-section-opposite}.

\begin{example}
\label{example-cosk0}
Supposons que la catégorie $\mathcal{C}$ possède les produits finis non
vides d'un objet avec lui-même. Un objet simplicial $0$-tronqué de
$\mathcal{C}$ est simplement un objet $X$ de $\mathcal{C}$. Dans ce cas,
nous affirmons que $\text{cosk}_0(X)$ est l'objet simplicial $U$ tel que
$U_n = X^{n + 1}$ soit le produit de $(n + 1)$ copies de $X$, muni de la
structure d'objet simplicial de l'exemple
\ref{example-fibre-products-simplicial-object}. En effet, un morphisme
$V \to U$, où $V$ est un objet simplicial, est donné par des morphismes
$V_n \to X^{n + 1}$ tels que tous les diagrammes
$$
\xymatrix{
V_n \ar[r] \ar[d]_{V([0] \to [n], 0 \mapsto i)} &
X^{n + 1} \ar[d]^{\text{pr}_i} \\
V_0 \ar[r] &
X
}
$$
soient commutatifs. Cela signifie manifestement que l'application
détermine, et est déterminée par, un unique morphisme $V_0 \to X$.
La formule (\ref{equation-cosk}) est donc vérifiée.
\end{example}

\noindent
Rappelons la catégorie $\Delta/[n]$, voir l'exemple
\ref{example-simplex-category}. Notons $(\Delta/[n])_{\leq m}$ la
sous-catégorie pleine de $\Delta/[n]$ formée des objets
$[k] \to [n]$ de $\Delta/[n]$ avec $k \leq m$. Autrement dit, on a le diagramme
commutatif suivant de catégories et de foncteurs :
$$
\xymatrix{
(\Delta/[n])_{\leq m} \ar[r] \ar[d] &
\Delta/[n] \ar[d] \\
\Delta_{\leq m} \ar[r] &
\Delta
}
$$
Étant donné un objet simplicial $m$-tronqué $U$ de $\mathcal{C}$,
on définit un foncteur
$$
U(n) : (\Delta/[n])_{\leq m}^{opp} \longrightarrow \mathcal{C}
$$
par les règles
\begin{eqnarray*}
([k] \to [n]) & \longmapsto & U_k \\
\psi : ([k'] \to [n]) \to ([k] \to [n]) &
\longmapsto &
U(\psi) : U_k \to U_{k'}
\end{eqnarray*}
À tout morphisme $\varphi : [n] \to [n']$ de $\Delta$ est associé un foncteur
$$
\overline{\varphi} :
(\Delta/[n])_{\leq m}
\longrightarrow
(\Delta/[n'])_{\leq m}
$$
qui envoie $\alpha : [k] \to [n]$ sur
$\varphi \circ \alpha : [k] \to [n']$. Le composé
$U(n') \circ \overline{\varphi}$ est égal au foncteur $U(n)$.

\begin{lemma}
\label{lemma-existence-cosk}
Si la catégorie $\mathcal{C}$ possède les limites finies, alors les
foncteurs $\text{cosk}_m$ existent pour tout $m$. De plus, pour tout
objet simplicial $m$-tronqué $U$, l'objet simplicial
$\text{cosk}_mU$ est décrit par la formule
$$
(\text{cosk}_mU)_n = \lim_{(\Delta/[n])_{\leq m}^{opp}} U(n)
$$
et, pour $\varphi : [n] \to [n']$, l'application
$\text{cosk}_mU(\varphi)$ provient de l'identification
$U(n') \circ \overline{\varphi} = U(n)$ ci-dessus, par Categories,
lemme \ref{categories-lemma-functorial-limit}.
\end{lemma}

\begin{proof}
Dans cette démonstration, notons $\text{cosk}_mU$ l'objet simplicial tel
que $(\text{cosk}_mU)_n$ soit égal à
$\lim_{(\Delta/[n])_{\leq m}^{opp}} U(n)$. Nous conclurons à la fin que
celui-ci satisfait bien à la propriété universelle requise.

\medskip\noindent
Soit $V$ un objet simplicial. Un morphisme
$\gamma : V \to \text{cosk}_mU$ est donné par une suite de morphismes
$\gamma_n : V_n \to (\text{cosk}_mU)_n$. Par définition d'une limite,
ces données reviennent à une famille de morphismes
$\gamma(\alpha) : V_n \to U_k$, où $\alpha$ parcourt les applications
$\alpha : [k] \to [n]$ avec $k \leq m$. Ces morphismes satisfont en outre
à la condition que les diagrammes
$$
\xymatrix{
V_n \ar[r]_{\gamma(\alpha)} &  U_k \\
V_{n'} \ar[r]^{\gamma(\alpha')} \ar[u]^{V(\varphi)} & U_{k'} \ar[u]_{U(\psi)}
}
$$
soient commutatifs, quels que soient $0 \leq k, k' \leq m$,
$0 \leq n, n'$ et $\psi : [k] \to [k']$,
$\varphi : [n] \to [n']$, $\alpha : [k] \to [n]$ et
$\alpha' : [k'] \to [n']$ dans $\Delta$, avec
$\varphi \circ \alpha = \alpha' \circ \psi$. En prenant
$n = k = k'$, $\varphi = \alpha'$ et
$\alpha = \psi = \text{id}_{[k]}$, on en déduit que
$\gamma(\alpha') = \gamma(\text{id}_{[k]}) \circ V(\alpha')$.
Autrement dit, les morphismes $\gamma(\text{id}_{[k]})$, $k \leq m$,
déterminent $\gamma$. Il est aisé de voir qu'ils forment un morphisme
$\text{sk}_m V \to U$.

\medskip\noindent
Réciproquement, étant donné un morphisme
$\gamma : \text{sk}_m V \to U$, on obtient une famille de morphismes
$\gamma(\alpha)$, où $\alpha$ parcourt les applications
$\alpha : [k] \to [n]$ avec $k \leq m$, en posant
$\gamma(\alpha) = \gamma(\text{id}_{[k]}) \circ V(\alpha)$.
Ces morphismes satisfont à toutes les conditions de commutativité
affichées ci-dessus et donnent donc un morphisme
$V \to \text{cosk}_m U$.
\end{proof}

\begin{lemma}
\label{lemma-trivial-cosk}
Soit $\mathcal{C}$ une catégorie et soit $U$ un objet simplicial
$m$-tronqué de $\mathcal{C}$. Pour $n \leq m$, la limite
$\lim_{(\Delta/[n])_{\leq m}^{opp}} U(n)$ existe et est canoniquement
isomorphe à $U_n$.
\end{lemma}

\begin{proof}
En effet, la catégorie $(\Delta/[n])_{\leq m}$ possède alors un objet
final, à savoir l'application identité $[n] \to [n]$.
\end{proof}

\begin{lemma}
\label{lemma-recover-cosk}
Soit $\mathcal{C}$ une catégorie possédant les limites finies et soit
$U$ un objet simplicial $n$-tronqué de $\mathcal{C}$. Le morphisme
$\text{sk}_n \text{cosk}_n U \to U$ est un isomorphisme.
\end{lemma}

\begin{proof}
Combiner les lemmes \ref{lemma-existence-cosk} et
\ref{lemma-trivial-cosk}.
\end{proof}

\noindent
Décrivons plus en détail un cas particulier du foncteur cosquelette.
Par abus de notation, nous noterons encore $\text{sk}_n$ le foncteur de
restriction
$\text{Simp}_{n'}(\mathcal{C}) \to \text{Simp}_n(\mathcal{C})$
pour tout $n' \geq n$. Nous allons décrire un adjoint à droite du foncteur
$\text{sk}_n : \text{Simp}_{n + 1}(\mathcal{C})
\to \text{Simp}_n(\mathcal{C})$.
Pour $n \geq 1$ et $0 \leq i < j \leq n + 1$, définissons
$\delta^{n + 1}_{i, j} : [n - 1] \to [n + 1]$ comme l'application
croissante qui omet $i$ et $j$. Remarquons que
$\delta^{n + 1}_{i, j} =
\delta^{n + 1}_j \circ \delta^n_i =
\delta^{n + 1}_i \circ \delta^n_{j - 1}$, see
d'après le lemme \ref{lemma-relations-face-degeneracy}. Cela motive le
lemme suivant.

\begin{lemma}
\label{lemma-formula-limit}
Soit $n$ un entier $\geq 1$ et soit $U$ un objet simplicial $n$-tronqué
de $\mathcal{C}$. Considérons le foncteur contravariant de $\mathcal{C}$
vers $\textit{Sets}$ qui associe à un objet $T$ l'ensemble
$$
\{ (f_0, \ldots, f_{n + 1}) \in \Mor_\mathcal{C}(T, U_n)
\mid
d^n_{j - 1} \circ f_i = d^n_i \circ f_j
\ \forall\ 0\leq i < j\leq n + 1\}
$$
Si ce foncteur est représentable par un objet $U_{n + 1}$ de
$\mathcal{C}$, alors
$$
U_{n + 1} = \lim_{(\Delta/[n + 1])_{\leq n}^{opp}} U(n)
$$
\end{lemma}

\begin{proof}
La limite, si elle existe, représente le foncteur qui associe à un objet
$T$ l'ensemble
$$
\{
(f_\alpha)_{\alpha : [k] \to [n + 1], k \leq n}
\mid
f_{\alpha \circ \psi} = U(\psi) \circ f_\alpha\ \forall
\ \psi : [k'] \to [k], \alpha : [k] \to [n + 1]
\}.
$$
Nous allons en fait montrer que ce foncteur est isomorphe à celui qui
figure dans l'énoncé. L'application dans un sens est donnée par la règle
$$
(f_\alpha)_{\alpha}
\longmapsto
(f_{\delta^{n + 1}_0}, \ldots, f_{\delta^{n + 1}_{n + 1}}).
$$
Cette famille satisfait aux conditions du lemme, car
$$
d^n_{j - 1} \circ f_{\delta^{n + 1}_i} =
f_{\delta^{n + 1}_i \circ \delta^n_{j - 1}} =
f_{\delta^{n + 1}_j \circ \delta^n_i} =
d^n_i \circ f_{\delta^{n + 1}_j}
$$
par les relations rappelées avant le lemme. Pour construire une
application dans l'autre sens, il faut associer à un système
$(f_0, \ldots, f_{n + 1})$ comme dans la formule du lemme un système
d'applications $f_\alpha$. Soit $\alpha : [k] \to [n + 1]$. Comme
$k \leq n$, l'application $\alpha$ n'est pas surjective. On peut donc
écrire $\alpha = \delta^{n + 1}_i \circ \psi$ pour un certain
$0 \leq i \leq n + 1$ et un certain $\psi : [k] \to [n]$. On doit poser
$$
f_\alpha = U(\psi) \circ f_i.
$$
Il faut vérifier que cette définition ne dépend pas du choix du couple
$(i, \psi)$. D'abord, $i$ étant fixé, il existe un unique $\psi$ qui
convienne. Ensuite, soit $(j, \phi)$ un autre couple. Alors
$i \not = j$ et l'on peut supposer $i < j$. Puisqu'aucun des indices
$i, j$ n'appartient à l'image de $\alpha$, on peut écrire
$\alpha = \delta^{n + 1}_{i, j} \circ \xi$ ; on a alors
$\psi = \delta^n_{j - 1} \circ \xi$ et
$\phi = \delta^n_i \circ \xi$. Ainsi
\begin{eqnarray*}
U(\psi) \circ f_i & = & U(\delta^n_{j - 1} \circ \xi) \circ f_i \\
& = & U(\xi) \circ d^n_{j - 1} \circ f_i \\
& = & U(\xi) \circ d^n_i \circ f_j \\
& = & U(\delta^n_i \circ \xi) \circ f_j \\
& = & U(\phi) \circ f_j
\end{eqnarray*}
comme voulu. Il reste à vérifier que les applications $f_\alpha$ ainsi
définies satisfont aux conditions imposées à un système
$(f_\alpha)_\alpha$. Pour cela, supposons que
$\psi : [k'] \to [k]$, $\alpha : [k] \to [n + 1]$, avec
$k, k' \leq n$. Posons $\alpha' = \alpha \circ \psi$.
Choisissons $i$ hors de l'image de $\alpha$. Alors $i$ n'appartient pas
non plus à l'image de $\alpha'$. Écrivons
$\alpha = \delta^{n + 1}_i \circ \phi$ (on ne peut employer la lettre
$\psi$ ici, puisqu'elle est déjà utilisée). On a manifestement
$\alpha' = \delta^{n + 1}_i \circ \phi \circ \psi$. Par la construction
ci-dessus,
$$
U(\psi) \circ f_\alpha = U(\psi) \circ U(\phi) \circ f_i
= U(\phi \circ \psi) \circ f_i = f_{\alpha \circ \psi} = f_{\alpha'}
$$
comme voulu. Nous laissons au lecteur l'agréable tâche de vérifier que
nos constructions sont des bijections inverses l'une de l'autre et sont
fonctorielles en $T$.
\end{proof}

\begin{lemma}
\label{lemma-work-out}
Soit $n$ un entier $\geq 1$ et soit $U$ un objet simplicial $n$-tronqué
de $\mathcal{C}$. Considérons le foncteur contravariant de $\mathcal{C}$
vers $\textit{Sets}$ qui associe à un objet $T$ l'ensemble
$$
\{ (f_0, \ldots, f_{n + 1}) \in \Mor_\mathcal{C}(T, U_n)
\mid
d^n_{j - 1} \circ f_i = d^n_i \circ f_j
\ \forall\ 0\leq i < j\leq n + 1\}
$$
Si ce foncteur est représentable par un objet $U_{n + 1}$ de
$\mathcal{C}$, alors il existe un objet simplicial $(n + 1)$-tronqué
$\tilde U$, avec $\text{sk}_n \tilde U = U$ et
$\tilde U_{n + 1} = U_{n + 1}$, tel que l'adjonction suivante soit vérifiée :
$$
\Mor_{\text{Simp}_{n + 1}(\mathcal{C})}(V, \tilde U)
=
\Mor_{\text{Simp}_n(\mathcal{C})}(\text{sk}_nV, U)
$$
\end{lemma}

\begin{proof}
Par le lemme \ref{lemma-trivial-cosk}, on a des identifications
$$
U_i = \lim_{(\Delta/[i])_{\leq n}^{opp}} U(i)
$$
pour $0 \leq i \leq n$. Par le lemme \ref{lemma-formula-limit}, on a
$$
U_{n + 1} = \lim_{(\Delta/[n + 1])_{\leq n}^{opp}} U(n).
$$
On peut donc définir, pour tout $\varphi : [i] \to [j]$ avec
$i, j \leq n + 1$, l'application correspondante
$\tilde U(\varphi) : \tilde U_j \to \tilde U_i$ exactement comme dans le
lemme \ref{lemma-existence-cosk}. On obtient ainsi un objet simplicial
$(n + 1)$-tronqué $\tilde U$ avec $\text{sk}_n \tilde U = U$.

\medskip\noindent
Pour établir l'adjonction, raisonnons comme suit. Étant donné un élément
$\gamma : \text{sk}_n V \to U$ du membre de droite, considérons les morphismes
$f_i = \gamma_n \circ d^{n + 1}_i : V_{n + 1} \to V_n \to U_n$.
Ils satisfont manifestement aux relations
$d^n_{j - 1} \circ f_i = d^n_i \circ f_j$ et définissent donc un unique
morphisme $V_{n + 1} \to U_{n + 1}$ par le choix de $U_{n + 1}$.
Réciproquement, étant donné un morphisme $\gamma' : V \to \tilde U$ du
membre de gauche, il suffit de le restreindre à $\Delta_{\leq n}$ pour
obtenir un élément du membre de droite. Nous laissons au lecteur le soin
de vérifier que ces constructions sont inverses l'une de l'autre.
\end{proof}

\begin{remark}
\label{remark-explicit-face-degeneracy}
Soient $U$ et $U_{n + 1}$ comme dans le lemme \ref{lemma-work-out}.
Sur les points à valeurs dans $T$, on décrit aisément les applications
de face et de dégénérescence de $\tilde U$. Explicitement, les
applications $d^{n + 1}_i : U_{n + 1} \to U_n$ sont données par
$$
(f_0, \ldots, f_{n + 1}) \longmapsto f_i.
$$
Et les applications $s^n_j : U_n \to U_{n + 1}$ sont données par
\begin{eqnarray*}
f & \longmapsto & (
s^{n - 1}_{j - 1} \circ d^{n - 1}_0 \circ f, \\
& &
s^{n - 1}_{j - 1} \circ d^{n - 1}_1 \circ f, \\
& &
\ldots\\
& &
s^{n - 1}_{j - 1} \circ d^{n - 1}_{j - 1} \circ f, \\
& &
f, \\
& &
f, \\
& &
s^{n - 1}_j \circ d^{n - 1}_{j + 1} \circ f, \\
& &
s^{n - 1}_j \circ d^{n - 1}_{j + 2} \circ f, \\
& &
\ldots\\
& &
s^{n - 1}_j \circ d^{n - 1}_n \circ f
)
\end{eqnarray*}
où nous laissons au lecteur le soin de vérifier que le membre de droite
appartient à l'ensemble affiché dans le lemme \ref{lemma-work-out}.
Pour $n = 0$, il y a une application, à savoir $f \mapsto (f, f)$.
Pour $n = 1$, il y en a deux, à savoir
$f \mapsto (f, f, s_0d_1f)$ et
$f \mapsto (s_0d_0f, f, f)$.
Pour $n = 2$, il y en a trois, à savoir
$f \mapsto (f, f, s_0d_1f, s_0d_2f)$,
$f \mapsto (s_0d_0f, f, f, s_1d_2f)$, et
$f \mapsto (s_1d_0f, s_1d_1f, f, f)$.
Et ainsi de suite.
\end{remark}

\begin{remark}
\label{remark-cosk-simplicial-sets}
Dans le cas des ensembles simpliciaux, la construction du lemme
\ref{lemma-work-out} est la suivante. Étant donné un ensemble simplicial
$n$-tronqué $U$, on construit canoniquement un ensemble simplicial
$(n + 1)$-tronqué $\tilde U$. On adjoint un ensemble de
$(n + 1)$-simplexes $U_{n + 1}$ par la formule du lemme. Un élément de
$U_{n + 1}$
est donc une famille numérotée $(f_0, \ldots, f_{n + 1})$ de
$n$-simplexes, qui se recollent comme les faces d'un
$(n + 1)$-simplexe. Autrement dit, la $i$-ième face de $f_j$ est la
$(j-1)$-ième face de $f_i$ pour $i < j$. Géométriquement, les
applications de face et de dégénérescence de $\tilde U$ sont évidentes.
Si $V$ est un ensemble simplicial $(n + 1)$-tronqué, ses
$(n + 1)$-simplexes fournissent des familles compatibles de $n$-simplexes
$(f_0, \ldots, f_{n + 1})$ avec $f_i \in V_n$. Il existe donc une
application naturelle
$\Mor(\text{sk}_nV, U) \to \Mor(V, \tilde U)$
inverse de l'application de restriction canonique dans l'autre sens.

\medskip\noindent
Il suffit d'ailleurs d'effectuer la combinatoire de la construction pour
les ensembles simpliciaux tronqués. En effet, pour tout objet $T$ de la
catégorie $\mathcal{C}$ et tout objet simplicial $n$-tronqué $U$ de
$\mathcal{C}$, on peut considérer l'ensemble simplicial $n$-tronqué
$\Mor(T, U)$. On lui applique la construction, on prend son ensemble de
$(n + 1)$-simplexes et on exige que celui-ci soit représentable. C'est une
bonne manière d'interpréter le lemme \ref{lemma-work-out}.
\end{remark}

\begin{remark}
\label{remark-inductive-coskeleton}
{\it Construction inductive des cosquelettes.}
Supposons que $\mathcal{C}$ soit une catégorie possédant les limites
finies et que $U$ soit un objet simplicial $m$-tronqué de $\mathcal{C}$.
On peut alors construire par récurrence des objets $n$-tronqués $U^n$
comme suit :
\begin{enumerate}
\item Pour commencer, poser $U^m = U$.
\item Étant donné $U^n$ pour $n \geq m$, poser
$U^{n + 1} = \tilde U^n$, où $\tilde U^n$ est construit à partir de
$U^n$ comme dans le lemme \ref{lemma-work-out}.
\end{enumerate}
Comme la construction du lemme \ref{lemma-work-out} laisse inchangé le
$n$-squelette de $U^n$, on peut définir $\text{cosk}_m U$ comme l'objet
simplicial tel que
$(\text{cosk}_m U)_n = U^n_n = U^{n + 1}_n = \ldots$.
Il résulte formellement du lemme \ref{lemma-work-out} que $U^n$ satisfait
à la formule
$$
\Mor_{\text{Simp}_n(\mathcal{C})}(V, U^n)
=
\Mor_{\text{Simp}_m(\mathcal{C})}(\text{sk}_mV, U)
$$
pour tout $n \geq m$. Il en résulte également, de manière formelle, que
$$
\Mor_{\text{Simp}(\mathcal{C})}(V, \text{cosk}_mU)
=
\Mor_{\text{Simp}_m(\mathcal{C})}(\text{sk}_mV, U)
$$
pour le choix de $\text{cosk}_mU$ ci-dessus.
\end{remark}

\begin{lemma}
\label{lemma-cosk-up}
Soit $\mathcal{C}$ une catégorie possédant les limites finies.
\begin{enumerate}
\item Pour tout $n$, le foncteur $\text{sk}_n : \text{Simp}(\mathcal{C})
\to \text{Simp}_n(\mathcal{C})$ possède un adjoint à droite
$\text{cosk}_n$.
\item Pour tout $n' \geq n$, le foncteur
$\text{sk}_n : \text{Simp}_{n'}(\mathcal{C}) \to \text{Simp}_n(\mathcal{C})$
possède un adjoint à droite, à savoir $\text{sk}_{n'}\text{cosk}_n$.
\item Pour tout $m \geq n \geq 0$ et tout objet simplicial $n$-tronqué
$U$ de $\mathcal{C}$, on a
$\text{cosk}_m \text{sk}_m \text{cosk}_n U = \text{cosk}_n U$.
\item Si $U$ est un objet simplicial de $\mathcal{C}$ tel que l'application
canonique
$U \to \text{cosk}_n \text{sk}_nU$
soit un isomorphisme pour un certain $n \geq 0$, alors l'application canonique
$U \to \text{cosk}_m \text{sk}_mU$
est un isomorphisme pour tout $m \geq n$.
\end{enumerate}
\end{lemma}

\begin{proof}
L'existence dans (1) résulte du lemme \ref{lemma-existence-cosk}.
Les assertions (2) et (3) découlent de la discussion de la remarque
\ref{remark-inductive-coskeleton}. L'assertion (4) est alors évidente.
\end{proof}

\begin{remark}
\label{remark-existence-cosk}
Toutes les limites finies ne sont pas nécessaires pour définir les
foncteurs cosquelette. Voici quelques remarques.
\begin{enumerate}
\item On a vu dans l'exemple \ref{example-cosk0} que, si $\mathcal{C}$
possède les produits de deux objets, alors $\text{cosk}_0$ existe.
\item Pour $k > 0$, le foncteur $\text{cosk}_k$ existe si
$\mathcal{C}$ possède les limites finies connexes.
\end{enumerate}
Cela ressort de la construction inductive des cosquelettes (remarques
\ref{remark-cosk-simplicial-sets} et
\ref{remark-inductive-coskeleton}), mais résulte aussi du fait que les
catégories $(\Delta/[n])_{\leq k}$, pour $k \geq 1$ et $n \geq k + 1$,
utilisées dans le lemme \ref{lemma-existence-cosk}, sont connexes.
Pour $n \leq k$, ces catégories ne sont pas nécessaires, d'après les
lemmes \ref{lemma-trivial-cosk} et \ref{lemma-recover-cosk}.
(Lorsque $k$ croît, les catégories $(\Delta/[n])_{\leq k}$ pour
$k \geq 1$ et $n \geq k + 1$ deviennent de plus en plus connexes au sens
topologique.)
\end{remark}

\begin{lemma}
\label{lemma-cosk-product}
Soient $U$, $V$ des objets simpliciaux $n$-tronqués d'une catégorie
$\mathcal{C}$. Alors
$$
\text{cosk}_n (U \times V) = \text{cosk}_nU \times \text{cosk}_nV
$$
dès que les membres de gauche et de droite existent.
\end{lemma}

\begin{proof}
Soit $W$ un objet simplicial. On a
\begin{eqnarray*}
\Mor(W, \text{cosk}_n (U \times V))
& = &
\Mor(\text{sk}_n W, U \times V) \\
& = &
\Mor(\text{sk}_n W, U)
\times
\Mor(\text{sk}_nW, V) \\
& = &
\Mor(W, \text{cosk}_n U)
\times
\Mor(W, \text{cosk}_n V) \\
& = &
\Mor(W, \text{cosk}_n U \times \text{cosk}_n V)
\end{eqnarray*}
Le lemme en découle.
\end{proof}

\begin{lemma}
\label{lemma-cosk-fibre-product}
Supposons que $\mathcal{C}$ possède les produits fibrés. Soient
$U \to V$ et $W \to V$ des morphismes d'objets simpliciaux $n$-tronqués
de la catégorie $\mathcal{C}$. Alors
$$
\text{cosk}_n (U \times_V W)
=
\text{cosk}_nU \times_{\text{cosk}_n V} \text{cosk}_nW
$$
dès que les membres de gauche et de droite existent.
\end{lemma}

\begin{proof}
La démonstration est omise ; elle est très semblable à celle du lemme
\ref{lemma-cosk-product}.
\end{proof}

\begin{lemma}
\label{lemma-cosk-above-object}
Soit $\mathcal{C}$ une catégorie possédant les limites finies et soit
$X \in \Ob(\mathcal{C})$. Le foncteur $\mathcal{C}/X \to \mathcal{C}$
commute aux foncteurs cosquelette $\text{cosk}_k$ pour $k \geq 1$.
\end{lemma}

\begin{proof}
Cela signifie que, si $U$ est un objet simplicial de $\mathcal{C}/X$,
que l'on peut voir comme un objet simplicial de $\mathcal{C}$ muni d'un
morphisme vers l'objet simplicial constant $X$, alors
$\text{cosk}_k U$ calculé dans $\mathcal{C}/X$ est le même que celui
calculé dans $\mathcal{C}$. Cela résulte, par exemple, de Categories,
lemme \ref{categories-lemma-connected-limit-over-X}, car les catégories
$(\Delta/[n])_{\leq k}$ pour $k \geq 1$ et $n \geq k + 1$ utilisées dans
le lemme \ref{lemma-existence-cosk} sont connexes. Pour $n \leq k$, ces
catégories ne sont pas nécessaires, d'après les lemmes
\ref{lemma-trivial-cosk} et \ref{lemma-recover-cosk}.
\end{proof}

\begin{lemma}
\label{lemma-simplex-cosk}
L'application canonique
$\Delta[n] \to \text{cosk}_1 \text{sk}_1 \Delta[n]$
est un isomorphisme.
\end{lemma}

\begin{proof}
Considérons un ensemble simplicial $U$ et un morphisme
$f : U \to \Delta[n]$. Celui-ci associe à chaque $u \in U_i$ une
application $f_u : [i] \to [n]$ de $\Delta$. Ces applications doivent en
outre vérifier $f_u \circ \varphi = f_{U(\varphi)(u)}$ pour tout
$\varphi : [j] \to [i]$. Notons $\epsilon^i_j : [0] \to [i]$
l'application qui envoie $0$ sur $j$, et notons $F : U_0 \to [n]$
l'application $u \mapsto f_u(0)$. Alors
$$
f_u(j) = F(\epsilon^i_j(u))
$$
pour tout $0 \leq j \leq i$ et tout $u \in U_i$. En particulier, la
connaissance de $F$ détermine les applications $f_u$ pour tout
$u\in U_i$ et tout $i$. Réciproquement, étant donnée une application
$F : U_0 \to [n]$, on peut poser, pour tout $i$, tout $u \in U_i$ et
tout $0 \leq j \leq i$,
$$
f_u(j) = F(\epsilon^i_j(u))
$$
Cela ne définit pas en général un morphisme $f$ d'ensembles simpliciaux
comme ci-dessus. La condition requise est que toutes les applications
$f_u$ soient croissantes. Elle équivaut manifestement à
$F(\epsilon^i_j(u)) \leq F(\epsilon^i_{j'}(u))$ dès que
$0 \leq j \leq j' \leq i$ et $u \in U_i$. Or les morphismes
$$
\epsilon^i_j, \epsilon^i_{j'} : [0] \to [i]
$$
se factorisent tous deux par l'application
$\epsilon^i_{j, j'} : [1] \to [i]$ définie par
$0 \mapsto j$, $1 \mapsto j'$. Il suffit donc de vérifier les inégalités
pour $i = 1$ et $u \in U_1$. Autrement dit,
$$
\Mor(U, \Delta[n])
=
\Mor(\text{sk}_1 U, \text{sk}_1 \Delta[n])
$$
comme voulu.
\end{proof}



\section{Augmentations}
\label{section-augmentation}

\begin{definition}
\label{definition-augmentation}
Soit $\mathcal{C}$ une catégorie et soit $U$ un objet simplicial de
$\mathcal{C}$. Une {\it augmentation $\epsilon : U \to X$ de $U$ vers
un objet $X$ de $\mathcal{C}$} est un morphisme de $U$ vers l'objet
simplicial constant $X$.
\end{definition}

\begin{lemma}
\label{lemma-augmentation-howto}
Soit $\mathcal{C}$ une catégorie, soit $X \in \Ob(\mathcal{C})$ et soit
$U$ un objet simplicial de $\mathcal{C}$. Se donner une augmentation de
$U$ vers $X$ revient à se donner un morphisme $\epsilon_0 : U_0 \to X$
tel que $\epsilon_0 \circ d^1_0 = \epsilon_0 \circ d^1_1$.
\end{lemma}

\begin{proof}
Un morphisme $\epsilon : U \to X$ fournit évidemment un
$\epsilon_0$ comme dans le lemme. Réciproquement, étant donné
$\epsilon_0$ comme dans le lemme, définissons
$\epsilon_n : U_n \to X$ en choisissant un morphisme quelconque
$\alpha : [0] \to [n]$ et en posant
$\epsilon_n = \epsilon_0 \circ U(\alpha)$. En effet, si
$\beta : [0] \to [n]$ est un autre choix, il existe un morphisme
$\gamma : [1] \to [n]$ tel que $\alpha$ et $\beta$ se factorisent tous
deux sous la forme $[0] \to [1] \to [n]$. La condition sur
$\epsilon_0$ montre donc que $\epsilon_n$ est bien définie. On vérifie
alors aisément que $(\epsilon_n) : U \to X$ est un morphisme d'objets
simpliciaux.
\end{proof}

\begin{lemma}
\label{lemma-cosk-minus-one}
Soit $\mathcal{C}$ une catégorie possédant les produits fibrés, soit
$f : Y\to X$ un morphisme de $\mathcal{C}$ et soit $U$ l'objet
simplicial de $\mathcal{C}$ dont le terme de degré $n$ est le produit
fibré de $(n + 1)$ copies
$Y \times_X Y \times_X \ldots \times_X Y$. Voir l'exemple
\ref{example-fibre-products-simplicial-object}. Pour tout objet simplicial
$V$ de $\mathcal{C}$, on a
\begin{align*}
\Mor_{\text{Simp}(\mathcal{C})}(V, U)
& =
\Mor_{\text{Simp}_1(\mathcal{C})}(\text{sk}_1 V, \text{sk}_1 U) \\
& =
\{g_0 : V_0 \to Y \mid f \circ g_0 \circ d^1_0 = f \circ g_0 \circ d^1_1\}
\end{align*}
En particulier, $U = \text{cosk}_1 \text{sk}_1 U$.
\end{lemma}

\begin{proof}
Supposons que $g : \text{sk}_1V \to \text{sk}_1U$ soit un morphisme
d'objets simpliciaux $1$-tronqués. Alors le diagramme
$$
\xymatrix{
V_1 \ar@<1ex>[r]^{d^1_0} \ar@<-1ex>[r]_{d^1_1} \ar[d]_{g_1} &
V_0 \ar[d]^{g_0} \\
Y \times_X Y \ar@<1ex>[r]^{pr_1} \ar@<-1ex>[r]_{pr_0} &
Y \ar[r] & X
}
$$
est commutatif, ce qui démontre la relation du lemme. Réciproquement,
il faut montrer qu'un morphisme $g_0$ satisfaisant à la relation
$f \circ g_0 \circ d^1_0 = f \circ g_0 \circ d^1_1$ détermine un unique
morphisme d'objets simpliciaux $g : V \to U$. Pour tout $n \geq 1$, posons
$g_{n, i} = g_0 \circ V([0] \to [n], 0 \mapsto i) :
V_n \to Y$. L'égalité ci-dessus implique que
$f \circ g_{n, i} = f \circ g_{n, i + 1}$ en vertu du diagramme commutatif
$$
\xymatrix{
[0] \ar[rd]_{0 \mapsto 0} \ar[rrrrrd]^{0 \mapsto i} \\
& [1] \ar[rrrr]^{0 \mapsto i, 1\mapsto i + 1} & & & & [n] \\
[0] \ar[ru]^{0 \mapsto 1} \ar[rrrrru]_{0 \mapsto i + 1}
}
$$
On obtient donc
$(g_{n, 0}, \ldots, g_{n, n}) : V_n \to Y \times_X\ldots \times_X Y = U_n$.
Nous laissons au lecteur le soin de vérifier qu'il s'agit d'un morphisme
d'objets simpliciaux. La dernière assertion du lemme équivaut à la
première égalité de sa formule affichée.
\end{proof}

\begin{remark}
\label{remark-augmentation}
Soit $\mathcal{C}$ une catégorie possédant les produits fibrés, soit $V$
un objet simplicial et soit $\epsilon : V \to X$ une augmentation.
Soit $U$ l'objet simplicial dont le terme de degré $n$ est le produit
fibré de $(n + 1)$ copies de $V_0$ au-dessus de $X$. En combinant simplement
les lemmes \ref{lemma-augmentation-howto} et
\ref{lemma-cosk-minus-one}, on obtient un morphisme canonique $V \to U$.
\end{remark}
\section{Adjoints à gauche des foncteurs squelette}
\label{section-adjoint-left}

\noindent
Dans cette section, nous construisons, dans certains cas, un adjoint à
gauche $i_{m!}$ du foncteur squelette $\text{sk}_m$. La formule
d'adjonction est
$$
\Mor_{\text{Simp}_m(\mathcal{C})}(U, \text{sk}_mV)
=
\Mor_{\text{Simp}(\mathcal{C})}(i_{m!}U, V).
$$
Cet adjoint à gauche existe lorsque la catégorie $\mathcal{C}$ possède
les colimites finies.

\medskip\noindent
Nous employons une construction analogue à celle de la section
\ref{section-skeleton}. Rappelons la catégorie $[n]/\Delta$ des objets
sous $[n]$, voir Categories, exemple
\ref{categories-example-category-under-X}. Ses objets sont les morphismes
$\alpha : [n] \to [k]$ et ses morphismes les triangles commutatifs.
Notons $([n]/\Delta)_{\leq m}$ la sous-catégorie pleine de
$[n]/\Delta$ formée des objets $[n] \to [k]$ avec $k \leq m$.
Étant donné un objet simplicial $m$-tronqué $U$ de $\mathcal{C}$,
on définit un foncteur
$$
U(n) : ([n]/\Delta)_{\leq m}^{opp} \longrightarrow \mathcal{C}
$$
par les règles
\begin{eqnarray*}
([n] \to [k]) & \longmapsto & U_k \\
\psi : ([n] \to [k']) \to ([n] \to [k])
& \longmapsto &
U(\psi) : U_k \to U_{k'}
\end{eqnarray*}
À tout morphisme $\varphi : [n] \to [n']$ de $\Delta$ est associé un foncteur
$$
\underline{\varphi} :
([n']/\Delta)_{\leq m}
\longrightarrow
([n]/\Delta)_{\leq m}
$$
qui envoie $\alpha : [n'] \to [k]$ sur
$\varphi \circ \alpha : [n] \to [k]$. Le composé
$U(n) \circ \underline{\varphi}$ est égal au foncteur $U(n')$.

\begin{lemma}
\label{lemma-left-adjoint-exists}
Soit $\mathcal{C}$ une catégorie possédant les colimites finies.
Les foncteurs $i_{m!}$ existent pour tout $m$. Soit $U$ un objet
simplicial $m$-tronqué de $\mathcal{C}$. L'objet simplicial $i_{m!}U$
est décrit par la formule
$$
(i_{m!}U)_n = \colim_{([n]/\Delta)_{\leq m}^{opp}} U(n)
$$
et, pour $\varphi : [n] \to [n']$, l'application
$i_{m!}U(\varphi)$ provient de l'identification
$U(n) \circ \underline{\varphi} = U(n')$ ci-dessus, par Categories,
lemme \ref{categories-lemma-functorial-colimit}.
\end{lemma}

\begin{proof}
Dans cette démonstration, notons $i_{m!}U$ l'objet simplicial dont le
terme de degré $n$ est donné par la formule affichée du lemme. Nous allons
montrer qu'il satisfait à la propriété d'adjonction.

\medskip\noindent
Soit $V$ un objet simplicial de $\mathcal{C}$ et soit
$\gamma : U \to \text{sk}_mV$. Un morphisme
$$
\colim_{([n]/\Delta)_{\leq m}^{opp}} U(n) \to T
$$
est donné par un système compatible de morphismes
$f_\alpha : U_k \to T$, où $\alpha : [n] \to [k]$ et $k \leq m$.
On obtient un tel système en prenant les composés
$$
U_k \xrightarrow{\gamma_k} V_k \xrightarrow{V(\alpha)} V_n.
$$
On obtient donc un morphisme induit $(i_{m!}U)_n \to V_n$.
Nous laissons au lecteur le soin de vérifier que ceux-ci forment un
morphisme d'objets simpliciaux $\gamma' : i_{m!}U \to V$.

\medskip\noindent
Réciproquement, étant donné un morphisme $\gamma' : i_{m!}U \to V$,
on obtient un morphisme $\gamma : U \to \text{sk}_m V$ en définissant
$\gamma_i : U_i \to V_i$ comme le composé
$$
U_i
\xrightarrow{\text{id}_{[i]}}
\colim_{([i]/\Delta)_{\leq m}^{opp}} U(i)
\xrightarrow{\gamma'_i}
V_i
$$
pour $0 \leq i \leq n$. Nous laissons au lecteur le soin de vérifier que
cette construction est inverse de la précédente.
\end{proof}

\begin{lemma}
\label{lemma-recovering-U}
Soit $\mathcal{C}$ une catégorie et soit $U$ un objet simplicial
$m$-tronqué de $\mathcal{C}$. Pour tout $n \leq m$, la colimite
$$
\colim_{([n]/\Delta)_{\leq m}^{opp}} U(n)
$$
existe et est égale à $U_n$.
\end{lemma}

\begin{proof}
En effet, la catégorie $([n]/\Delta)_{\leq m}$ possède un objet initial,
à savoir $\text{id} : [n] \to [n]$.
\end{proof}

\begin{lemma}
\label{lemma-recovering-U-for-real}
Soit $\mathcal{C}$ une catégorie possédant les colimites finies et soit
$U$ un objet simplicial $m$-tronqué de $\mathcal{C}$. L'application
$U \to \text{sk}_m i_{m!}U$ est un isomorphisme.
\end{lemma}

\begin{proof}
Combiner les lemmes \ref{lemma-left-adjoint-exists} et
\ref{lemma-recovering-U}.
\end{proof}

\begin{lemma}
\label{lemma-imshriek-sets}
Si $U$ est un ensemble simplicial $m$-tronqué et si $n > m$, alors tous
les $n$-simplexes de $i_{m!}U$ sont dégénérés.
\end{lemma}

\begin{proof}
Cela résulte de la construction de $i_{m!}U$ dans le lemme
\ref{lemma-left-adjoint-exists}, mais on peut aussi raisonner directement.
Posons $V = i_{m!}U$. Soit $V' \subset V$ le sous-ensemble simplicial tel
que $V'_i = V_i$ pour $i \leq m$ et dont tous les $i$-simplexes sont
dégénérés pour $i > m$, voir le lemme
\ref{lemma-simplicial-set-n-skel-sub}. Par la formule d'adjonction, puisque
$\text{sk}_m V' = U$, l'injection $V' \to V$ possède un inverse. Ainsi
$V' = V$.
\end{proof}

\begin{lemma}
\label{lemma-n-skeleton-sets}
Soit $U$ un ensemble simplicial et soit $n \geq 0$ un entier. Le
morphisme $i_{n!} \text{sk}_n U \to U$ identifie
$i_{n!} \text{sk}_n U$ à l'ensemble simplicial $U' \subset U$ défini
dans le lemme \ref{lemma-simplicial-set-n-skel-sub}.
\end{lemma}

\begin{proof}
Par le lemme \ref{lemma-imshriek-sets}, les seuls simplexes non dégénérés
de $i_{n!} \text{sk}_n U$ sont en degrés $\leq n$. L'application
$i_{n!} \text{sk}_n U \to U$ est un isomorphisme en degrés $\leq n$.
On en conclut que l'application $i_{n!} \text{sk}_n U \to U$ envoie les
simplexes non dégénérés sur des simplexes non dégénérés et que deux
simplexes non dégénérés distincts n'ont pas la même image. Le lemme
\ref{lemma-injective-map-simplicial-sets} s'applique donc. Ainsi
$i_{n!} \text{sk}_n U \to U$ est injective, d'où le résultat.
\end{proof}

\begin{remark}
\label{remark-sk-literature}
Dans certains textes, le foncteur composé
$$
\text{Simp}(\mathcal{C})
\xrightarrow{\text{sk}_m}
\text{Simp}_m(\mathcal{C})
\xrightarrow{i_{m!}}
\text{Simp}(\mathcal{C})
$$
est noté $\text{sk}_m$. Cela se justifie pour les ensembles simpliciaux,
car le lemme \ref{lemma-n-skeleton-sets} affirme alors que
$i_{m!} \text{sk}_m V$ est simplement le sous-ensemble simplicial de
$V$ formé de tous les $i$-simplexes de $V$, $i \leq m$, et de leurs
dégénérescences. Dans ces textes, il est également usuel de noter le composé
$$
\text{Simp}(\mathcal{C})
\xrightarrow{\text{sk}_m}
\text{Simp}_m(\mathcal{C})
\xrightarrow{\text{cosk}_m}
\text{Simp}(\mathcal{C})
$$
par $\text{cosk}_m$.
\end{remark}

\begin{lemma}
\label{lemma-glue-simplex}
Soient $U \subset V$ des ensembles simpliciaux. Supposons que
$n \geq 0$ et $x \in V_n$, $x \not \in U_n$, soient tels que
\begin{enumerate}
\item $V_i = U_i$ pour $i < n$,
\item $V_n = U_n \cup \{x\}$,
\item tout $z \in V_j$, $z \not \in U_j$, pour $j > n$, est dégénéré.
\end{enumerate}
Soit $\Delta[n] \to V$ l'unique morphisme qui envoie le $n$-simplexe non
dégénéré de $\Delta[n]$ sur $x$. Alors le diagramme
$$
\xymatrix{
\Delta[n] \ar[r] & V \\
i_{(n - 1)!} \text{sk}_{n - 1} \Delta[n] \ar[r] \ar[u] & U \ar[u]
}
$$
est cocartésien.
\end{lemma}

\begin{proof}
Pour abréger, notons
$\partial \Delta[n] = i_{(n - 1)!} \text{sk}_{n - 1} \Delta[n]$.
Il existe une application naturelle
$U \amalg_{\partial \Delta[n]} \Delta[n] \to V$. Il faut montrer qu'elle
est bijective en degré $j$, pour tout $j$. C'est clair pour $j \leq n$.
Soit $j > n$. La troisième condition signifie que tout $z \in V_j$,
$z \not \in U_j$, est un simplexe dégénéré, disons
$z = s^{j - 1}_i(z')$. Bien entendu, $z' \not \in U_{j - 1}$.
Par récurrence, $z'$ est une dégénérescence de $x$. Ainsi, tout
$j$-simplexe de $V$ appartient à $U$ ou est une dégénérescence de $x$.
L'application $U \amalg_{\partial \Delta[n]} \Delta[n] \to V$ est donc
surjective. Remarquons qu'un simplexe non dégénéré de
$U \amalg_{\partial \Delta[n]} \Delta[n]$ est soit l'image d'un simplexe
non dégénéré de $U$, soit l'image de l'unique $n$-simplexe non dégénéré
de $\Delta[n]$. Comme $x$ est manifestement non dégénéré, l'application
$U \amalg_{\partial \Delta[n]} \Delta[n] \to V$ envoie les simplexes non
dégénérés sur des simplexes non dégénérés et est injective sur les
simplexes non dégénérés. Elle est donc injective par le lemme
\ref{lemma-injective-map-simplicial-sets}.
\end{proof}

\begin{lemma}
\label{lemma-add-simplices}
Soient $U \subset V$ des ensembles simpliciaux tels que $U_n, V_n$
soient finis et non vides pour tout $n$. Supposons que $U$ et $V$
n'aient qu'un nombre fini de simplexes non dégénérés. Alors il existe une
suite de sous-ensembles simpliciaux
$$
U = W^0 \subset W^1 \subset W^2 \subset \ldots W^r = V
$$
telle que le lemme \ref{lemma-glue-simplex} s'applique à chacune des
inclusions $W^i \subset W^{i + 1}$.
\end{lemma}

\begin{proof}
Soit $n$ le plus petit entier tel que $V$ possède un simplexe non dégénéré
qui n'appartient pas à $U$. Soit $x \in V_n$, $x\not \in U_n$, un tel
simplexe non dégénéré. Soit $W \subset V$ l'ensemble des éléments qui
appartiennent à $U$ ou qui sont une dégénérescence, éventuellement
itérée, de $x$ (autrement dit, qui sont de la forme $V(\varphi)(x)$,
avec $\varphi : [m] \to [n]$ surjective). On vérifie aisément que $W$ est
un ensemble simplicial. L'inclusion $U \subset W$ satisfait aux
conditions du lemme \ref{lemma-glue-simplex}. De plus, le nombre de
simplexes non dégénérés de $V$ qui ne sont pas contenus dans $W$ est
exactement inférieur d'une unité au nombre de simplexes non dégénérés de
$V$ qui ne sont pas contenus dans $U$. Le résultat s'ensuit par récurrence
sur ce nombre.
\end{proof}

\begin{lemma}
\label{lemma-imshriek-abelian}
Soit $\mathcal{A}$ une catégorie abélienne. Soit $U$ un objet simplicial
$m$-tronqué de $\mathcal{A}$. Pour $n > m$, on a $N(i_{m!}U)_n = 0$.
\end{lemma}

\begin{proof}
Posons $V = i_{m!}U$. Soit $V' \subset V$ le sous-objet simplicial de
$V$ tel que $V'_i = V_i$ pour $i \leq m$ et $N(V'_i) = 0$ pour
$i > m$, voir le lemme \ref{lemma-simplicial-abelian-n-skel-sub}.
Par la formule d'adjonction, puisque $\text{sk}_m V' = U$, l'injection
$V' \to V$ possède un inverse. Ainsi $V' = V$.
\end{proof}

\begin{lemma}
\label{lemma-n-skeleton-abelian}
Soit $\mathcal{A}$ une catégorie abélienne, soit $U$ un objet simplicial
de $\mathcal{A}$ et soit $n \geq 0$ un entier. Le morphisme
$i_{n!} \text{sk}_n U \to U$ identifie $i_{n!} \text{sk}_n U$ au
sous-objet simplicial $U' \subset U$ défini dans le lemme
\ref{lemma-simplicial-abelian-n-skel-sub}.
\end{lemma}

\begin{proof}
Par le lemme \ref{lemma-imshriek-abelian}, on a
$N(i_{n!} \text{sk}_n U)_i = 0$ pour $i > n$. L'application
$i_{n!} \text{sk}_n U \to U$ est un isomorphisme en degrés $\leq n$,
voir le lemme \ref{lemma-recovering-U-for-real}. On en conclut que
l'application $i_{n!} \text{sk}_n U \to U$ induit des applications
injectives $N(i_{n!} \text{sk}_n U)_i \to N(U)_i$ pour tout $i$. Le
lemme \ref{lemma-injective-map-simplicial-abelian} s'applique donc. Ainsi
$i_{n!} \text{sk}_n U \to U$ est injective, d'où le résultat.
\end{proof}

\noindent
Voici une autre interprétation du foncteur cosquelette à l'aide des
résultats précédents.

\begin{lemma}
\label{lemma-cosk-shriek}
Soit $\mathcal{C}$ une catégorie possédant les sommes finies et les
limites finies. Soit $V$ un objet simplicial de $\mathcal{C}$. Alors
$$
(\text{cosk}_n \text{sk}_n V)_{n + 1}
=
\Hom(i_{n !}\text{sk}_n \Delta[n + 1], V)_0.
$$
\end{lemma}

\begin{proof}
Par le lemme \ref{lemma-morphism-from-coproduct}, l'objet de gauche
représente le foncteur qui associe à $X$ le premier ensemble de la suite
d'égalités suivante :
\begin{eqnarray*}
\Mor(X \times \Delta[n + 1], \text{cosk}_n \text{sk}_n V)
& = &
\Mor(X \times \text{sk}_n \Delta[n + 1], \text{sk}_n V) \\
& = &
\Mor(X \times i_{n !} \text{sk}_n \Delta[n + 1], V).
\end{eqnarray*}
L'objet de droite dans la formule du lemme est représenté par le foncteur
qui associe à $X$ le dernier ensemble de la suite d'égalités. Cela démontre
le résultat.

\medskip\noindent
Dans la suite d'égalités, nous avons utilisé
$\text{sk}_n (X \times \Delta[n + 1]) = X \times \text{sk}_n \Delta[n + 1]$
et
$i_{n!}(X \times \text{sk}_n \Delta[n + 1]) =
X \times i_{n !} \text{sk}_n \Delta[n + 1]$.
La première égalité est évidente. Pour tout objet simplicial $W$,
éventuellement tronqué, de $\mathcal{C}$ et tout objet $X$ de
$\mathcal{C}$, notons provisoirement $\Mor_\mathcal{C}(X, W)$ l'ensemble
simplicial, éventuellement tronqué,
$[n] \mapsto \Mor_\mathcal{C}(X, W_n)$. Il résulte des définitions que
$\Mor(U \times X, W) =
\Mor(U, \Mor_\mathcal{C}(X, W))$ pour tout ensemble simplicial $U$,
éventuellement tronqué. Ainsi
\begin{eqnarray*}
\Mor(X \times i_{n !} \text{sk}_n \Delta[n + 1], W)
& = &
\Mor(i_{n !} \text{sk}_n \Delta[n + 1], \Mor_\mathcal{C}(X, W)) \\
& = &
\Mor(\text{sk}_n \Delta[n + 1],
\text{sk}_n\Mor_\mathcal{C}(X, W)) \\
& = &
\Mor(X \times \text{sk}_n \Delta[n + 1], \text{sk}_nW) \\
& = &
\Mor(i_{n!}(X \times \text{sk}_n \Delta[n + 1]), W).
\end{eqnarray*}
Cela démontre la seconde égalité utilisée et achève la démonstration.
\end{proof}




\section{Objets simpliciaux dans les catégories abéliennes}
\label{section-abelian}

\noindent
Rappelons que la notion de catégorie abélienne est définie dans Homology,
section \ref{homology-section-abelian-categories}.

\begin{lemma}
\label{lemma-abelian}
Soit $\mathcal{A}$ une catégorie abélienne.
\begin{enumerate}
\item Les catégories $\text{Simp}(\mathcal{A})$ et
$\text{CoSimp}(\mathcal{A})$ sont abéliennes.
\item Un morphisme d'objets (co)simpliciaux $f : A \to B$ est injectif
si et seulement si chaque $f_n : A_n \to B_n$ est injectif.
\item Un morphisme d'objets (co)simpliciaux $f : A \to B$ est surjectif
si et seulement si chaque $f_n : A_n \to B_n$ est surjectif.
\item Une suite d'objets (co)simpliciaux
$$
A \xrightarrow{f} B \xrightarrow{g} C
$$
est exacte en $B$ si et seulement si chaque suite
$$
A_i \xrightarrow{f_i} B_i \xrightarrow{g_i} C_i
$$
est exacte en $B_i$.
\end{enumerate}
\end{lemma}

\begin{proof}
La préadditivité est immédiate. Un objet final est donné par
$U_n = 0$ en tout degré. L'existence des produits directs a été établie
dans les lemmes \ref{lemma-product} et
\ref{lemma-product-cosimplicial-objects}. Les noyaux et conoyaux
s'obtiennent degré par degré.
\end{proof}

\noindent
Pour un objet $A$ de $\mathcal{A}$ et un entier $k$, considérons l'objet
simplicial $k$-tronqué $U$ défini par
\begin{enumerate}
\item $U_i = 0$ pour $i < k$,
\item $U_k = A$,
\item tous les morphismes $U(\varphi)$ sont nuls, à l'exception de
$U(\text{id}_{[k]}) = \text{id}_A$.
\end{enumerate}
Comme $\mathcal{A}$ possède les limites et les colimites finies,
$\text{cosk}_k U$ et $i_{k!}U$ existent tous deux. Nous allons les
décrire, ainsi que l'application canonique
$i_{k!}U \to \text{cosk}_kU$.

\begin{lemma}
\label{lemma-eilenberg-maclane-object}
Avec $A$, $k$ et $U$ comme ci-dessus, de sorte que $U_i = 0$, $i < k$,
et $U_k = A$.
\begin{enumerate}
\item Étant donné un objet simplicial $k$-tronqué $V$, on a
$$
\Mor(U, V)
=
\{ f : A \to V_k \mid d^k_i \circ f = 0, \ i = 0, \ldots, k \}
$$
et
$$
\Mor(V, U)
=
\{ f : V_k \to A \mid f \circ s^{k - 1}_i = 0, \ i = 0, \ldots, k - 1 \}.
$$
\item L'objet $i_{k!} U$ a pour terme de degré $n$
$\bigoplus_\alpha A$, où $\alpha$ parcourt tous les morphismes surjectifs
$\alpha : [n] \to [k]$.
\item Pour tout $\varphi : [m] \to [n]$, l'application
$i_{k!} U(\varphi)$ est décrite par l'application
$\bigoplus_\alpha A \to \bigoplus_{\alpha'} A$
qui envoie la composante correspondant à $\alpha : [n] \to [k]$ sur
zéro si $\alpha \circ \varphi$ n'est pas surjective, et par l'identité
sur la composante correspondant à $\alpha \circ \varphi$ si elle est
surjective.
\item L'objet $\text{cosk}_k U$ a pour terme de degré $n$
$\bigoplus_\beta A$, où $\beta$ parcourt tous les morphismes injectifs
$\beta : [k] \to [n]$.
\item Pour tout $\varphi : [m] \to [n]$, l'application
$\text{cosk}_k U(\varphi)$ est décrite par l'application
$\bigoplus_\beta A \to \bigoplus_{\beta'} A$
qui envoie la composante correspondant à $\beta : [k] \to [n]$ sur zéro
si $\beta$ ne se factorise pas par $\varphi$, et par l'identité sur
chacune des composantes correspondant à $\beta'$ telles que
$\beta = \varphi \circ \beta'$ dans le cas contraire.
\item L'application canonique
$
c : i_{k !} U \to \text{cosk}_k U
$
possède en degré $n$ un coefficient $(\alpha, \beta)$, $A \to A$, nul
si $\alpha \circ \beta$ n'est pas l'identité et égal à $\text{id}_A$
sinon.
\item L'application canonique
$
c : i_{k !} U \to \text{cosk}_k U
$
est injective.
\end{enumerate}
\end{lemma}

\begin{proof}
La démonstration de (1) est laissée au lecteur.

\medskip\noindent
Prenons les règles de (2) et (3) comme définition d'un objet simplicial,
noté $\tilde U$. Nous allons montrer qu'il constitue une réalisation de
$i_{k!}U$, ce qui démontrera simultanément (2) et (3). Il faut montrer
qu'étant donné un morphisme $f : U \to \text{sk}_kV$, il existe un unique
morphisme $\tilde f : \tilde U \to V$ qui redonne $f$ après passage au
$k$-squelette.
D'après (1), $f$ correspond à un morphisme $f_k : A \to V_k$ dont l'image
est contenue dans le noyau de $d^k_i$ pour tout $i$. Pour toute surjection
$\alpha : [n] \to [k]$, définissons $\tilde f_\alpha : A \to V_n$ comme
le composé $\tilde f_\alpha = V(\alpha) \circ f_k : A \to V_n$.
Définissons $\tilde f_n : \tilde U_n \to V_n$ comme la somme des
$\tilde f_\alpha$ pour les surjections $\alpha : [n] \to [k]$.
Une telle famille de $\tilde f_\alpha$ définit un morphisme d'objets
simpliciaux si et seulement si, pour tout $\varphi : [m] \to [n]$, le diagramme
$$
\xymatrix{
\bigoplus_{\alpha : [n] \to [k]\text{ surjective}} A
\ar[r]_-{\tilde f_n}
\ar[d]_{(3)} &
V_n \ar[d]^{V(\varphi)} \\
\bigoplus_{\alpha' : [m] \to [k]\text{ surjective}} A
\ar[r]^-{\tilde f_m} &
V_m
}
$$
soit commutatif. Le choix $\varphi = \alpha$ montre que
$\tilde f_\alpha$ est uniquement déterminé par $f_k$. La commutativité se
vérifie en général séparément sur chaque facteur direct du coin supérieur
gauche. Elle est claire pour les facteurs correspondant aux $\alpha$ tels
que $\alpha \circ \varphi$ soit surjective, car ils sont envoyés par
$\text{id}_A$ sur le facteur avec $\alpha' = \alpha \circ \varphi$, et l'on a
$\tilde f_{\alpha'} = V(\alpha') \circ f_k =
V(\alpha \circ \varphi) \circ f_k = V(\varphi) \circ \tilde f_\alpha$.
Pour ceux où $\alpha \circ \varphi$ n'est pas surjective, il faut montrer
que $V(\varphi) \circ \tilde f_\alpha = 0$. Par définition, ce morphisme
est égal à
$V(\varphi) \circ V(\alpha) \circ f_k = V(\alpha \circ \varphi) \circ f_k$.
Comme $\alpha \circ \varphi$ n'est pas surjective, on peut l'écrire sous
la forme $\delta^k_i \circ \psi$, et l'on en déduit que
$V(\varphi) \circ V(\alpha) \circ f_k =
V(\psi) \circ d^k_i \circ f_k = 0$, comme ci-dessus.

\medskip\noindent
Prenons les règles de (4) et (5) comme définition d'un objet simplicial,
noté $\tilde U$. Nous allons montrer qu'il constitue une réalisation de
$\text{cosk}_k U$, ce qui démontrera simultanément (4) et (5). L'argument
est entièrement dual de la démonstration de (2) et (3), mais nous le donnons.
Il faut montrer qu'étant donné un morphisme $f : \text{sk}_kV \to U$, il
existe un unique morphisme $\tilde f : V \to \tilde U$ qui redonne $f$
après passage au $k$-squelette. D'après (1), $f$ correspond à un morphisme
$f_k : V_k \to A$ nul sur l'image de $s^{k - 1}_i$ pour tout $i$. Pour
toute injection $\beta : [k] \to [n]$, définissons
$\tilde f_\beta : V_n \to A$ comme le composé
$\tilde f_\beta = f_k \circ V(\beta) : V_n \to A$. Définissons
$\tilde f_n : V_n \to \tilde U_n$ comme la somme des $\tilde f_\beta$
pour les injections $\beta : [k] \to [n]$. Une telle famille de
$\tilde f_\beta$ définit un morphisme d'objets simpliciaux si et seulement
si, pour tout $\varphi : [m] \to [n]$, le diagramme
$$
\xymatrix{
V_n
\ar[d]_{V(\varphi)}
\ar[r]_-{\tilde f_n}
&
\bigoplus_{\beta : [k] \to [n]\text{ injective}} A
\ar[d]^{(5)}
\\
V_m
\ar[r]^-{\tilde f_m}
&
\bigoplus_{\beta' : [k] \to [m]\text{ injective}} A
}
$$
soit commutatif. Le choix $\varphi = \beta$ montre que
$\tilde f_\beta$ est uniquement déterminé par $f_k$. La commutativité se
vérifie en général séparément sur chaque facteur direct du coin inférieur
droit. Elle est claire pour les facteurs correspondant aux $\beta'$ tels
que $\varphi \circ \beta'$ soit injective, car l'unique facteur envoyé
dans chacun d'eux est celui avec $\beta = \varphi \circ \beta'$, et l'on
a alors
$\tilde f_{\beta'} \circ V(\varphi) =
f_k \circ V(\beta') \circ V(\varphi) =
f_k \circ V(\beta) = \tilde f_\beta$. Pour ceux où
$\varphi \circ \beta'$ n'est pas injective, il faut montrer que
$\tilde f_{\beta'} \circ V(\varphi) = 0$. Par définition, ce morphisme est égal à
$f_k \circ V(\beta') \circ V(\varphi) =
f_k \circ V(\varphi \circ \beta')$.
Comme $\varphi \circ \beta'$ n'est pas injective, on peut l'écrire sous
la forme $\psi \circ \sigma^{k - 1}_i$, et l'on en déduit que
$f_k \circ V(\beta') \circ V(\varphi) =
f_k \circ s^{k - 1}_i \circ V(\psi) = 0$, comme ci-dessus.

\medskip\noindent
Le composé $i_{k!}U \to \text{cosk}_kU$ est l'unique application
d'objets simpliciaux qui soit l'identité sur
$A = U_k = (i_{k!}U)_k = (\text{cosk}_kU)_k$. Il suffit donc de vérifier
que la règle proposée définit un morphisme d'objets simpliciaux. Pour cela,
il faut montrer que, pour tout $\varphi : [m] \to [n]$, le diagramme
$$
\xymatrix{
\bigoplus_{\alpha : [n] \to [k]\text{ surjective}} A
\ar[d]_{(3)}
\ar[r]_{(6)}
&
\bigoplus_{\beta : [k] \to [n]\text{ injective}} A
\ar[d]^{(5)}
\\
\bigoplus_{\alpha' : [m] \to [k]\text{ surjective}} A
\ar[r]^{(6)}
&
\bigoplus_{\beta' : [k] \to [m]\text{ injective}} A
}
$$
soit commutatif. On peut l'interpréter en termes de matrices ne contenant
que des $0$ et des $1$, comme suit. La matrice de (3) possède une entrée
$(\alpha', \alpha)$ non nulle si et seulement si
$\alpha' = \alpha \circ \varphi$. De même, la matrice de (5) possède une
entrée $(\beta', \beta)$ non nulle si et seulement si
$\beta = \varphi \circ \beta'$. La matrice supérieure de (6) possède une
entrée $(\alpha, \beta)$ non nulle si et seulement si
$\alpha \circ \beta = \text{id}_{[k]}$, et de même pour la matrice
inférieure de (6). La commutativité du diagramme revient alors à calculer
l'entrée $(\alpha, \beta')$ des deux composés et à constater qu'elles sont
égales. Cela revient à l'égalité suivante :
$$
\# \left\{
\beta
\mid
\beta = \varphi \circ \beta'
\text{ et }
\alpha \circ \beta = \text{id}_{[k]}
\right\}
=
\# \left\{
\alpha'
\mid
\alpha' = \alpha \circ \varphi
\text{ et }
\alpha' \circ \beta' = \text{id}_{[k]}
\right\}
$$
dont la démonstration peut sans risque être laissée au lecteur.

\medskip\noindent
Enfin, démontrons (7). Cela résulte directement des lemmes
\ref{lemma-injective-map-simplicial-abelian},
\ref{lemma-recover-cosk}, \ref{lemma-recovering-U-for-real}
et \ref{lemma-imshriek-abelian}.
\end{proof}

\begin{definition}
\label{definition-eilenberg-maclane}
Soit $\mathcal{A}$ une catégorie abélienne. Soit $A$ un objet de
$\mathcal{A}$ et soit $k$ un entier $\geq 0$. L'{\it objet
d'Eilenberg--Mac Lane $K(A, k)$} est l'objet $K(A, k) = i_{k!}U$ décrit
dans le lemme \ref{lemma-eilenberg-maclane-object}.
\end{definition}


\begin{lemma}
\label{lemma-extension}
Soit $\mathcal{A}$ une catégorie abélienne. Soit $A$ un objet de
$\mathcal{A}$ et soit $k$ un entier $\geq 0$. Considérons l'objet
simplicial $E$ défini par les règles suivantes :
\begin{enumerate}
\item $E_n = \bigoplus_\alpha A$, où la somme porte sur les
$\alpha : [n] \to [k + 1]$ dont l'image est $[k]$ ou $[k + 1]$.
\item Étant donné $\varphi : [m] \to [n]$, l'application
$E_n \to E_m$ envoie le facteur correspondant à $\alpha$, par
$\text{id}_A$, sur le facteur correspondant à $\alpha \circ \varphi$,
pourvu que $\Im(\alpha \circ \varphi)$ soit égal à $[k]$ ou $[k + 1]$.
\end{enumerate}
Alors il existe une suite exacte courte
$$
0 \to K(A, k) \to E \to K(A, k + 1) \to 0
$$
qui est scindée exacte degré par degré.
\end{lemma}

\begin{proof}
Les applications $K(A, k)_n \to E_n$, resp.\
$E_n \to K(A, k + 1)_n$, sont données par l'inclusion des sommes
directes, resp.\ la projection des sommes directes, qui résultent
manifestement des inclusions des ensembles d'indices. Ce sont clairement
des applications d'objets simpliciaux.
\end{proof}

\begin{lemma}
\label{lemma-abelian-limit-skeleta}
Soit $\mathcal{A}$ une catégorie abélienne. Pour tout objet simplicial
$V$ de $\mathcal{A}$, on a
$$
V = \colim_n i_{n!}\text{sk}_n V
$$
où toutes les applications de transition sont injectives.
\end{lemma}

\begin{proof}
En effet, chaque $V_m$ est égal à $(i_{n!}\text{sk}_n V)_m$ dès que
$n \geq m$. Voir aussi le lemme \ref{lemma-n-skeleton-abelian} pour les
applications de transition.
\end{proof}

\section{Objets simpliciaux et complexes de chaînes}
\label{section-complexes}

\noindent
Soit $\mathcal{A}$ une catégorie abélienne. Voir Homology, section
\ref{homology-section-complexes}, pour les conventions et les notations
relatives aux complexes de chaînes. Soit $U$ un objet simplicial de
$\mathcal{A}$. Le {\it complexe de chaînes associé $s(U)$ à $U$},
parfois appelé {\it complexe de Moore}, est le complexe de chaînes
$$
\ldots \to U_2 \to U_1 \to U_0 \to 0 \to 0 \to \ldots
$$
dont les différentielles $d_n : U_n \to U_{n - 1}$ sont données par la formule
$$
d_n = \sum\nolimits_{i = 0}^n (-1)^i d^n_i.
$$
C'est bien un complexe car, d'après les relations de la remarque
\ref{remark-relations}, on a
\begin{eqnarray*}
d_n \circ d_{n + 1}
& = &
(\sum\nolimits_{i = 0}^n (-1)^i d^n_i) \circ
(\sum\nolimits_{j = 0}^{n + 1} (-1)^j d^{n + 1}_j) \\
& = &
\sum\nolimits_{0 \leq i < j \leq n + 1}
(-1)^{i + j} d^n_{j - 1} \circ d^{n + 1}_i
+
\sum\nolimits_{n \geq i \geq j \geq 0}
(-1)^{i + j} d^n_i \circ d^{n + 1}_j \\
& = & 0.
\end{eqnarray*}
Les signes s'annulent ! On note $s(U)$ le complexe de chaînes associé.
La construction est manifestement fonctorielle et définit donc un foncteur
$$
s : \text{Simp}(\mathcal{A}) \longrightarrow \text{Ch}_{\geq 0}(\mathcal{A}).
$$
On a ainsi la formule déconcertante mais correcte $s(U)_n = U_n$.

\begin{lemma}
\label{lemma-s-exact}
Le foncteur $s$ est exact.
\end{lemma}

\begin{proof}
Cela résulte du lemme \ref{lemma-abelian}.
\end{proof}

\begin{lemma}
\label{lemma-homology-extension}
Soit $\mathcal{A}$ une catégorie abélienne. Soit $A$ un objet de
$\mathcal{A}$, soit $k$ un entier et soit $E$ l'objet décrit dans le
lemme \ref{lemma-extension}. Alors le complexe $s(E)$ est acyclique.
\end{lemma}

\begin{proof}
Pour un morphisme $\alpha : [n] \to [k + 1]$, définissons
$\alpha' : [n + 1] \to [k + 1]$ comme l'application telle que
$\alpha'|_{[n]} = \alpha$ et $\alpha'(n + 1) = k + 1$. Remarquons que,
si l'image de $\alpha$ est $[k]$ ou $[k + 1]$, l'image de $\alpha'$ est
$[k + 1]$. Considérons la famille d'applications
$h_n : E_n \to E_{n + 1}$ qui envoie le facteur correspondant à $\alpha$
sur celui correspondant à $\alpha'$ par l'identité de $A$. Calculons
$d_{n + 1} \circ h_n - h_{n - 1} \circ d_n$. Faisons-le d'abord lorsque
la catégorie $\mathcal{A}$ est celle des groupes abéliens. Notons
$x_\alpha$ l'élément $x \in A$ dans le facteur de $E_n$ correspondant à
l'application $\alpha$ de l'ensemble d'indices. Convenons aussi que
$x_\alpha$ désigne l'élément nul de $E_n$ lorsque $\Im(\alpha)$ n'est ni
$[k]$ ni $[k + 1]$. Avec ces conventions,
$$
d_{n + 1}(h_n(x_\alpha)) =
\sum\nolimits_{i = 0}^{n + 1} (-1)^i x_{\alpha' \circ \delta^{n + 1}_i}
$$
et
$$
h_{n - 1}(d_n(x_\alpha)) =
\sum\nolimits_{i = 0}^n (-1)^i x_{(\alpha \circ \delta_i^n)'}
$$
On vérifie aisément que
$\alpha' \circ \delta^{n + 1}_i = (\alpha \circ \delta_i^n)'$
pour $i = 0, \ldots, n$. On vérifie également que
$\alpha' \circ \delta^{n + 1}_{n + 1} = \alpha$. Ainsi,
$$
(d_{n + 1} \circ h_n - h_{n - 1} \circ d_n)(x_\alpha)
=
(-1)^{n + 1} x_\alpha
$$
Ces identités restent valables pour toute catégorie abélienne
$\mathcal{A}$, puisqu'elles valent dans le groupe abélien simplicial
$[n] \mapsto \Hom(A, E_n)$ ; les détails sont laissés au lecteur.
On en conclut que l'identité de $E$ est homotope à zéro, l'homotopie étant
donnée par la famille d'applications
$h'_n = (-1)^{n + 1}h_n : E_n \to E_{n + 1}$. Ainsi $E$ est acyclique,
par exemple d'après Homology, lemme
\ref{homology-lemma-map-homology-homotopy}.
\end{proof}

\begin{lemma}
\label{lemma-homology-eilenberg-maclane}
Soit $\mathcal{A}$ une catégorie abélienne. Soit $A$ un objet de
$\mathcal{A}$ et soit $k$ un entier. On a
$H_i(s(K(A, k))) = A$ si $i = k$, et $0$ sinon.
\end{lemma}

\begin{proof}
Démontrons d'abord l'énoncé pour $k = 0$. Dans ce cas,
$K(A, 0)_n = A$ pour tout $n$. Toutes les applications de ce groupe
abélien simplicial sont en outre $\text{id}_A$ ; autrement dit,
$K(A, 0)$ est l'objet simplicial constant de valeur $A$. Les
différentielles $d_n = \sum_{i = 0}^n (-1)^i \text{id}_A = 0$ si $n$
est impair et $ = \text{id}_A$ si $n$ est pair. Ainsi $s(K(A, 0))$ est
de la forme
$$
\ldots \to A \xrightarrow{0} A \xrightarrow{1} A \xrightarrow{0} A \to 0
$$
et le résultat est clair.

\medskip\noindent
Montrons ensuite le résultat pour tout $k$ par récurrence. Le résultat
étant acquis pour $k$, considérons la suite exacte courte
$$
0 \to K(A, k) \to E \to K(A, k + 1) \to 0
$$
du lemme \ref{lemma-extension}. Par le lemme \ref{lemma-abelian}, la
suite associée de complexes de chaînes est exacte. Par le lemme
\ref{lemma-homology-extension}, $s(E)$ est acyclique. Le résultat pour
$k + 1$ découle donc de la suite exacte longue d'homologie ; voir
Homology, lemme \ref{homology-lemma-long-exact-sequence-chain}.
\end{proof}

\noindent
On peut associer un second complexe de chaînes à un objet simplicial de
$\mathcal{A}$. Rappelons que, par le lemme
\ref{lemma-splitting-abelian-category}, tout objet simplicial $U$ de
$\mathcal{A}$ est canoniquement scindé avec
$N(U_m) = \bigcap_{i = 0}^{m - 1} \Ker(d^m_i)$.
On définit le {\it complexe de chaînes normalisé $N(U)$} comme le
complexe de chaînes
$$
\ldots \to N(U_2) \to N(U_1) \to N(U_0) \to 0 \to 0 \to \ldots
$$
dont la différentielle $d_n : N(U_n) \to N(U_{n - 1})$ est la restriction
de $(-1)^nd^n_n$ au facteur direct $N(U_n)$ de $U_n$. Remarquons que le
lemme \ref{lemma-N-d-in-N} implique
$d^n_n(N(U_n)) \subset N(U_{n - 1})$. C'est un complexe car
$d^n_n \circ d^{n + 1}_{n + 1} = d^n_n \circ d^{n + 1}_n$
et $d^{n + 1}_n$ est nulle sur $N(U_{n + 1})$ par définition. On obtient
ainsi un second foncteur
$$
N : \text{Simp}(\mathcal{A}) \longrightarrow \text{Ch}_{\geq 0}(\mathcal{A}).
$$
Voici la raison du signe dans la différentielle.

\begin{lemma}
\label{lemma-map-associated-complexes}
Soit $\mathcal{A}$ une catégorie abélienne et soit $U$ un objet
simplicial de $\mathcal{A}$. L'application canonique $N(U_n) \to U_n$
induit un morphisme de complexes $N(U) \to s(U)$.
\end{lemma}

\begin{proof}
C'est clair, car la différentielle sur $s(U)_n = U_n$ est
$\sum (-1)^i d^n_i$, et les applications $d^n_i$, $i < n$, sont nulles
sur $N(U_n)$, tandis que la restriction de $(-1)^nd^n_n$ est, par
définition, la différentielle de $N(U)$.
\end{proof}

\begin{lemma}
\label{lemma-N-K}
Soit $\mathcal{A}$ une catégorie abélienne. Soit $A$ un objet de
$\mathcal{A}$ et soit $k$ un entier. On a
$N(K(A, k))_i = A$ si $i = k$, et $0$ sinon.
\end{lemma}

\begin{proof}
Il est clair que $N(K(A, k))_i = 0$ lorsque $i < k$, car
$K(A, k)_i = 0$ dans ce cas. Il est également clair que
$N(K(A, k))_k = A$, puisque $K(A, k)_{k - 1} = 0$ et
$K(A, k)_k = A$. Pour $i > k$, on a $N(K(A, k))_i = 0$ par le lemme
\ref{lemma-imshriek-abelian} et la définition de $K(A, k)$ ; voir la
définition \ref{definition-eilenberg-maclane}.
\end{proof}

\begin{lemma}
\label{lemma-decompose-associated-complexes}
Soit $\mathcal{A}$ une catégorie abélienne et soit $U$ un objet
simplicial de $\mathcal{A}$. Le morphisme canonique de complexes de
chaînes $N(U) \to s(U)$ est scindé. En fait,
$$
s(U) = N(U) \oplus D(U)
$$
pour un certain complexe $D(U)$. La construction $U \mapsto D(U)$ est
fonctorielle.
\end{lemma}

\begin{proof}
Définissons $D(U)_n$ comme l'image de
$$
\bigoplus\nolimits_{\varphi : [n] \to [m]\text{ surjective}, \ m < n} N(U_m)
\xrightarrow{\bigoplus U(\varphi)}
U_n
$$
qui est un sous-objet de $U_n$ supplémentaire de $N(U_n)$ d'après le
lemme \ref{lemma-splitting-abelian-category} et la définition
\ref{definition-split}. Montrons que $D(U)$ est un sous-complexe.
Fixons une application surjective $\varphi : [n] \to [m]$ avec $m < n$
et considérons le composé
$$
N(U_m) \xrightarrow{U(\varphi)} U_n \xrightarrow{d_n} U_{n - 1}.
$$
Ce composé est la somme des applications
$$
N(U_m) \xrightarrow{U(\varphi \circ \delta^n_i)} U_{n - 1}
$$
affectées du signe $(-1)^i$, $i = 0, \ldots, n$.

\medskip\noindent
Montrons d'abord par récurrence croissante sur $m$, pour
$0 \leq m < n - 1$, que toutes les applications
$U(\varphi \circ \delta^n_i)$ envoient $N(U_m)$ dans $D(U)_{n - 1}$.
(Le cas $m = n - 1$ est traité ci-dessous.) Lorsque l'application
$\varphi \circ \delta^n_i : [n - 1] \to [m]$ est surjective, l'image de
$N(U_m)$ par $U(\varphi \circ \delta^n_i)$ est contenue par définition
dans $D(U)_{n - 1}$. Si
$\varphi \circ \delta^n_i : [n - 1] \to [m]$ n'est pas surjective,
posons $j = \varphi(i)$ et remarquons que $i$ est l'unique indice dont
l'image par $\varphi$ est $j$. On peut écrire
$\varphi \circ \delta^n_i = \delta^m_j \circ \psi$
pour un certain $\psi : [n - 1] \to [m - 1]$. Ainsi,
$U(\varphi \circ \delta^n_i) = U(\psi) \circ d^m_j$,
qui est nulle sur $N(U_m)$ sauf si $j = m$. Si $j = m$, alors
$d^m_m(N(U_m)) \subset N(U_{m - 1})$, et donc
$U(\varphi \circ \delta^n_i)(N(U_m)) \subset
U(\psi \circ \delta^n_i)(N(U_{m - 1}))$ ; on conclut par l'hypothèse de
récurrence.

\medskip\noindent
Pour achever la preuve que $D(U)$ est un sous-complexe, il reste à traiter
le composé
$$
N(U_m) \xrightarrow{U(\varphi)} U_n \xrightarrow{d_n} U_{n - 1}.
$$
dans le cas $m =  n - 1$. Alors $\varphi = \sigma^{n - 1}_j$ pour un
certain $0 \leq j \leq n - 1$, et $U(\varphi) = s^{n - 1}_j$. Le composé
est donc donné par la somme
$$
\sum (-1)^i d^n_i \circ s^{n - 1}_j
$$
Rappelons, d'après la remarque \ref{remark-relations}, que
$d^n_j \circ s^{n - 1}_j = d^n_{j + 1} \circ s^{n - 1}_j = \text{id}$
et que ces deux termes s'annulent, car ils ont des signes opposés.
L'application $d^n_n \circ s^{n - 1}_j$, si $j < n - 1$, est égale à
$s^{n - 2}_j \circ d^{n - 1}_{n - 1}$. Comme
$d^{n - 1}_{n - 1}$ envoie $N(U_{n - 1})$ dans $N(U_{n - 2})$, l'image
$d^n_n ( s^{n - 1}_j (N(U_{n - 1})))$ est contenue dans
$s^{n - 2}_j(N(U_{n - 2}))$, qui est contenu par définition dans
$D(U_{n - 1})$. Pour toutes les autres combinaisons $(i, j)$, on a soit
$d^n_i \circ s^{n - 1}_j = s^{n - 2}_{j - 1} \circ d^{n - 1}_i$
(si $i < j$), soit
$d^n_i \circ s^{n - 1}_j = s^{n - 2}_j \circ d^{n - 1}_{i - 1}$
(si $n > i > j + 1$), et dans ces cas l'application est nulle par la
définition de $N(U_{n - 1})$.
\end{proof}

\begin{remark}
\label{remark-degenerate-subcomplex}
Dans la situation du lemme \ref{lemma-decompose-associated-complexes}, le
sous-complexe $D(U) \subset s(U)$ peut aussi être défini comme le
sous-complexe de termes
$$
D(U)_n = \Im\left(
\bigoplus\nolimits_{\varphi : [n] \to [m]\text{ surjective}, \ m < n} U_m
\xrightarrow{\bigoplus U(\varphi)}
U_n\right)
$$
En effet, puisque $U_m$ est la somme directe du sous-objet $N(U_m)$ et
des images des $N(U_k)$ pour les surjections $[m] \to [k]$ avec $k < m$,
cette définition coïncide manifestement avec celle de $D(U)_n$ donnée
dans la démonstration du lemme \ref{lemma-decompose-associated-complexes}.
Ainsi, si $U$ est un groupe abélien simplicial, les éléments de $D(U)_n$
sont exactement les sommes de $n$-simplexes dégénérés.
\end{remark}

\begin{lemma}
\label{lemma-N-exact}
Le foncteur $N$ est exact.
\end{lemma}

\begin{proof}
Cela résulte du lemme \ref{lemma-s-exact} et de la décomposition
fonctorielle du lemme \ref{lemma-decompose-associated-complexes}.
\end{proof}

\begin{lemma}
\label{lemma-quasi-isomorphism}
Soit $\mathcal{A}$ une catégorie abélienne et soit $V$ un objet
simplicial de $\mathcal{A}$. Le morphisme canonique de complexes de
chaînes $N(V) \to s(V)$ est un quasi-isomorphisme. Autrement dit, le
complexe $D(V)$ du lemme \ref{lemma-decompose-associated-complexes} est
acyclique.
\end{lemma}

\begin{proof}
Remarquons que le résultat vaut pour $K(A, k)$, pour tout objet $A$ et
tout $k \geq 0$, d'après les lemmes
\ref{lemma-homology-eilenberg-maclane} et \ref{lemma-N-K}. Considérons
l'hypothèse $IH_{n, m}$ : pour tout $V$ tel que $V_j = 0$ pour
$j \leq m$, et tout $i \leq n$, l'application $N(V) \to s(V)$ induit un
isomorphisme $H_i(N(V)) \to H_i(s(V))$.

\medskip\noindent
Pour amorcer la récurrence, remarquons que $IH_{n, n}$ est trivialement
vraie, car alors $N(V)_n = 0$ et $s(V)_n = 0$.

\medskip\noindent
Supposons $IH_{n, m}$, avec $m \leq n$. Soit $V$ un objet simplicial tel
que $V_j = 0$ pour $j < m$. Par le lemme
\ref{lemma-eilenberg-maclane-object} et la définition
\ref{definition-eilenberg-maclane}, on a
$K(V_m, m) = i_{m!} \text{sk}_mV$. Par le lemme
\ref{lemma-n-skeleton-abelian}, le morphisme naturel
$$
K(V_m, m) = i_{m!} \text{sk}_mV \to V
$$
est injectif. On obtient donc une suite exacte courte
$$
0 \to K(V_m, m) \to V \to W \to 0
$$
pour un certain $W$ tel que $W_i = 0$ pour $i = 0, \ldots, m$. Cette
suite exacte courte induit un morphisme de suites exactes courtes de
complexes associés
$$
\xymatrix{
0 \ar[r] &
N(K(V_m, m)) \ar[r] \ar[d] &
N(V) \ar[r] \ar[d] &
N(W) \ar[r] \ar[d] &
0 \\
0 \ar[r] &
s(K(V_m, m)) \ar[r] &
s(V) \ar[r] &
s(W) \ar[r] &
0
}
$$
voir les lemmes \ref{lemma-s-exact} et \ref{lemma-N-exact}. Le résultat
pour $V$ se déduit donc du résultat aux extrémités.
\end{proof}






\section{Dold-Kan}
\label{section-dold-kan}

\noindent
Dans cette section, nous démontrons le théorème de Dold--Kan, qui relie
les objets simpliciaux d'une catégorie abélienne aux complexes de chaînes.

\begin{lemma}
\label{lemma-N-faithful}
Soit $\mathcal{A}$ une catégorie abélienne.
Le foncteur $N$ est fidèle et reflète les
isomorphismes, les injections et les surjections.
\end{lemma}

\begin{proof}
La fidélité résulte immédiatement de la
décomposition canonique du lemme \ref{lemma-splitting-abelian-category}.
L'assertion concernant la réflexion des injections, des surjections et des
isomorphismes résulte du
lemme \ref{lemma-injective-map-simplicial-abelian}.
\end{proof}

\begin{lemma}
\label{lemma-S-N}
Soient $\mathcal{A}$ et $\mathcal{B}$ des catégories abéliennes.
Soient $N : \mathcal{A} \to \mathcal{B}$ et
$S : \mathcal{B} \to \mathcal{A}$ des foncteurs.
Supposons que
\begin{enumerate}
\item les foncteurs $S$ et $N$ soient exacts,
\item il existe un isomorphisme $g : N \circ S \to \text{id}_\mathcal{B}$
vers le foncteur identité de $\mathcal{B}$,
\item $N$ soit fidèle, et
\item $S$ soit essentiellement surjectif.
\end{enumerate}
Alors $S$ et $N$ sont des équivalences de catégories quasi inverses.
\end{lemma}

\begin{proof}
Il suffit de construire un isomorphisme fonctoriel
$S(N(A)) \cong A$. Pour cela, choisissons $B$ et un
isomorphisme $f : A  \to S(B)$.
Considérons l'application
$$
f^{-1} \circ g_{S(B)} \circ S(N(f)) :
S(N(A)) \to S(N(S(B))) \to S(B) \to A.
$$
On vérifie aisément qu'elle ne dépend pas du
choix de $f, B$ et qu'elle donne l'isomorphisme cherché
$S \circ N \to \text{id}_\mathcal{A}$.
\end{proof}

\begin{theorem}
\label{theorem-dold-kan}
Soit $\mathcal{A}$ une catégorie abélienne.
Le foncteur $N$ induit une équivalence de
catégories
$$
N :
\text{Simp}(\mathcal{A})
\longrightarrow
\text{Ch}_{\geq 0}(\mathcal{A})
$$
\end{theorem}

\begin{proof}
Nous allons décrire un foncteur en sens inverse,
inspiré par la construction du lemme \ref{lemma-extension}
(à ceci près que nous introduisons un signe pour obtenir les bonnes
différentielles). Soit $A_\bullet$ un complexe de chaînes, de différentielles
$d_{A, n} : A_n \to A_{n - 1}$. Pour chaque $n \geq 0$, posons
$$
I_n =
\Big\{
\alpha : [n] \to \{0, 1, 2, \ldots\}
\mid
\Im(\alpha) = [k]\text{ pour un certain }k
\Big\}.
$$
Pour $\alpha \in I_n$, nous notons $k(\alpha)$ l'unique
entier tel que $\Im(\alpha) = [k]$.
Nous définissons un objet simplicial $S(A_\bullet)$ comme suit :
\begin{enumerate}
\item $S(A_\bullet)_n = \bigoplus_{\alpha \in I_n} A_{k(\alpha)}$, que
nous écrirons
$\bigoplus_{\alpha \in I_n} A_{k(\alpha)} \cdot \alpha$
afin d'évoquer « $\alpha$ » comme vecteur de base du
facteur direct correspondant,
\item étant donné $\varphi : [m] \to [n]$, nous définissons
$S(A_\bullet)(\varphi)$ par sa restriction au
facteur direct $A_{k(\alpha)} \cdot \alpha$
de $S(A_\bullet)_n$, de la manière suivante :
\begin{enumerate}
\item si $\alpha \circ \varphi \not \in I_m$, nous la prenons égale à zéro,
\item si $\alpha \circ \varphi \in I_m$, mais si $k(\alpha \circ \varphi)$
n'est égal ni à $k(\alpha)$ ni à $k(\alpha) - 1$, nous la prenons
encore égale à zéro,
\item si $\alpha \circ \varphi \in I_m$
et $k(\alpha \circ \varphi) = k(\alpha)$, nous utilisons
l'application identité vers le facteur direct
$A_{k(\alpha \circ \varphi)} \cdot (\alpha \circ \varphi)$
de $S(A_\bullet)_m$, et
\item si $\alpha \circ \varphi \in I_m$
et $k(\alpha \circ \varphi) = k(\alpha) - 1$,
nous utilisons $(-1)^{k(\alpha)} d_{A, k(\alpha)}$ vers le facteur direct
$A_{k(\alpha \circ \varphi)}\cdot (\alpha \circ \varphi)$
de $S(A_\bullet)_m$.
\end{enumerate}
\end{enumerate}
Montrons que $S(A_\bullet)$ est un objet simplicial de $\mathcal{A}$.
Pour cela, supposons données des applications $\varphi : [m] \to [n]$ et
$\psi : [n] \to [p]$. Nous allons montrer que
$S(A_\bullet)(\varphi) \circ S(A_\bullet)(\psi) =
S(A_\bullet)(\psi \circ \varphi)$. Choose $\beta \in I_p$ and set
$\alpha = \beta \circ \psi$ and $\gamma = \alpha \circ \varphi$
considérées comme des applications $\alpha : [n] \to \{0, 1, 2, \ldots\}$
et $\gamma : [m] \to \{0, 1, 2, \ldots\}$. On a le diagramme
$$
\xymatrix{
[m] \ar[r]_\varphi \ar[d]_\gamma &
[n] \ar[r]_\psi \ar[d]_\alpha &
[p] \ar[d]^\beta \\
\Im(\gamma) \ar[r] &
\Im(\alpha) \ar[r] &
[k(\beta)]
}
$$
Nous allons montrer que les restrictions des applications
$S(A_\bullet)(\varphi) \circ S(A_\bullet)(\psi)$ et
$S(A_\bullet)(\psi \circ \varphi)$
au facteur direct $A_{k(\beta)} \cdot \beta$ coïncident.
Il y a plusieurs cas à considérer :
\begin{enumerate}
\item Supposons que $\alpha \not \in I_n$ ; la restriction de
$S(A_\bullet)(\psi)$ à $A_{k(\beta)} \cdot \beta$ est donc nulle.
Alors soit $\gamma \not \in I_m$, soit
$[k(\gamma)] = \Im(\gamma) \subset \Im(\alpha) \subset [k(\beta)]$
et le sous-ensemble $\Im(\alpha)$ de $[k(\beta)]$ présente une lacune, de sorte que
$k(\gamma) < k(\beta) - 1$. Dans les deux cas, la restriction de
$S(A_\bullet)(\psi \circ \varphi)$ à $A_{k(\beta)} \cdot \beta$
est elle aussi nulle.
\item Supposons que $\alpha \in I_n$ et $k(\alpha) < k(\beta) - 1$ ; la restriction de
$S(A_\bullet)(\psi)$ à $A_{k(\beta)} \cdot \beta$ est donc nulle.
Alors soit $\gamma \not \in I_m$, soit
$[k(\gamma)] \subset [k(\alpha)] \subset [k(\beta)]$
et il s'ensuit que $k(\gamma) < k(\beta) - 1$. Dans les deux cas, la restriction de
$S(A_\bullet)(\psi \circ \varphi)$ à $A_{k(\beta)} \cdot \beta$
est elle aussi nulle.
\item Supposons que $\alpha \in I_n$ et $k(\alpha) = k(\beta)$ ; la restriction de
$S(A_\bullet)(\psi)$ à $A_{k(\beta)} \cdot \beta$ est alors
l'application identité de $A_{k(\beta)} \cdot \beta$ vers
$A_{k(\alpha)} \cdot \alpha$. Dans ce cas, comme $\Im(\alpha) =[k(\beta)]$,
la règle qui décrit la restriction de $S(A_\bullet)(\psi \circ \varphi)$
au facteur direct $A_{k(\beta)} \cdot \beta$ est exactement la même
que celle qui décrit la restriction de $S(A_\bullet)(\varphi)$
au facteur direct $A_{k(\alpha)} \cdot \alpha$ ; elles coïncident donc.
\item Supposons que $\alpha \in I_n$ et $k(\alpha) = k(\beta) - 1$ ;
la restriction de $S(A_\bullet)(\psi)$ à $A_{k(\beta)} \cdot \beta$
est alors donnée par $(-1)^{k(\beta)}d_{A, k(\beta)}$ vers $A_{k(\alpha)} \cdot \alpha$.
Distinguons plusieurs sous-cas :
\begin{enumerate}
\item Si $\gamma \not \in I_m$, la restriction de
$S(A_\bullet)(\psi \circ \varphi)$ au facteur direct $A_{k(\beta)} \cdot \beta$
et celle de $S(A_\bullet)(\varphi)$ au facteur direct
$A_{k(\alpha)} \cdot \alpha$ sont toutes deux nulles ; elles coïncident donc.
\item Si $\gamma \in I_m$, mais $k(\gamma) < k(\alpha) - 1$,
les deux restrictions sont encore nulles et coïncident donc.
\item Si $\gamma \in I_m$ et $k(\gamma) = k(\alpha)$, alors
$\Im(\gamma) = \Im(\alpha)$. Dans ce cas, la restriction de
$S(A_\bullet)(\psi \circ \varphi)$
au facteur direct $A_{k(\beta)} \cdot \beta$ est donnée par
$(-1)^{k(\beta)}d_{A, k(\beta)}$ vers $A_{k(\gamma)} \cdot \gamma$
et la restriction de $S(A_\bullet)(\varphi)$ au facteur direct
$A_{k(\alpha)} \cdot \alpha$ est l'application identité
$A_{k(\alpha)} \cdot \alpha \to A_{k(\gamma)} \cdot \gamma$.
Les deux applications coïncident donc.
\item Enfin, si $\gamma \in I_m$ et $k(\gamma) = k(\alpha) - 1$, la restriction de
$S(A_\bullet)(\varphi)$ au facteur direct
$A_{k(\alpha)} \cdot \alpha$ est donnée par $(-1)^{k(\alpha)} d_{A, k(\alpha)}$
comme application
$A_{k(\alpha)} \cdot \alpha \to A_{k(\gamma)} \cdot \gamma$.
Comme $A_\bullet$ est un complexe, la composée
$A_{k(\beta)} \cdot \beta \to
A_{k(\alpha)} \cdot \alpha \to A_{k(\gamma)} \cdot \gamma$
est nulle, ce qui coïncide avec la restriction de
$S(A_\bullet)(\psi \circ \varphi)$
au facteur direct $A_{k(\beta)} \cdot \beta$,
car $k(\gamma) = k(\beta) - 2 < k(\beta) - 1$.
\end{enumerate}
\end{enumerate}
Ainsi, $S(A_\bullet)$ est un objet simplicial de $\mathcal{A}$.

\medskip\noindent
Construisons un isomorphisme $A_\bullet \to N(S(A_\bullet))$
fonctoriel en $A_\bullet$. Rappelons que
$$
S(A_\bullet) =  N(S(A_\bullet)) \oplus D(S(A_\bullet))
$$
comme complexes de chaînes, d'après le lemme
\ref{lemma-decompose-associated-complexes}.
D'autre part, il résulte de la remarque \ref{remark-degenerate-subcomplex}
et de la construction de $S(A_\bullet)$ que
$$
D(S(A_\bullet))_n =
\bigoplus\nolimits_{\alpha \in I_n,\ k(\alpha) < n} A_{k(\alpha)} \cdot \alpha
\subset
\bigoplus\nolimits_{\alpha \in I_n} A_{k(\alpha)} \cdot \alpha
$$
Or, si $\alpha \in I_n$, alors
$k(\alpha) \geq n \Leftrightarrow \alpha = \text{id}_{[n]} : [n] \to [n]$.
Ainsi, le facteur direct $A_n \cdot \text{id}_{[n]}$ de $S(A_\bullet)_n$
est un supplémentaire de $D(S(A_\bullet))_n$.
Toutes les applications $d^n_i : S(A_\bullet)_n \to S(A_\bullet)_{n - 1}$
s'annulent sur le facteur direct $A_n \cdot \text{id}_{[n]}$,
sauf $d^n_n$, qui y induit $(-1)^n d_{A, n}$ de
$A_n \cdot \text{id}_{[n]}$ vers $A_{n - 1} \cdot \text{id}_{[n - 1]}$.
Nous en concluons que $A_n \cdot \text{id}_{[n]}$ est nécessairement
égal au facteur direct $N(S(A_\bullet))_n$ et, de plus, que
la restriction de la différentielle
$d_n = \sum (-1)^id^n_i : S(A_\bullet)_n \to S(A_\bullet)_{n - 1}$
au facteur direct $A_n \cdot \text{id}_{[n]}$ est bien celle que nous cherchons.

\medskip\noindent
Enfin, il nous reste à montrer que $S \circ N$ est isomorphe au
foncteur identité. Soit $U$ un objet simplicial de $\mathcal{A}$.
Nous pouvons alors définir une application évidente
$$
S(N(U))_n = \bigoplus\nolimits_{\alpha \in I_n} N(U)_{k(\alpha)} \cdot \alpha
\longrightarrow
U_n
$$
en utilisant $U(\alpha) : N(U)_{k(\alpha)} \to U_n$ sur le facteur direct
correspondant à $\alpha$. D'après la définition \ref{definition-split},
c'est un isomorphisme.
Pour achever la démonstration, il faut vérifier sa compatibilité avec
les applications des objets simpliciaux. Soient donc $\varphi : [m] \to [n]$
et $\alpha \in I_n$. Posons $\beta = \alpha \circ \varphi$.
On a le diagramme
$$
\xymatrix{
[m] \ar[r]_\varphi \ar[d]_\beta &
[n] \ar[d]_\alpha \\
\Im(\beta) \ar[r] &
[k(\alpha)]
}
$$
Il y a plusieurs cas à considérer :
\begin{enumerate}
\item Supposons que $\beta \not \in I_m$. Il existe alors un indice
$0 \leq j < k(\alpha)$ tel que $j \not \in \Im(\alpha \circ \varphi)$ ;
nous pouvons donc choisir une factorisation
$\alpha \circ \varphi = \delta^{k(\alpha)}_j \circ \psi$
pour une certaine $\psi : [m] \to [k(\alpha) - 1]$.
Il s'ensuit que $U(\varphi)$ s'annule sur l'image du facteur direct
$N(U)_{k(\alpha)} \cdot \alpha$, car $U(\varphi) \circ U(\alpha) =
U(\alpha \circ \varphi) = U(\psi) \circ d^{k(\alpha)}_j$
est nulle sur $N(U)_{k(\alpha)}$ par construction de $N$.
C'est précisément la règle donnée ci-dessus pour $S(N(U))$.
\item Supposons que $\beta \in I_m$ et $k(\beta) < k(\alpha) - 1$.
Nous raisonnons exactement comme dans le cas (1), avec $j = k(\alpha) - 1$.
\item Supposons que $\beta \in I_m$ et $k(\beta) = k(\alpha)$.
Le facteur direct $N(U)_{k(\alpha)} \cdot \alpha$ est alors envoyé par
l'identité sur le facteur direct $N(U)_{k(\beta)} \cdot \beta$.
C'est aussi l'effet de $U(\varphi)$, puisque, dans ce cas,
$U(\varphi) \circ U(\alpha) = U(\beta)$.
\item Supposons que $\beta \in I_m$ et $k(\beta) = k(\alpha) - 1$.
Nous utilisons alors la différentielle $(-1)^{k(\alpha)} d_{N(U), k(\alpha)}$
pour envoyer le facteur direct $N(U)_{k(\alpha)} \cdot \alpha$
sur le facteur direct $N(U)_{k(\beta)} \cdot \beta$.
D'autre part, comme $\Im(\beta) = [k(\beta)]$ dans ce cas, nous avons
$\alpha \circ \varphi = \delta^{k(\alpha)}_{k(\alpha)} \circ \beta$.
Ainsi, la composée de $U(\varphi)$ avec la restriction de
$U(\alpha)$ à $N(U)_{k(\alpha)}$ est égale à
$U(\beta)$ précomposée avec la restriction de
$d^{k(\alpha)}_{k(\alpha)}$ à $N(U)_{k(\alpha)}$.
Puisque $d_{N(U), k(\alpha)} = \sum (-1)^i d^{k(\alpha)}_i$
et puisque $d^{k(\alpha)}_i$ s'annule sur
$N(U)_{k(\alpha)}$ lorsque $i < k(\alpha)$,
nous obtenons bien l'égalité voulue.
\end{enumerate}
Ceci achève la démonstration du théorème.
\end{proof}
\section{Dold--Kan pour les objets cosimpliciaux}
\label{section-dold-kan-cosimplicial}

\noindent
Soit $\mathcal{A}$ une catégorie abélienne.
D'après Homologie, lemme \ref{homology-lemma-abelian-opposite},
$\mathcal{A}^{opp}$ est elle aussi abélienne. Il résulte
formellement des définitions que
$$
\text{CoSimp}(\mathcal{A}) = \text{Simp}(\mathcal{A}^{opp})^{opp}.
$$
Ainsi, le théorème de Dold--Kan
(théorème \ref{theorem-dold-kan})
implique que $\text{CoSimp}(\mathcal{A})$ est équivalente à
la catégorie
$\text{Ch}_{\geq 0}(\mathcal{A}^{opp})^{opp}$. D'autre part, il
résulte formellement des définitions que
$$
\text{CoCh}_{\geq 0}(\mathcal{A}) =
\text{Ch}_{\geq 0}(\mathcal{A}^{opp})^{opp}.
$$
En composant ces flèches, nous obtenons une équivalence
$$
Q :
\text{CoSimp}(\mathcal{A})
\longrightarrow
\text{CoCh}_{\geq 0}(\mathcal{A}).
$$
Dans cette section, nous décrivons $Q$.

\medskip\noindent
Définissons d'abord le
{\it complexe de cochaînes $s(U)$ associé à un objet
cosimplicial $U$}. C'est le complexe de cochaînes dont les termes sont
nuls aux degrés négatifs et tel que $s(U)^n = U_n$ pour $n \geq 0$.
Comme différentielles, nous utilisons les applications
$d^n : s(U)^n \to s(U)^{n + 1}$ définies par
$d^n = \sum_{i = 0}^{n + 1} (-1)^i \delta^{n + 1}_i$.
Autrement dit, le complexe $s(U)$ est de la forme
$$
\xymatrix{
0 \ar[r] &
U_0 \ar[rr]^{\delta^1_0 - \delta^1_1} & &
U_1 \ar[rr]^{\delta^2_0 - \delta^2_1 + \delta^2_2} & &
U_2 \ar[r] &
\ldots
}
$$
On l'appelle aussi parfois le {\it complexe de Moore} associé
à $U$.

\medskip\noindent
D'autre part, étant donné un objet
cosimplicial $U$ de $\mathcal{A}$, posons
$Q(U)^0 = U_0$ et
$$
Q(U)^n = \Coker(
\xymatrix{
\bigoplus_{i = 0}^{n - 1} U_{n - 1} \ar[r]^-{\delta^n_i} &
U_n
}).
$$
La différentielle $d^n : Q(U)^n \to Q(U)^{n + 1}$
est induite par $(-1)^{n + 1}\delta^{n + 1}_{n + 1}$, c'est-à-dire
en insérant le morphisme
$(-1)^{n + 1}\delta^{n + 1}_{n + 1}$
dans un diagramme
commutatif
$$
\xymatrix{
U_n \ar[rr]_{(-1)^{n + 1}\delta^{n + 1}_{n + 1}} \ar[d] & &
U_{n + 1} \ar[d] \\
Q(U)^n \ar[rr]^{d^n} & &
Q(U)^{n + 1}.
}
$$
Nous laissons au lecteur le soin de vérifier que ce diagramme a bien
un sens, c'est-à-dire que l'image de $\delta^n_i$ est envoyée dans
le noyau de la flèche verticale de droite pour $i = 0, \ldots, n - 1$.
(C'est l'énoncé dual du lemme \ref{lemma-N-d-in-N}.)
Ainsi, notre complexe de cochaînes $Q(U)$ est de la forme
$$
0 \to Q(U)^0 \to Q(U)^1 \to Q(U)^2 \to \ldots
$$
On l'appelle le {\it complexe de cochaînes normalisé associé
à $U$}.
L'énoncé dual du théorème de Dold--Kan \ref{theorem-dold-kan} est le suivant.

\begin{lemma}
\label{lemma-dual-dold-kan}
Soit $\mathcal{A}$ une catégorie abélienne.
\begin{enumerate}
\item Le foncteur
$s : \text{CoSimp}(\mathcal{A}) \to \text{CoCh}_{\geq 0}(\mathcal{A})$
est exact.
\item Les applications $s(U)^n \to Q(U)^n$ définissent un morphisme
de complexes de cochaînes.
\item Il existe une décomposition en somme directe fonctorielle
$s(U) = D(U) \oplus Q(U)$ dans $\text{CoCh}_{\geq 0}(\mathcal{A})$.
\item Le foncteur $Q$ est exact.
\item Le morphisme de complexes $s(U) \to Q(U)$ est un quasi-isomorphisme.
\item Le foncteur $U \mapsto Q(U)^\bullet$ définit
une équivalence de catégories
$\text{CoSimp}(\mathcal{A}) \to \text{CoCh}_{\geq 0}(\mathcal{A})$.
\end{enumerate}
\end{lemma}

\begin{proof}
Omis. Ces résultats sont toutefois les énoncés exactement duaux des
lemmes \ref{lemma-s-exact}, \ref{lemma-map-associated-complexes},
\ref{lemma-decompose-associated-complexes},
\ref{lemma-N-exact}, \ref{lemma-quasi-isomorphism}, et
du théorème \ref{theorem-dold-kan}.
\end{proof}

\section{Homotopies}
\label{section-homotopy}

\noindent
Considérons les ensembles simpliciaux $\Delta[0]$ et $\Delta[1]$.
Rappelons qu'il existe deux morphismes
$$
e_0, e_1 : \Delta[0] \longrightarrow \Delta[1],
$$
provenant des morphismes $[0] \to [1]$ qui envoient
$0$ sur un élément de $[1] = \{0, 1\}$. Rappelons aussi que chaque
ensemble $\Delta[1]_k$ est fini. Par conséquent, si la catégorie
$\mathcal{C}$ possède les coproduits finis, nous pouvons
former le produit
$$
U \times \Delta[1]
$$
pour tout objet simplicial $U$ de $\mathcal{C}$ ; voir la
définition \ref{definition-product-with-simplicial-set}.
Notons que $\Delta[0]$ possède la propriété suivante : $\Delta[0]_k = \{*\}$
est un singleton pour tout $k \geq 0$. Ainsi, $U \times \Delta[0]
= U$. Les morphismes $e_0, e_1$ ci-dessus donnent donc lieu à des morphismes
$$
e_0, e_1 : U \to U \times \Delta[1].
$$

\begin{definition}
\label{definition-homotopy}
Soit $\mathcal{C}$ une catégorie possédant les coproduits finis.
Supposons que $U$ et $V$ soient deux objets simpliciaux de $\mathcal{C}$.
Soient $a, b : U \to V$ deux morphismes.
\begin{enumerate}
\item Nous disons qu'un morphisme
$$
h : U \times \Delta[1] \longrightarrow V
$$
est une {\it homotopie de $a$ vers $b$} si $a = h \circ e_0$ et
$b = h \circ e_1$.
\item Nous disons que les morphismes $a$ et $b$ sont {\it homotopes}, ou
{\it dans la même classe d'homotopie},
s'il existe une suite de morphismes $a = a_0, a_1, \ldots, a_n = b$
de $U$ vers $V$ telle que, pour chaque $i = 1, \ldots, n$, il existe soit
une homotopie de $a_{i - 1}$ vers $a_i$, soit une homotopie
de $a_i$ vers $a_{i - 1}$.
\end{enumerate}
\end{definition}

\noindent
La relation « il existe une homotopie de $a$ vers $b$ » n'est en général ni
transitive ni symétrique ; nous verrons dans l'exemple
\ref{example-trivial-homotopy} qu'elle est réflexive. Bien entendu, « être
homotopes » est une relation d'équivalence sur l'ensemble $\Mor(U, V)$ ;
c'est la relation d'équivalence engendrée par la relation
« il existe une homotopie de $a$ vers $b$ ».
On peut en fait définir des homotopies entre des couples d'applications
d'objets simpliciaux dans une catégorie quelconque. Nous le ferons dans la
remarque \ref{remark-homotopy-better}, après avoir explicité ce que signifie
la donnée d'un morphisme
$h : U \times \Delta[1] \to V$.

\medskip\noindent
Soit $\mathcal{C}$ une catégorie possédant les coproduits finis.
Soient $U$, $V$ des objets simpliciaux de $\mathcal{C}$.
Soient $a, b : U \to V$ des morphismes. Supposons en outre
que $h : U \times \Delta[1] \to V$ soit une homotopie
de $a$ vers $b$. Pour tout $n \geq 0$, écrivons
$$
\Delta[1]_n = \{\alpha^n_0, \ldots, \alpha^n_{n + 1}\}
$$
où $\alpha^n_i : [n] \to [1]$ est l'application
définie par
$$
\alpha^n_i(j)
=
\left\{
\begin{matrix}
0 & \text{si} & j < i\\
1 & \text{si} & j \geq i
\end{matrix}
\right.
$$
Ainsi,
$$
h_n : (U \times \Delta[1])_n = \coprod U_n \cdot \alpha^n_i
\longrightarrow
V_n
$$
possède une composante $h_{n, i} : U_n \to V_n$ qui est sa restriction au
facteur direct correspondant à $\alpha^n_i$, pour tout
$i = 0, \ldots, n + 1$.

\begin{lemma}
\label{lemma-relations-homotopy}
Dans la situation ci-dessus, les relations suivantes sont satisfaites :
\begin{enumerate}
\item On a $h_{n, 0} = b_n$ et $h_{n, n + 1} = a_n$.
\item On a $d^n_j \circ h_{n, i} = h_{n - 1, i - 1} \circ d^n_j$
pour $i > j$.
\item On a $d^n_j \circ h_{n, i} = h_{n - 1, i} \circ d^n_j$
pour $i \leq j$.
\item On a $s^n_j \circ h_{n, i} = h_{n + 1, i + 1} \circ s^n_j$
pour $i > j$.
\item On a $s^n_j \circ h_{n, i} = h_{n + 1, i} \circ s^n_j$
pour $i \leq j$.
\end{enumerate}
Réciproquement, tout système d'applications $h_{n, i}$ satisfaisant aux
propriétés ci-dessus définit un morphisme
$h$ qui est une homotopie de $a$ vers $b$.
\end{lemma}

\begin{proof}
Omis. Pour démontrer la dernière assertion, on utilise le fait que,
d'après le lemme \ref{lemma-face-degeneracy-category},
se donner un morphisme d'objets simpliciaux revient à se donner
une suite de morphismes $h_n$ commutant avec tous les $d^n_j$ et $s^n_j$.
\end{proof}

\begin{example}
\label{example-trivial-homotopy}
Supposons, dans la situation ci-dessus, que $a = b$. Il existe alors une
homotopie {\it triviale} de $a$ vers $b$, à savoir celle pour laquelle
$h_{n, i} = a_n = b_n$.
\end{example}

\begin{remark}
\label{remark-homotopy-better}
Soit $\mathcal{C}$ une catégorie quelconque (sans aucune hypothèse).
Soient $U$ et $V$ des objets simpliciaux de $\mathcal{C}$, et soient
$a, b : U \to V$ des morphismes d'objets simpliciaux de $\mathcal{C}$.
Une {\it homotopie de $a$ vers $b$} est la donnée de
morphismes\footnote{Dans la littérature, on utilise souvent les applications
$h_{n + 1, i} \circ s_i : U_n \to V_{n + 1}$ à la place des
applications $h_{n, i}$. Les relations qu'elles satisfont sont bien entendu
différentes de celles du lemme \ref{lemma-relations-homotopy}.}
$h_{n, i} : U_n \to V_n$, for $n \geq 0$, $i = 0, \ldots, n + 1$
satisfaisant aux relations du lemme \ref{lemma-relations-homotopy}.
Comme dans la définition \ref{definition-homotopy}, nous disons que les
morphismes $a$ et $b$ sont {\it homotopes} s'il existe une suite de morphismes
$a = a_0, a_1, \ldots, a_n = b$ de $U$ vers $V$ telle que, pour chaque
$i = 1, \ldots, n$, il existe soit une homotopie de $a_{i - 1}$ vers $a_i$,
soit une homotopie de $a_i$ vers $a_{i - 1}$.
Il est clair que si $F : \mathcal{C} \to \mathcal{C}'$ est un foncteur
quelconque et si $\{h_{n, i}\}$ est une homotopie de $a$ vers $b$, alors
$\{F(h_{n, i})\}$ est une homotopie de $F(a)$ vers $F(b)$.
De même, si $a$ et $b$ sont homotopes, alors $F(a)$ et $F(b)$
le sont aussi.
Comme le lemme affirme que la nouvelle notion coïncide avec l'ancienne
lorsque les coproduits finis existent, nous en déduisons notamment que les
foncteurs préservent la notion initiale dès que les deux catégories possèdent
les coproduits finis.
\end{remark}

\begin{remark}
\label{remark-homotopy-pre-post-compose}
Soit $\mathcal{C}$ une catégorie quelconque. Supposons que deux morphismes
$a, a' : U \to V$ d'objets simpliciaux soient homotopes. Alors, pour tout
morphisme $b : V \to W$, les deux applications
$b \circ a, b \circ a' : U \to W$ sont homotopes.
De même, pour tout morphisme $c : X \to U$, les deux applications
$a \circ c, a' \circ c : X \to V$ sont homotopes. En fait, les applications
$b \circ a \circ c, b \circ a' \circ c : X \to W$ sont homotopes.
En effet, si les applications $h_{n, i} : U_n \to V_n$ définissent une
homotopie de $a$ vers $a'$, alors les applications
$b \circ h_{n, i} \circ c$ définissent une homotopie de
$b \circ a \circ c$ vers $b \circ a' \circ c$. Nous obtenons ainsi une
nouvelle catégorie $\text{hSimp}(\mathcal{C})$, qui a les mêmes objets que
$\text{Simp}(\mathcal{C})$, mais dont les morphismes sont les classes
d'homotopie de morphismes de $\text{Simp}(\mathcal{C})$.
Il existe donc un foncteur canonique
$$
\text{Simp}(\mathcal{C})
\longrightarrow
\text{hSimp}(\mathcal{C})
$$
qui est essentiellement surjectif et surjectif sur les ensembles de morphismes.
\end{remark}

\begin{definition}
\label{definition-homotopy-equivalent}
Soient $U$ et $V$ deux objets simpliciaux d'une catégorie $\mathcal{C}$.
Nous disons qu'un morphisme $a : U \to V$ est une
{\it équivalence d'homotopie} s'il existe un morphisme $b : V \to U$ tel que
$a \circ b$ soit homotope à $\text{id}_V$ et $b \circ a$ homotope à
$\text{id}_U$. Nous disons que $U$ et $V$ ont le {\it même type d'homotopie}
s'il existe une équivalence d'homotopie $a : U \to V$.
\end{definition}

\begin{example}
\label{example-simplex-contractible}
L'ensemble simplicial $\Delta[m]$ a le même type d'homotopie que $\Delta[0]$.
En effet, considérons l'unique morphisme $f : \Delta[m] \to \Delta[0]$ et
le morphisme $g : \Delta[0] \to \Delta[m]$ donné par l'inclusion du dernier
$0$-simplexe de $\Delta[m]$. On a $f \circ g = \text{id}$.
Nous allons construire une homotopie
$h : \Delta[m] \times \Delta[1] \to \Delta[m]$
de $\text{id}_{\Delta[m]}$ vers $g \circ f$. L'homotopie $h$ est donnée par les applications
$$
\Mor_\Delta([n], [m]) \times \Mor_\Delta([n], [1]) \to \Mor_\Delta([n], [m])
$$
qui envoient $(\varphi, \alpha)$ sur
$$
k \mapsto
\left\{
\begin{matrix}
\varphi(k) & \text{si} & \alpha(k) = 0 \\
m & \text{si} & \alpha(k) = 1
\end{matrix}
\right.
$$
Notons que ceci ne fonctionne que parce que nous avons choisi pour $g$
l'inclusion du dernier $0$-simplexe. Si nous avions pris pour $g$ l'inclusion du
premier $0$-simplexe, nous aurions pu construire une homotopie de
$g \circ f$ vers $\text{id}_{\Delta[m]}$. Cela illustre l'asymétrie propre
aux homotopies dans la catégorie des ensembles simpliciaux.
\end{example}

\noindent
Le lemme suivant affirme que $U \times \Delta[1]$
a le même type d'homotopie que $U$.

\begin{lemma}
\label{lemma-contractible}
Soit $\mathcal{C}$ une catégorie possédant les coproduits finis.
Soit $U$ un objet simplicial de $\mathcal{C}$.
Considérons les applications $e_1, e_0 : U \to U \times \Delta[1]$
et $\pi : U \times \Delta[1] \to U$ ; voir le
lemme \ref{lemma-back-to-U}.
\begin{enumerate}
\item On a $\pi \circ e_1 = \pi \circ e_0 = \text{id}_U$, et
\item les morphismes $\text{id}_{U \times \Delta[1]}$
et $e_0 \circ \pi$ sont homotopes.
\item Les morphismes $\text{id}_{U \times \Delta[1]}$
et $e_1 \circ \pi$ sont homotopes.
\end{enumerate}
\end{lemma}

\begin{proof}
La première assertion est immédiate. Considérons l'application d'ensembles simpliciaux
$\Delta[1] \times \Delta[1] \longrightarrow \Delta[1]$
qui, en degré $n$, associe à un couple $(\beta_1, \beta_2)$,
où $\beta_i : [n] \to [1]$, le morphisme
$\beta : [n] \to [1]$ défini par la règle
$$
\beta(i) = \max\{\beta_1(i), \beta_2(i)\}.
$$
Il s'agit d'un morphisme d'ensembles simpliciaux, car l'action
$\Delta[1](\varphi) : \Delta[1]_n \to \Delta[1]_m$
de $\varphi : [m] \to [n]$ se fait par précomposition.
Avec la notation $\Delta[1]_n = \{\alpha^n_0, \ldots, \alpha^n_{n + 1}\}$
introduite ci-dessus, on a $\beta = \alpha^n_0$ si
$\beta_2 = \alpha^n_0$, et $\beta = \beta_1$ si
$\beta_2 = \alpha^n_{n + 1}$.
Comme $\alpha^n_0$, resp.\ $\alpha^n_{n + 1}$, est l'application constante
de valeur $1$, resp.\ $0$, il en résulte que le morphisme induit
$$
U \times \Delta[1] \times \Delta[1]
\longrightarrow
U \times \Delta[1]
$$
du lemme \ref{lemma-back-to-U}
est une homotopie de $\text{id}_{U \times \Delta[1]}$ vers $e_1 \circ \pi$.
On procède de même pour $e_0 \circ \pi$ (en utilisant le minimum au lieu du maximum).
\end{proof}

\begin{lemma}
\label{lemma-fibre-products-simplicial-object-w-section}
Soit $f : Y \to X$ un morphisme d'une catégorie $\mathcal{C}$ possédant les
produits fibrés. Supposons que $f$ admette une section $s$. Considérons
l'objet simplicial $U$ construit dans
l'exemple \ref{example-fibre-products-simplicial-object} à partir de $f$.
Le morphisme $U \to U$ qui, en chaque degré, est l'endomorphisme
$(s \circ f)^{n + 1}$ de $Y \times_X \ldots \times_X Y$, donné par
$s \circ f$ sur chaque facteur, est homotope à l'identité de $U$.
En particulier, $U$ a le même type d'homotopie que l'objet simplicial
constant $X$.
\end{lemma}

\begin{proof}
Posons $g^0 = \text{id}_Y$ et $g^1 = s \circ f$.
Nous utilisons les morphismes
\begin{eqnarray*}
Y \times_X \ldots \times_X Y \times \Mor([n], [1])
& \to &
Y \times_X \ldots \times_X Y \\
(y_0, \ldots, y_n) \times \alpha
& \mapsto &
(g^{\alpha(0)}(y_0), \ldots, g^{\alpha(n)}(y_n))
\end{eqnarray*}
où nous adoptons le point de vue du foncteur des points pour définir les applications.
On peut aussi dire que
$h_{n, 0} = \text{id}$, $h_{n, n + 1} = (s \circ f)^{n + 1}$ et
$h_{n, i} = \text{id}_Y^{i + 1} \times (s \circ f)^{n + 1 - i}$.
Nous laissons au lecteur le soin de vérifier qu'elles satisfont aux relations
du lemme \ref{lemma-relations-homotopy}. Elles définissent donc
l'homotopie cherchée. Voir aussi la remarque \ref{remark-homotopy-better},
qui montre qu'aucune autre hypothèse sur la catégorie $\mathcal{C}$
n'est nécessaire.
\end{proof}

\begin{lemma}
\label{lemma-products-homotopy}
Soit $\mathcal{C}$ une catégorie et soit $T$ un ensemble. Pour $t \in T$,
soient $X_t$, $Y_t$ des objets simpliciaux de $\mathcal{C}$. Supposons que
$X = \prod_{t \in T} X_t$ et $Y = \prod_{t \in T} Y_t$ existent.
\begin{enumerate}
\item Si $X_t$ et $Y_t$ ont le même type d'homotopie pour tout $t \in T$
et si $T$ est fini, alors $X$ et $Y$ ont le même type d'homotopie.
\end{enumerate}
Pour $t \in T$, soient $a_t, b_t : X_t \to Y_t$ des morphismes.
Posons $a = \prod a_t : X \to Y$ et $b = \prod b_t : X \to Y$.
\begin{enumerate}
\item[(2)] S'il existe une homotopie de $a_t$ vers $b_t$ pour
tout $t \in T$, alors il existe une homotopie de $a$ vers $b$.
\item[(3)] Si $T$ est fini et si $a_t, b_t : X_t \to Y_t$ sont homotopes
pour tout $t \in T$, alors $a$ et $b$ sont homotopes.
\end{enumerate}
\end{lemma}

\begin{proof}
Si $h_t = (h_{t, n , i})$ est une homotopie de $a_t$ vers $b_t$
(voir la remarque \ref{remark-homotopy-better}), alors
$h = (\prod_t h_{t, n, i})$ est une homotopie de $\prod a_t$ vers
$\prod b_t$. Cela démontre (2).

\medskip\noindent
Démontrons (3). Choisissons $t \in T$. Il existe un entier $n \geq 0$
et une chaîne $a_t = a_{t, 0}, a_{t, 1}, \ldots, a_{t, n} = b_t$ telle que,
pour tout $1 \leq i \leq n$, il existe soit une homotopie de
$a_{t, i - 1}$ vers $a_{t, i}$, soit une homotopie de
$a_{t, i}$ vers $a_{t, i - 1}$. Si $n = 0$, choisissons un autre $t$.
(La démonstration est achevée si $a_t = b_t$ pour tout $t \in T$.)
Supposons donc $n > 0$. D'après l'exemple \ref{example-trivial-homotopy},
il existe une homotopie de $b_{t'}$ vers $b_{t'}$ pour tout
$t' \in T \setminus \{t\}$. Ainsi, d'après (2), il existe soit une homotopie de
$a_{t, n - 1} \times \prod_{t'} b_{t'}$ vers $b$, soit une homotopie de
$b$ vers $a_{t, n - 1} \times \prod_{t'} b_{t'}$.
On peut ainsi diminuer $n$ de $1$. Cela démontre (3).

\medskip\noindent
La partie (1) résulte de la partie (3) et des définitions.
\end{proof}




\section{Homotopies dans les catégories abéliennes}
\label{section-homotopy-abelian}

\noindent
Soit $\mathcal{A}$ une catégorie additive. Soient $U$, $V$ des objets
simpliciaux de $\mathcal{A}$ et soient $a, b : U \to V$ des morphismes.
Supposons en outre que $h : U \times \Delta[1] \to V$ soit une homotopie
de $a$ vers $b$. Montrons que les deux morphismes de complexes de chaînes
$s(a), s(b) : s(U) \longrightarrow s(V)$ sont homotopes au sens de
Homologie, section \ref{homology-section-complexes}.
Avec la notation introduite dans la section \ref{section-homotopy},
nous définissons
$$
s(h)_n : U_n \longrightarrow V_{n + 1}
$$
par la formule
\begin{equation}
\label{equation-homotopy-to-homotopy}
s(h)_n = \sum\nolimits_{i = 0}^n (-1)^{i + 1} h_{n + 1, i + 1} \circ s^n_i.
\end{equation}
Calculons $d_{n + 1} \circ s(h)_n + s(h)_{n - 1} \circ d_n$.
On a d'abord
\begin{eqnarray*}
d_{n + 1} \circ s(h)_n & = &
\sum\nolimits_{j = 0}^{n + 1}
\sum\nolimits_{i = 0}^n
(-1)^{j + i + 1}
d^{n + 1}_j \circ h_{n + 1, i + 1} \circ s^n_i \\
& = &
\sum\nolimits_{1 \leq i + 1 \leq j \leq n + 1}
(-1)^{j + i + 1}
h_{n, i + 1} \circ d^{n + 1}_j \circ s^n_i \\
& &
+
\sum\nolimits_{n \geq i \geq j \geq 0}
(-1)^{i + j + 1}
h_{n, i} \circ d^{n + 1}_j \circ s^n_i \\
& = &
\sum\nolimits_{1 \leq i + 1 < j \leq n + 1}
(-1)^{j + i + 1}
h_{n, i + 1} \circ s^{n - 1}_i \circ d^n_{j - 1} \\
& &
+
\sum\nolimits_{1 \leq i + 1 = j \leq n + 1}
(-1)^{j + i + 1}
h_{n, i + 1} \\
& &
+
\sum\nolimits_{n \geq i = j \geq 0}
(-1)^{i + j + 1}
h_{n, i} \\
& &
+
\sum\nolimits_{n \geq i > j \geq 0}
(-1)^{i + j + 1}
h_{n, i} \circ s^{n - 1}_{i - 1} \circ d^n_j
\end{eqnarray*}
Nous laissons au lecteur le soin de vérifier que la première et la dernière
des quatre sommes s'annulent exactement avec tous les termes de
$$
s(h)_{n - 1} \circ d_n
=
\sum_{i = 0}^{n - 1} \sum_{j = 0}^n
(-1)^{i + 1 + j} h_{n, i + 1} \circ s^{n - 1}_i \circ d^n_j.
$$
Nous obtenons donc
\begin{eqnarray*}
d_{n + 1} \circ s(h)_n + s(h)_{n - 1} \circ d_n
& = &
\sum_{j = 1}^{n + 1} (-1)^{2j} h_{n, j}
+
\sum_{i = 0}^n (-1)^{2i + 1} h_{n, i} \\
& = &
h_{n, n + 1} - h_{n , 0} \\
& = &
a_n - b_n
\end{eqnarray*}
comme souhaité.

\begin{lemma}
\label{lemma-homotopy-s-N}
Soit $\mathcal{A}$ une catégorie additive. Soient $a, b : U \to V$ des
morphismes d'objets simpliciaux de $\mathcal{A}$. Si $a$, $b$ sont homotopes,
alors $s(a), s(b) : s(U) \to s(V)$ sont des applications homotopes de
complexes de chaînes. Si $\mathcal{A}$ est abélienne, alors
$N(a), N(b) : N(U) \to N(V)$ sont également des applications homotopes de
complexes de chaînes.
\end{lemma}

\begin{proof}
Nous pouvons choisir une suite $a = a_0, a_1, \ldots, a_n = b$ de morphismes
de $U$ vers $V$ telle que, pour chaque $i = 1, \ldots, n$, il existe soit une
homotopie de $a_i$ vers $a_{i - 1}$, soit une homotopie de
$a_{i - 1}$ vers $a_i$. Le calcul ci-dessus montre qu'alors soit
$s(a_i)$ est homotope à $s(a_{i - 1})$ comme application de complexes de
chaînes, soit $s(a_{i - 1})$ est homotope à $s(a_i)$. Ces deux assertions
sont bien entendu équivalentes ; de plus, être homotopes est une relation
d'équivalence sur l'ensemble des applications de complexes de chaînes ; voir
Homologie, section \ref{homology-section-complexes}.
Cela démontre que $s(a)$ et $s(b)$ sont homotopes comme applications de
complexes de chaînes.

\medskip\noindent
Passons à $N(a)$ et $N(b)$. Il résulte du
lemme \ref{lemma-decompose-associated-complexes} que
$N(a)$, $N(b)$ sont les composées
$$
N(U) \to s(U) \to s(V) \to N(V)
$$
où l'on utilise $s(a)$, $s(b)$ au milieu. L'assertion résulte donc de
Homologie, lemme \ref{homology-lemma-compose-homotopy}.
\end{proof}

\begin{lemma}
\label{lemma-homotopy-equivalence-s-N}
Soit $\mathcal{A}$ une catégorie additive. Soit $a : U \to V$ un morphisme
d'objets simpliciaux de $\mathcal{A}$. Si $a$ est une équivalence d'homotopie,
alors $s(a) : s(U) \to s(V)$ est une équivalence d'homotopie de complexes de
chaînes. Si, de plus, $\mathcal{A}$ est abélienne, alors
$N(a) : N(U) \to N(V)$ est lui aussi une équivalence d'homotopie de complexes de chaînes.
\end{lemma}

\begin{proof}
Omis. Voir le lemme \ref{lemma-homotopy-s-N} ci-dessus.
\end{proof}




\section{Homotopies et objets cosimpliciaux}
\label{section-homotopy-cosimplicial}

\noindent
Soit $\mathcal{C}$ une catégorie possédant les produits finis. Soit $V$ un
objet cosimplicial, et considérons $\Hom(\Delta[1], V)$ ; voir la
section \ref{section-hom-from-simplicial-sets-into-cosimplicial}.
Les morphismes $e_0, e_1 : \Delta[0] \to \Delta[1]$
induisent deux morphismes $e_0, e_1 : \Hom(\Delta[1], V) \to V$.

\begin{definition}
\label{definition-homotopy-cosimplicial}
Soit $\mathcal{C}$ une catégorie possédant les produits finis.
Soient $U$ et $V$ deux objets cosimpliciaux de $\mathcal{C}$.
Soient $a, b : U \to V$ deux morphismes d'objets cosimpliciaux
de $\mathcal{C}$.
\begin{enumerate}
\item Nous disons qu'un morphisme
$$
h : U \longrightarrow \Hom(\Delta[1], V)
$$
tel que $a = e_0 \circ h$ et $b = e_1 \circ h$ est une
{\it homotopie de $a$ vers $b$}.
\item Nous disons que $a$ et $b$ sont {\it homotopes}, ou
{\it dans la même classe d'homotopie}, s'il existe une suite
$a = a_0, a_1, \ldots, a_n = b$ de morphismes de $U$ vers $V$
telle que, pour chaque $i = 1, \ldots, n$, il existe soit une
homotopie de $a_i$ vers $a_{i - 1}$, soit une homotopie
de $a_{i - 1}$ vers $a_i$.
\end{enumerate}
\end{definition}

\noindent
C'est la notion duale de celle introduite pour les objets simpliciaux dans la
section \ref{section-homotopy}. Pour l'expliquer, considérons une
homotopie $h : U \to \Hom(\Delta[1], V)$ de $a$ vers $b$ comme dans la
définition. Rappelons que $\Delta[1]_n$ est un ensemble fini.
La composante de degré $n$ de $h$ est un morphisme
$$
h_n = (h_{n, \alpha}) :
U
\longrightarrow
\Hom(\Delta[1], V)_n =
\prod\nolimits_{\alpha \in \Delta[1]_n} V_n
$$
Les morphismes $h_{n, \alpha} : U_n \to V_n$ de $\mathcal{C}$ possèdent la
propriété suivante : pour tout morphisme $f : [n] \to [m]$
de $\Delta$, on a
\begin{equation}
\label{equation-property-homotopy-cosimplicial}
h_{m, \alpha} \circ U(f) = V(f) \circ h_{n, \alpha \circ f}
\end{equation}
De plus, la condition $a = e_0 \circ h$ signifie que
$a_n = h_{n, 0 : [n] \to [1]}$, où $0 : [n] \to [1]$ est l'application
constante de valeur $0$. De même, la condition $b = e_1 \circ h$ signifie que
$b_n = h_{n, 1 : [n] \to [1]}$, où $1 : [n] \to [1]$ est l'application
constante de valeur $1$.
Réciproquement, étant donnée une famille de morphismes $\{h_{n, \alpha}\}$
telle que (\ref{equation-property-homotopy-cosimplicial}) soit satisfaite
pour tout morphisme $f$ de $\Delta$ et telle que
$a_n = h_{n, 0 : [n] \to [1]}$ et $b_n = h_{n, 1 : [n] \to [1]}$
pour tout $n \geq 0$, nous obtenons une homotopie $h$ de $a$ vers $b$ en posant
$h = \prod_{\alpha \in \Delta[1]_n} h_{n, \alpha}$.

\begin{remark}
\label{remark-homotopy-cosimplicial-better}
Soit $\mathcal{C}$ une catégorie quelconque (sans aucune hypothèse).
Soient $U$ et $V$ des objets cosimpliciaux de $\mathcal{C}$, et soient
$a, b : U \to V$ des morphismes d'objets cosimpliciaux de $\mathcal{C}$.
Une {\it homotopie de $a$ vers $b$} est la donnée de morphismes
$h_{n, \alpha} : U_n \to V_n$, for $n \geq 0$, $\alpha \in \Delta[1]_n$
satisfaisant à (\ref{equation-property-homotopy-cosimplicial})
pour tout morphisme $f$ de $\Delta$ et tels que
$a_n = h_{n, 0 : [n] \to [1]}$ et $b_n = h_{n, 1 : [n] \to [1]}$
pour tout $n \geq 0$.
Comme dans la définition \ref{definition-homotopy-cosimplicial},
nous disons que les morphismes $a$ et $b$ sont {\it homotopes} s'il existe
une suite de morphismes $a = a_0, a_1, \ldots, a_n = b$ de $U$ vers $V$
telle que, pour chaque $i = 1, \ldots, n$, il existe soit une homotopie de
$a_{i - 1}$ vers $a_i$, soit une homotopie de $a_i$ vers $a_{i - 1}$.
Il est clair que si $F : \mathcal{C} \to \mathcal{C}'$ est un foncteur
quelconque et si $\{h_{n, \alpha}\}$ est une homotopie de $a$ vers $b$, alors
$\{F(h_{n, \alpha})\}$ est une homotopie de $F(a)$ vers $F(b)$.
De même, si $a$ et $b$ sont homotopes, alors $F(a)$ et $F(b)$ le sont aussi.
Cette nouvelle notion coïncide avec l'ancienne lorsque les produits finis
existent. Nous en déduisons notamment que les foncteurs préservent la notion
initiale dès que les deux catégories possèdent les produits finis.
\end{remark}

\begin{lemma}
\label{lemma-compare-homotopies}
Soit $\mathcal{C}$ une catégorie. Supposons que $U$ et $V$ soient deux
objets cosimpliciaux de $\mathcal{C}$. Soient $a, b : U \to V$ des morphismes
d'objets cosimpliciaux. Rappelons que $U$, $V$ correspondent à des objets
simpliciaux $U'$, $V'$ de $\mathcal{C}^{opp}$.
De plus, $a, b$ correspondent à des morphismes $a', b' : V' \to U'$.
Les assertions suivantes sont équivalentes :
\begin{enumerate}
\item Il existe une homotopie $h = \{h_{n, \alpha}\}$ de
$a$ vers $b$ au sens de la remarque \ref{remark-homotopy-cosimplicial-better}.
\item Il existe une homotopie $h = \{h_{n, i}\}$ de $a'$ vers $b'$
au sens de la remarque \ref{remark-homotopy-better}.
\end{enumerate}
Ainsi, $a$ est homotope à $b$ au sens de la
remarque \ref{remark-homotopy-cosimplicial-better} si et seulement si
$a'$ est homotope à $b'$ au sens de la
remarque \ref{remark-homotopy-better}.
\end{lemma}

\begin{proof}
Si $\mathcal{C}$ possède les produits finis, alors $\mathcal{C}^{opp}$
possède les coproduits finis et nous pouvons utiliser les
définitions \ref{definition-homotopy-cosimplicial}
et \ref{definition-homotopy} à la place des
remarques \ref{remark-homotopy-cosimplicial-better} et
\ref{remark-homotopy-better}.
Dans ce cas, $h : U \to \Hom(\Delta[1], V)$ équivaut à un morphisme
$h' : \Hom(\Delta[1], V)' \to U'$.
Comme les produits et les coproduits s'échangent eux aussi, il est immédiat
que $(\Hom(\Delta[1], V))' = V' \times \Delta[1]$.
De plus, les versions « primées » des morphismes
$e_0, e_1 : \Hom(\Delta[1], V) \to V$ sont les morphismes
$e_0, e_1 : V' \to V' \times \Delta[1]$.
Ainsi, $e_0 \circ h = a$ se traduit par $h' \circ e_0 = a'$, et
$e_1 \circ h = b$ se traduit de même par $h' \circ e_1 = b'$.
Cela démontre le lemme dans ce cas.

\medskip\noindent
Dans le cas général, il faut traduire les relations données par
(\ref{equation-property-homotopy-cosimplicial}) en celles du
lemme \ref{lemma-relations-homotopy}. Nous omettons les détails.

\medskip\noindent
La dernière assertion résulte formellement de l'équivalence de (1) et (2).
\end{proof}

\begin{lemma}
\label{lemma-functorial-homotopy}
\begin{slogan}
Les foncteurs préservent les morphismes homotopes d'objets (co)simpliciaux.
\end{slogan}
Soient $\mathcal{C}, \mathcal{C}', \mathcal{D}, \mathcal{D}'$ des catégories.
Nous employons la terminologie des remarques
\ref{remark-homotopy-cosimplicial-better} et \ref{remark-homotopy-better}.
\begin{enumerate}
\item Soient $a, b : U \to V$ des morphismes d'objets simpliciaux
de $\mathcal{D}$. Soit $F : \mathcal{D} \to \mathcal{D}'$ un foncteur
covariant. Si $a$ et $b$ sont homotopes, alors $F(a)$, $F(b)$ sont des
morphismes homotopes $F(U) \to F(V)$ d'objets simpliciaux.
\item Soient $a, b : U \to V$ des morphismes d'objets cosimpliciaux
de $\mathcal{C}$. Soit $F : \mathcal{C} \to \mathcal{C}'$ un foncteur
covariant. Si $a$ et $b$ sont homotopes, alors $F(a)$, $F(b)$ sont des
morphismes homotopes $F(U) \to F(V)$ d'objets cosimpliciaux.
\item Soient $a, b : U \to V$ des morphismes d'objets simpliciaux de $\mathcal{D}$.
Soit $F : \mathcal{D} \to \mathcal{C}$ un foncteur contravariant.
Si $a$ et $b$ sont homotopes, alors $F(a)$, $F(b)$ sont des morphismes
homotopes $F(V) \to F(U)$ d'objets cosimpliciaux.
\item Soient $a, b : U \to V$ des morphismes d'objets cosimpliciaux de
$\mathcal{C}$. Soit $F : \mathcal{C} \to \mathcal{D}$ un foncteur
contravariant. Si $a$ et $b$ sont homotopes, alors $F(a)$, $F(b)$ sont des
morphismes homotopes $F(V) \to F(U)$ d'objets simpliciaux.
\end{enumerate}
\end{lemma}

\begin{proof}
D'après le lemme \ref{lemma-compare-homotopies}, nous pouvons transformer
$F$ en un foncteur covariant entre deux catégories ; il nous reste alors à
montrer que ce foncteur préserve les couples d'applications homotopes.
Ceci est expliqué dans la remarque \ref{remark-homotopy-better}.
\end{proof}

\begin{lemma}
\label{lemma-push-outs-simplicial-object-w-section}
Soit $f : X \to Y$ un morphisme d'une catégorie $\mathcal{C}$ possédant les
sommes amalgamées. Supposons qu'il existe un morphisme $s : Y \to X$ tel que
$s \circ f = \text{id}_X$. Considérons l'objet cosimplicial $U$ construit dans
l'exemple \ref{example-push-outs-simplicial-object} à partir de $f$.
Le morphisme $U \to U$ qui, en chaque degré, est l'endomorphisme de
$Y \amalg_X \ldots \amalg_X Y$ donné par $f \circ s$ sur chaque facteur
est homotope à l'identité de $U$. En particulier, $U$ a le même type
d'homotopie que l'objet cosimplicial constant $X$.
\end{lemma}

\begin{proof}
Ce lemme est dual du
lemme \ref{lemma-fibre-products-simplicial-object-w-section}.
Il résulte donc du lemme \ref{lemma-compare-homotopies}.
\end{proof}

\begin{lemma}
\label{lemma-homotopy-s-Q}
\begin{slogan}
Le foncteur de Dold--Kan (cosimplicial) envoie des applications homotopes sur des applications homotopes.
\end{slogan}
Soit $\mathcal{A}$ une catégorie additive. Soient $a, b : U \to V$ des
morphismes d'objets cosimpliciaux de $\mathcal{A}$. Si $a$, $b$ sont
homotopes, alors $s(a), s(b) : s(U) \to s(V)$ sont des applications homotopes
de complexes de cochaînes. Si, de plus, $\mathcal{A}$ est abélienne, alors
$Q(a), Q(b) : Q(U) \to Q(V)$ sont des applications homotopes de complexes de cochaînes.
\end{lemma}

\begin{proof}
Soit $(-)' : \mathcal{A} \to \mathcal{A}^{opp}$ le foncteur contravariant
$A \mapsto A$. D'après le lemme \ref{lemma-compare-homotopies}, les
applications $a'$ et $b'$ sont homotopes.
Le lemme \ref{lemma-homotopy-s-N} montre que $s(a')$ et $s(b')$ sont des
applications homotopes de complexes de chaînes. Comme
$s(a') = (s(a))'$ et $s(b') = (s(b))'$, nous en concluons que
$s(a)$ et $s(b)$ sont eux aussi homotopes, en appliquant le foncteur additif
contravariant $(-)'' : \mathcal{A}^{opp} \to \mathcal{A}$.
Le résultat pour les complexes $Q$ s'obtient de la même manière,
en utilisant $Q(U)' = N(U')$.
\end{proof}

\begin{lemma}
\label{lemma-homotopy-equivalence-s-Q}
Soit $\mathcal{A}$ une catégorie additive. Soit $a : U \to V$ un morphisme
d'objets cosimpliciaux de $\mathcal{A}$. Si $a$ est une équivalence
d'homotopie, alors $s(a) : s(U) \to s(V)$ est une équivalence d'homotopie de
complexes de chaînes. Si, de plus, $\mathcal{A}$ est abélienne, alors
$Q(a) : Q(U) \to Q(V)$ est lui aussi une équivalence d'homotopie de complexes de chaînes.
\end{lemma}

\begin{proof}
Omis. Voir le lemme \ref{lemma-homotopy-s-Q} ci-dessus.
\end{proof}




\section{Autres homotopies dans les catégories abéliennes}
\label{section-backwards-homotopy}

\noindent
Soit $\mathcal{A}$ une catégorie abélienne.
Dans cette section, nous montrons qu'une homotopie entre des
morphismes de $\text{Ch}_{\geq 0}(\mathcal{A})$ provient toujours d'un
morphisme $U \times \Delta[1] \to V$ dans la catégorie des objets
simpliciaux. Cela donnera, en un certain sens, une réciproque du
lemme \ref{lemma-homotopy-s-N}. Nous établissons d'abord quelques résultats
sur les homotopies entre morphismes de complexes de chaînes.

\begin{lemma}
\label{lemma-represent-homotopy}
Soit $\mathcal{A}$ une catégorie abélienne et soit $A$ un complexe de chaînes.
Considérons le foncteur covariant
$$
B \longmapsto
\{
(a, b, h)
\mid
a, b : A \to B\text{ et }h\text{ une homotopie entre }a, b
\}
$$
Il existe un complexe de chaînes $\diamond A$ tel que
$\Mor_{\text{Ch}(\mathcal{A})}(\diamond A, -)$ soit isomorphe au foncteur
affiché. La construction $A \mapsto \diamond A$ est fonctorielle.
\end{lemma}

\begin{proof}
Posons $\diamond A_n = A_n \oplus A_n \oplus A_{n - 1}$ et
définissons $d_{\diamond A, n}$ par la matrice
$$
d_{\diamond A, n}
=
\left(
\begin{matrix}
d_{A, n} & 0 & \text{id}_{A_{n - 1}} \\
0 & d_{A, n} & -\text{id}_{A_{n - 1}} \\
0 & 0 & -d_{A, n - 1}
\end{matrix}
\right) :
A_n \oplus A_n \oplus A_{n - 1} \to
A_{n - 1} \oplus A_{n - 1} \oplus A_{n - 2}
$$
Si $\mathcal{A}$ est la catégorie des groupes abéliens et si
$(x, y, z) \in A_n \oplus A_n \oplus A_{n - 1}$, alors
$d_{\diamond A, n}(x, y, z) = (d_n(x) + z, d_n(y) - z, -d_{n - 1}(z))$.
On vérifie aisément que $d^2 = 0$. Il existe manifestement deux applications
$\diamond a, \diamond b : A \to \diamond A$ (le premier et le second facteur)
et une application $\diamond A \to A[-1]$ qui donnent une suite exacte courte
$$
0 \to
A \oplus A \to
\diamond A \to
A[-1] \to
0
$$
scindée en chaque degré. De plus, il existe une suite d'applications
$\diamond h_n : A_n \to \diamond A_{n + 1}$, à savoir l'identité de $A_n$
vers le facteur $A_n$ de $\diamond A_{n + 1}$, telle que $\diamond h$ soit
une homotopie entre $\diamond a$ et $\diamond b$.

\medskip\noindent
Nous en concluons que tout morphisme $f : \diamond A \to B$ donne un triplet
$(a, b, h)$ en posant $a = f \circ \diamond a$,
$b = f \circ \diamond b$ et $h_n = f_{n + 1} \circ \diamond h_n$.
Réciproquement, un triplet $(a, b, h)$ fournit un morphisme
$f : \diamond A \to B$ en posant
$$
f_n = (a_n, b_n, h_{n - 1}).
$$
Pour vérifier qu'il s'agit d'un morphisme de complexes de chaînes, il faut
effectuer un calcul. Nous ne le faisons que lorsque $\mathcal{A}$ est la
catégorie des groupes abéliens. Soit
$(x, y, z) \in \diamond A_n = A_n \oplus A_n \oplus A_{n - 1}$. Alors
\begin{eqnarray*}
f_{n - 1}(d_n(x, y, z)) & = &
f_{n - 1}(d_n(x) + z, d_n(y) - z, -d_{n - 1}(z)) \\
& = &
a_{n - 1}(d_n(x)) + a_{n - 1}(z) + b_{n - 1}(d_n(y)) - b_{n - 1}(z) - h_{n - 2}(d_{n - 1}(z))
\end{eqnarray*}
et
\begin{eqnarray*}
d_n(f_n(x, y, z)) & = &
d_n(a_n(x) + b_n(y) + h_{n - 1}(z)) \\
& = &
d_n(a_n(x)) + d_n(b_n(y)) + d_n(h_{n - 1}(z))
\end{eqnarray*}
et ces expressions sont égales par définition d'une homotopie.
\end{proof}

\noindent
Notons que l'extension
$$
0 \to
A \oplus A \to
\diamond A \to
A[-1] \to
0
$$
est munie de sections des morphismes $\diamond A_n \to A[-1]_n$ telles que
le morphisme associé $\delta : A[-1] \to (A \oplus A)[-1]$, voir
Homologie, lemme \ref{homology-lemma-ses-termwise-split}, soit égal au
morphisme $(1, -1) : A[-1] \to A[-1] \oplus A[-1]$.

\begin{lemma}
\label{lemma-map-into-diamond}
Soit $\mathcal{A}$ une catégorie abélienne. Soit
$$
0 \to A \oplus A \to B \to C \to 0
$$
une suite exacte courte de complexes de chaînes de $\mathcal{A}$.
Supposons donnés en outre des morphismes $s_n : C_n \to B_n$ qui scindent
la suite exacte courte associée en degré $n$.
Soit $\delta(s) : C \to (A \oplus A)[-1] = A[-1] \oplus A[-1]$
le morphisme de complexes associé ; voir
Homologie, lemme \ref{homology-lemma-ses-termwise-split}.
Si $\delta(s)$ se factorise par le morphisme
$(1, -1) : A[-1] \to A[-1] \oplus A[-1]$, alors il existe un unique
morphisme $B \to \diamond A$ qui s'insère dans un diagramme commutatif
$$
\xymatrix{
0 \ar[r] &
A \oplus A \ar[d] \ar[r] &
B \ar[r] \ar[d] &
C \ar[d] \ar[r] &
0 \\
0 \ar[r] &
A \oplus A \ar[r] &
\diamond A \ar[r] &
A[-1] \ar[r] &
0
}
$$
où les applications verticales sont compatibles avec les scindages $s_n$
ainsi qu'avec les scindages de $\diamond A_n \to A[-1]_n$.
\end{lemma}

\begin{proof}
Notons $(p_n, q_n) : B_n \to A_n \oplus A_n$ le morphisme $\pi_n$ de
Homologie, lemme \ref{homology-lemma-ses-termwise-split}.
Notons aussi $(a, b) : A \oplus A \to B$ et $r : B \to C$ les applications
de la suite exacte courte. Écrivons la factorisation de $\delta(s)$ sous la
forme $\delta(s) = (1, -1) \circ f$. Cela signifie que
$p_{n - 1} \circ d_{B, n} \circ s_n = f_n$ et
$q_{n - 1} \circ d_{B, n} \circ s_n = - f_n$.
Prenons pour application $B_n \to \diamond A_n = A_n \oplus A_n \oplus A_{n - 1}$
l'application $(p_n, q_n, f_n \circ r_n)$.

\medskip\noindent
Il faut maintenant vérifier que ceci définit bien un morphisme de complexes.
Nous ne le faisons que dans le cas des groupes abéliens.
Choisissons $x \in B_n$. Alors $x = a_n(x_1) + b_n(x_2) + s_n(x_3)$,
et il suffit de vérifier séparément, pour chaque terme, que notre définition
commute avec la différentielle. Pour le terme
$a_n(x_1)$, on a $(p_n, q_n, f_n \circ r_n)(a_n(x_1)) =
(x_1, 0, 0)$, et le résultat est évident. Il en va de même pour
le terme $b_n(x_2)$. Pour le terme $s_n(x_3)$, on a
\begin{eqnarray*}
(p_{n - 1}, q_{n - 1}, f_{n - 1} \circ r_{n - 1})(d_n(s_n(x_3))) & = &
(p_{n - 1}, q_{n - 1}, f_{n - 1} \circ r_{n - 1})( \\
& & \ \ \ \ \ a_{n - 1}(f_n(x_3)) - b_{n - 1}(f_n(x_3)) + s_{n - 1}(d_n(x_3))) \\
& = &
(f_n(x_3), -f_n(x_3), f_{n - 1}(d_n(x_3)))
\end{eqnarray*}
par définition de $f_n$. D'autre part,
\begin{eqnarray*}
d_n(p_n, q_n, f_n \circ r_n)(s_n(x_3)) & = & d_n(0, 0, f_n(x_3)) \\
& = &
(f_n(x_3), - f_n(x_3), d_{A[-1], n}(f_n(x_3)))
\end{eqnarray*}
Le résultat en découle, puisque $f$ est un morphisme de complexes.
\end{proof}

\begin{lemma}
\label{lemma-backwards-homotopy}
Soit $\mathcal{A}$ une catégorie abélienne. Soient $U$, $V$ des objets
simpliciaux de $\mathcal{A}$ et soient $a, b : U \to V$ deux morphismes.
Supposons que les applications correspondantes de complexes de chaînes
$N(a), N(b) : N(U) \to N(V)$ soient homotopes par une homotopie
$\{N_n : N(U)_n \to N(V)_{n + 1}\}$.
Alors il existe une homotopie de $a$ vers $b$ au sens de la
définition \ref{definition-homotopy}. De plus, on peut choisir l'homotopie
$h : U \times \Delta[1] \to V$ de sorte que $N_n = N(h)_n$, où $N(h)$ est
l'homotopie issue de $h$ comme dans la section \ref{section-homotopy-abelian}.
\end{lemma}

\begin{proof}
Soit $(\diamond N(U), \diamond a, \diamond b, \diamond h)$ comme dans le
lemme \ref{lemma-represent-homotopy} et sa démonstration. D'après ce lemme,
il existe un morphisme $\diamond N(U) \to N(V)$ représentant le triplet
$(N(a), N(b), \{N_n\})$. Nous allons montrer qu'il existe un morphisme
$\psi : N(U \times \Delta[1]) \to \diamond{N(U)}$
tel que $\diamond a = \psi \circ N(e_0)$ et
$\diamond b = \psi \circ N(e_1)$. Nous montrerons de plus que l'homotopie entre
$N(e_0), N(e_1) :
N(U) \to N(U \times \Delta[1])$ provenant de
(\ref{equation-homotopy-to-homotopy}) et du
lemme \ref{lemma-homotopy-s-N}, avec
$h = \text{id}_{U \times \Delta[1]}$, est envoyée par
$\psi$ sur l'homotopie canonique $\diamond h$ entre les deux applications
$\diamond a, \diamond b : N(U) \to \diamond{N(U)}$.
Cela impliquera évidemment le lemme.

\medskip\noindent
Notons que le foncteur
$N : \text{Simp}(\mathcal{A}) \to \text{Ch}_{\geq 0}(\mathcal{A})$ est un
facteur direct du foncteur
$s : \text{Simp}(\mathcal{A}) \to \text{Ch}_{\geq 0}(\mathcal{A})$.
De plus, le foncteur $\diamond$ est compatible aux sommes directes.
Il suffit donc de construire un morphisme
$\Psi : s(U \times \Delta[1]) \to \diamond{s(U)}$ possédant les propriétés
correspondantes. C'est ce que nous faisons ci-dessous.

\medskip\noindent
D'après la définition \ref{definition-homotopy}, les morphismes
$e_0 : U \to U \times \Delta[1]$ et $e_1 : U \to U \times \Delta[1]$
sont homotopes par l'homotopie $\text{id}_{U \times \Delta[1]}$.
Le lemme \ref{lemma-homotopy-s-N} fournit une homotopie explicite
$\{h_n : s(U)_n \to s(U \times \Delta[1])_{n + 1}\}$
entre les morphismes de complexes de chaînes
$s(e_0) : s(U) \to s(U \times \Delta[1])$ et
$s(e_1) : s(U) \to s(U \times \Delta[1])$. D'après le
lemme \ref{lemma-map-into-diamond}, nous obtenons un morphisme correspondant
$$
\Phi : \diamond{s(U)} \to s(U \times \Delta[1])
$$
D'après la construction, la restriction de $\Phi_n$ au facteur direct
$s(U)[-1]_n = s(U)_{n - 1}$ de $\diamond{s(U)}_n$ est égale à
$h_{n - 1}$. Or
$$
h_{n - 1} = \sum\nolimits_{i = 0}^{n - 1}
(-1)^{i + 1} s^{n - 1}_i \cdot \alpha^n_{i + 1} :
U_{n - 1} \to \bigoplus\nolimits_j U_n \cdot \alpha^n_j.
$$
avec des notations évidentes.

\medskip\noindent
D'autre part, les morphismes $e_i : U \to U \times \Delta[1]$ induisent un
morphisme $(e_0, e_1) : U \oplus U \to U \times \Delta[1]$.
Notons $W$ son conoyau. Si nous écrivons
$(U \times \Delta[1])_n = \bigoplus_{\alpha : [n] \to [1]} U_n \cdot \alpha$,
nous pouvons alors identifier
$W_n = \bigoplus_{i = 1}^n U_n \cdot \alpha^n_i$, où
$\alpha^n_i$ est comme dans la section \ref{section-homotopy}.
Nous avons un diagramme commutatif
$$
\xymatrix{
0 \ar[r] &
U \oplus U \ar[rd]_{(1, 1)} \ar[r] &
U \times \Delta[1] \ar[d]^\pi \ar[r] &
W \ar[r] &
0 \\
& & U & &
}
$$
Après application du foncteur $s$, nous obtenons donc un diagramme
commutatif analogue. Choisissons ensuite les scindages
$\sigma_n : s(W)_n \to s(U \times \Delta[1])_n$ en envoyant le facteur
$U_n \cdot \alpha^n_i \subset W_n$, par $(-1, 1)$, dans les facteurs
$U_n \cdot \alpha^n_0 \oplus U_n \cdot \alpha^n_i
\subset (U \times \Delta[1])_n$. Notons que
$s(\pi)_n \circ \sigma_n = 0$. Il s'ensuit que
$(1, 1) \circ \delta(\sigma)_n = 0$. Par conséquent,
$\delta(\sigma)$ se factorise comme dans le
lemme \ref{lemma-map-into-diamond}. Ce lemme fournit un morphisme canonique
$\Psi : s(U \times \Delta[1]) \to \diamond{s(U)}$.

\medskip\noindent
Pour calculer $\Psi$, calculons d'abord le morphisme
$\delta(\sigma) : s(W) \to s(U)[-1] \oplus s(U)[-1]$.
D'après Homologie, lemme \ref{homology-lemma-ses-termwise-split}, et sa
démonstration, il faut pour cela calculer
$$
d_{s(U \times \Delta[1]), n} \circ \sigma_n
-
\sigma_{n - 1} \circ d_{s(W), n}
$$
et l'écrire comme un morphisme à valeurs dans
$U_{n - 1} \cdot \alpha^{n - 1}_0 \oplus U_{n - 1} \cdot \alpha^{n - 1}_n$.
Nous ne le faisons que lorsque $\mathcal{A}$ est la catégorie des groupes
abéliens. Nous employons la notation abrégée $x_{\alpha}$, pour $x \in U_n$,
afin de désigner l'élément $x$ du facteur $U_n \cdot \alpha$ de
$(U \times \Delta[1])_n$. Rappelons que
$$
d_{s(U \times \Delta[1]), n}
=
\sum\nolimits_{i = 0}^n (-1)^i d^n_i
$$
où $d^n_i$ envoie le facteur $U_n \cdot \alpha$ dans le facteur
$U_{n - 1} \cdot (\alpha \circ \delta^n_i)$ par le morphisme $d^n_i$
de l'objet simplicial $U$. Avec la notation ci-dessus, cela signifie que
$$
d_{s(U \times \Delta[1]), n}(x_\alpha) =
\sum\nolimits_{i = 0}^n (-1)^i (d^n_i(x))_{\alpha \circ \delta^n_i}
$$
Partant de $x_\alpha \in W_n$, autrement dit de $\alpha = \alpha^n_j$
pour un certain $j \in \{1, \ldots, n\}$, on a
$\sigma_n(x_\alpha) = x_\alpha - x_{\alpha^n_0}$, et par conséquent
$$
(d_{s(U \times \Delta[1]), n} \circ \sigma_n)(x_\alpha)
=
\sum\nolimits_{i = 0}^n (-1)^i (d^n_i(x))_{\alpha \circ \delta^n_i}
-
\sum\nolimits_{i = 0}^n (-1)^i (d^n_i(x))_{\alpha^n_0 \circ \delta^n_i}
$$
Pour calculer $d_{s(W), n}(x_\alpha)$, il faut omettre tous les termes pour
lesquels $\alpha \circ \delta^n_i = \alpha^{n - 1}_0, \alpha^{n - 1}_n$.
Nous obtenons donc
$$
\begin{matrix}
(\sigma_{n - 1} \circ d_{s(W), n})(x_\alpha) = \\
\sum\nolimits_{i = 0, \ldots, n\text{ et }
\alpha \circ \delta^n_i \not = \alpha^{n - 1}_0\text{ ou }\alpha^{n - 1}_n}
\Big((-1)^i (d^n_i(x))_{\alpha \circ \delta^n_i}
-
(-1)^i (d^n_i(x))_{\alpha^{n - 1}_0}
\Big)
\end{matrix}
$$
Manifestement, la différence des deux termes est la somme
$$
\sum\nolimits_{i = 0, \ldots, n\text{ et }
\alpha \circ \delta^n_i = \alpha^{n - 1}_0\text{ ou }\alpha^{n - 1}_n}
\Big((-1)^i (d^n_i(x))_{\alpha \circ \delta^n_i}
-
(-1)^i (d^n_i(x))_{\alpha^{n - 1}_0}
\Big)
$$
Bien entendu, si $\alpha \circ \delta^n_i = \alpha^{n - 1}_0$,
le terme disparaît. Rappelons que $\alpha = \alpha^n_j$
pour un certain $j \in \{1, \ldots, n\}$. La seule possibilité pour que
$\alpha^n_j \circ \delta^n_i = \alpha^{n - 1}_n$ est
$j = n$ et $i = n$. Nous obtenons donc $0$ sauf si $j = n$, auquel cas
nous obtenons
$(-1)^n(d^n_n(x))_{\alpha^{n - 1}_n} - (-1)^n(d^n_n(x))_{\alpha^{n - 1}_0}$.
Autrement dit, nous en concluons que le morphisme
$$
\delta(\sigma)_n :
W_n
\to
(s(U)[-1] \oplus s(U)[-1])_n = U_{n - 1} \oplus U_{n - 1}
$$
est nul sur tous les facteurs directs sauf $U_n \cdot \alpha^n_n$, et que,
sur ce facteur, il est égal à $((-1)^nd^n_n, -(-1)^nd^n_n)$.
(En effet, le premier des deux facteurs correspond au facteur portant
$\alpha^{n - 1}_n$, car c'est l'application $[n - 1] \to [1]$ qui envoie
tout élément sur $0$, et qui correspond donc à $e_0$.)

\medskip\noindent
Nous obtenons un diagramme canonique
$$
\xymatrix{
0 \ar[r] &
s(U) \oplus s(U) \ar[r] \ar[d] &
\diamond{s(U)} \ar[r] \ar[d]^{\Phi}&
s(U)[-1] \ar[r] \ar[d] &
0 \\
0 \ar[r] &
s(U) \oplus s(U) \ar[r] \ar[d] &
s(U \times \Delta[1]) \ar[r] \ar[d]^\Psi &
s(W) \ar[r] \ar[d] &
0 \\
0 \ar[r] &
s(U) \oplus s(U) \ar[r] &
\diamond{s(U)} \ar[r] &
s(U)[-1] \ar[r] &
0
}
$$
Nous affirmons que $\Psi \circ \Phi$ est l'identité.
Pour le voir, il suffit de montrer que la composée induite par $\Phi$ et
$\delta(\sigma)$,
$s(U)[-1] \to s(W) \to s(U)[-1] \oplus s(U)[-1]$, est l'identité sur le
premier facteur et l'opposé de l'identité sur le second.
D'après les calculs ci-dessus, elle vaut
$((-1)^nd^n_n, -(-1)^nd^n_n) \circ (-1)^n s^{n - 1}_{n - 1} = (1, -1)$,
comme souhaité.
\end{proof}









\section{Fibrations de Kan triviales}
\label{section-trivial-kan}

\noindent
Rappelons que, pour $n \geq 0$, l'ensemble simplicial $\Delta[n]$ est donné
par la règle $[k] \mapsto \Mor_\Delta([k], [n])$ ; voir
l'exemple \ref{example-simplex-simplicial-set}. Rappelons que $\Delta[n]$
possède un unique $n$-simplexe non dégénéré et que tous ses simplexes non
dégénérés sont des faces de ce $n$-simplexe. En fait, les simplexes non
dégénérés de $\Delta[n]$ correspondent exactement aux morphismes injectifs
$[k] \to [n]$, que nous pouvons identifier aux sous-ensembles de $[n]$.
Rappelons en outre que $\Mor(\Delta[n], X) = X_n$ pour tout ensemble
simplicial $X$ (lemme \ref{lemma-simplex-map}). Posons
$$
\partial \Delta[n] = i_{(n - 1)!}\text{sk}_{n - 1}\Delta[n]
$$
et appelons-le le {\it bord} de $\Delta[n]$.
Le lemme \ref{lemma-n-skeleton-sets} montre que
$\partial \Delta[n] \subset \Delta[n]$ est le sous-ensemble simplicial ayant
les mêmes simplexes non dégénérés aux degrés $\leq n - 1$, mais ne contenant
pas le $n$-simplexe non dégénéré.

\begin{definition}
\label{definition-trivial-kan}
Une application $X \to Y$ d'ensembles simpliciaux est appelée une
{\it fibration de Kan triviale} si $X_0 \to Y_0$ est surjective et si, pour
tout $n \geq 1$ et tout diagramme commutatif dont les flèches pleines sont
données,
$$
\xymatrix{
\partial \Delta[n] \ar[r] \ar[d] & X \ar[d] \\
\Delta[n] \ar[r] \ar@{-->}[ru] & Y
}
$$
il existe une flèche pointillée qui rend le diagramme commutatif.
\end{definition}

\noindent
Une fibration de Kan triviale satisfait une propriété de relèvement très générale.

\begin{lemma}
\label{lemma-trivial-kan}
Soit $f : X \to Y$ une fibration de Kan triviale d'ensembles simpliciaux.
Pour tout diagramme commutatif dont les flèches pleines sont données
$$
\xymatrix{
Z \ar[r]_b \ar[d] & X \ar[d] \\
W \ar[r]^a \ar@{-->}[ru] & Y
}
$$
d'ensembles simpliciaux, où $Z \to W$ est injective en chaque degré,
il existe une flèche pointillée qui rend le diagramme commutatif.
\end{lemma}

\begin{proof}
Supposons que $Z \not = W$. Soit $n$ le plus petit entier tel que
$Z_n \not = W_n$. Soit $x \in W_n$, $x \not \in Z_n$. Notons
$Z' \subset W$ le sous-ensemble simplicial contenant $Z$, $x$ et toutes les
dégénérescences de $x$. Soit $\varphi : \Delta[n] \to Z'$ le morphisme
correspondant à $x$ (lemme \ref{lemma-simplex-map}).
Alors $\varphi|_{\partial \Delta[n]}$ est à valeurs dans $Z$, puisque tous
les simplexes non dégénérés de $\partial \Delta[n]$ sont envoyés dans $Z$.
Par hypothèse, nous pouvons prolonger
$b \circ \varphi|_{\partial \Delta[n]}$ en
$\beta : \Delta[n] \to X$.
D'après le lemme \ref{lemma-glue-simplex}, l'ensemble simplicial $Z'$ est la
somme amalgamée de $\Delta[n]$ et $Z$ le long de $\partial \Delta[n]$.
Ainsi, $b$ et $\beta$ définissent un morphisme $b' : Z' \to X$.
Autrement dit, nous avons prolongé le morphisme $b$ à un sous-ensemble
simplicial plus grand de $W$.

\medskip\noindent
La démonstration s'achève par une application du lemme de Zorn (omise).
\end{proof}

\begin{lemma}
\label{lemma-trivial-kan-base-change}
Soit $f : X \to Y$ une fibration de Kan triviale d'ensembles simpliciaux.
Soit $Y' \to Y$ un morphisme d'ensembles simpliciaux.
Alors $X \times_Y Y' \to Y'$ est une fibration de Kan triviale.
\end{lemma}

\begin{proof}
Ceci résulte immédiatement des propriétés fonctorielles du produit fibré
(lemme \ref{lemma-fibre-product}) et des définitions.
\end{proof}

\begin{lemma}
\label{lemma-trivial-kan-composition}
La composée de deux fibrations de Kan triviales est une fibration de Kan triviale.
\end{lemma}

\begin{proof}
Omis.
\end{proof}

\begin{lemma}
\label{lemma-limit-trivial-kan}
Soit $\ldots \to U^2 \to U^1 \to U^0$ une suite de fibrations de Kan
triviales. Soit $U = \lim U^t$, défini par $U_n = \lim U_n^t$.
Alors $U \to U^0$ est une fibration de Kan triviale.
\end{lemma}

\begin{proof}
Omis. Indication : utiliser le fait que la limite projective d'une suite
dénombrable de surjections d'ensembles est non vide.
\end{proof}

\begin{lemma}
\label{lemma-product-trivial-kan}
\begin{slogan}
Les produits de fibrations de Kan triviales sont des fibrations de Kan triviales.
\end{slogan}
Soit $X_i \to Y_i$ une famille de fibrations de Kan triviales. Alors
$\prod X_i \to \prod Y_i$ est une fibration de Kan triviale.
\end{lemma}

\begin{proof}
Omis.
\end{proof}

\begin{lemma}
\label{lemma-filtered-colimit-trivial-kan}
Une colimite filtrante de fibrations de Kan triviales est une fibration de Kan triviale.
\end{lemma}

\begin{proof}
Omis. Indication : voir la description des colimites filtrantes d'ensembles dans
Catégories, section \ref{categories-section-directed-colimits}.
\end{proof}

\begin{lemma}
\label{lemma-trivial-kan-homotopy}
Soit $f : X \to Y$ une fibration de Kan triviale d'ensembles simpliciaux.
Alors $f$ est une équivalence d'homotopie.
\end{lemma}

\begin{proof}
D'après le lemme \ref{lemma-trivial-kan}, nous pouvons choisir un inverse à
droite $g : Y \to X$ de $f$. Considérons le diagramme
$$
\xymatrix{
\partial \Delta[1] \times X \ar[d] \ar[r] & X \ar[d] \\
\Delta[1] \times X \ar[r] \ar@{-->}[ru] & Y
}
$$
La flèche horizontale supérieure est donnée par $\text{id}_X$ et $g \circ f$,
où l'on utilise
$(\partial \Delta[1] \times X)_n = X_n \amalg X_n$ pour tout $n \geq 0$.
La flèche horizontale inférieure est donnée par l'application
$\Delta[1] \to \Delta[0]$ et par $f : X \to Y$. Le diagramme est commutatif,
car $f \circ g \circ f = f$. D'après le lemme \ref{lemma-trivial-kan},
nous pouvons construire la flèche pointillée, ce qui conclut.
\end{proof}






\section{Fibrations de Kan}
\label{section-kan}

\noindent
Soient $n$, $k$ des entiers tels que $0 \leq k \leq n$ et $1 \leq n$.
Soient $\sigma_0, \ldots, \sigma_n$
les $n + 1$ faces de l'unique $n$-simplexe non dégénéré $\sigma$
de $\Delta[n]$, c'est-à-dire $\sigma_i = d_i\sigma$. On note
$$
\Lambda_k[n] \subset \Delta[n]
$$
la {\it $k$-ième corne} du $n$-simplexe $\Delta[n]$. C'est le
sous-ensemble simplicial de $\Delta[n]$ engendré par
$\sigma_0, \ldots, \hat \sigma_k, \ldots, \sigma_n$. Autrement dit,
l'image de l'inclusion affichée contient tous les simplexes non dégénérés
de $\Delta[n]$, à l'exception de $\sigma$ et de $\sigma_k$.

\begin{definition}
\label{definition-kan}
Une application $X \to Y$ d'ensembles simpliciaux est appelée une
{\it fibration de Kan} si, pour tous $k, n$ tels que $1 \leq n$,
$0 \leq k \leq n$ et tout diagramme commutatif dont les flèches pleines
sont données
$$
\xymatrix{
\Lambda_k[n] \ar[r] \ar[d] & X \ar[d] \\
\Delta[n] \ar[r] \ar@{-->}[ru] & Y
}
$$
il existe une flèche pointillée qui rend le diagramme commutatif. Un
{\it complexe de Kan} est un ensemble simplicial $X$ tel que $X \to *$
soit une fibration de Kan, où $*$ est l'ensemble simplicial constant
associé à un singleton.
\end{definition}

\noindent
Remarquons que $\Lambda_k[n]$ est toujours non vide. Ainsi, un morphisme de
l'ensemble simplicial vide vers n'importe quel ensemble simplicial est
toujours une fibration de Kan. Il résulte du lemme
\ref{lemma-trivial-kan} qu'une fibration de Kan triviale est une fibration
de Kan.

\begin{lemma}
\label{lemma-kan-base-change}
Soit $f : X \to Y$ une fibration de Kan d'ensembles simpliciaux.
Soit $Y' \to Y$ un morphisme d'ensembles simpliciaux.
Alors $X \times_Y Y' \to Y'$ est une fibration de Kan.
\end{lemma}

\begin{proof}
Ceci résulte immédiatement des propriétés fonctorielles du produit fibré
(lemme \ref{lemma-fibre-product}) et des définitions.
\end{proof}

\begin{lemma}
\label{lemma-kan-composition}
La composée de deux fibrations de Kan est une fibration de Kan.
\end{lemma}

\begin{proof}
Omis.
\end{proof}

\begin{lemma}
\label{lemma-limit-kan}
Soit $\ldots \to U^2 \to U^1 \to U^0$ une suite de fibrations de Kan.
Soit $U = \lim U^t$, défini en posant $U_n = \lim U_n^t$.
Alors $U \to U^0$ est une fibration de Kan.
\end{lemma}

\begin{proof}
Omis. Indication : utiliser le fait que la limite projective d'une suite
dénombrable de surjections d'ensembles est non vide.
\end{proof}

\begin{lemma}
\label{lemma-product-kan}
Soit $X_i \to Y_i$ une famille de fibrations de Kan. Alors
$\prod X_i \to \prod Y_i$ est une fibration de Kan.
\end{lemma}

\begin{proof}
Omis.
\end{proof}

\noindent
Le lemme suivant est dû à J.C.\ Moore, voir \cite{Moore-Cartan}.

\begin{lemma}
\label{lemma-simplicial-group-kan}
Soit $X$ un groupe simplicial. Alors $X$ est un complexe de Kan.
\end{lemma}

\begin{proof}
La démonstration suivante n'est essentiellement qu'une traduction en français
de la démonstration de la référence mentionnée ci-dessus.
Avec la terminologie expliquée dans l'introduction de cette section,
supposons que $f : \Lambda_k[n] \to X$ soit un morphisme défini sur une corne. Posons
$x_i = f(\sigma_i) \in X_{n - 1}$ pour $i = 0, \ldots, \hat k, \ldots, n$.
Cela signifie que, pour $i < j$, on a $d_i x_j = d_{j - 1} x_i$
dès que $i, j \not = k$.
Il faut trouver un $x \in X_n$ tel que $x_i = d_ix$ pour
$i = 0, \ldots, \hat k, \ldots, n$.

\medskip\noindent
Montrons d'abord qu'il existe un $u \in X_n$ tel que $d_i u = x_i$
pour $i < k$. C'est trivial pour $k = 0$. Si $k > 0$, on définit par
récurrence un élément $u^r \in X_n$ tel que $d_i u^r = x_i$ pour
$0 \leq i \leq r$. On commence avec $u^0 = s_0x_0$. Si $r < k - 1$,
on pose
$$
y^r = s_{r + 1}((d_{r + 1}u^r)^{-1}x_{r + 1}),\quad
u^{r + 1} = u^r y^r.
$$
Un calcul facile montre que $d_iy^r = 1$ (élément neutre du groupe
$X_{n - 1}$) pour $i \leq r$ et que
$d_{r + 1}y^r = (d_{r + 1}u^r)^{-1}x_{r + 1}$. Il s'ensuit que
$d_iu^{r + 1} = x_i$ pour $i \leq r + 1$. Enfin, prenons $u = u^{k - 1}$
pour obtenir le $u$ promis.

\medskip\noindent
Montrons ensuite, par récurrence sur l'entier $r$, $0 \leq r \leq n - k$,
qu'il existe un $x^r \in X_n$ tel que
$$
d_i x^r = x_i\quad\text{pour }i < k\text{ et }i > n - r.
$$
On commence avec $x^0 = u$ pour $r = 0$. Ayant défini $x^r$ pour
$r \leq n - k - 1$, on pose
$$
z^r = s_{n - r - 1}((d_{n - r}x^r)^{-1}x_{n - r}),\quad x^{r + 1} = x^rz^r
$$
Un calcul simple, utilisant les relations données, montre que
$d_iz^r = 1$ pour $i < k$ et $i > n - r$, et que
$d_{n - r}(z^r) = (d_{n - r}x^r)^{-1}x_{n - r}$. Il s'ensuit que
$d_ix^{r + 1} = x_i$ pour $i < k$ et $i > n - r - 1$. Enfin, on
prend $x = x^{n - k}$, ce qui achève la démonstration.
\end{proof}

\begin{lemma}
\label{lemma-surjection-simplicial-abelian-groups-kan}
Soit $f : X \to Y$ un homomorphisme de groupes abéliens simpliciaux
qui est surjectif en chaque degré. Alors $f$ est une fibration de Kan
d'ensembles simpliciaux.
\end{lemma}

\begin{proof}
Considérons un diagramme commutatif dont les flèches pleines sont données
$$
\xymatrix{
\Lambda_k[n] \ar[r]_a \ar[d] & X \ar[d] \\
\Delta[n] \ar[r]^b \ar@{-->}[ru] & Y
}
$$
comme dans la définition \ref{definition-kan}. L'application $a$ correspond
à $x_0, \ldots, \hat x_k, \ldots, x_n \in X_{n - 1}$ satisfaisant
$d_i x_j = d_{j - 1} x_i$ pour $i < j$, $i, j \not = k$.
L'application $b$ correspond à un élément $y \in Y_n$ tel
que $d_iy = f(x_i)$ pour $i \not = k$. Il s'agit de produire un
$x \in X_n$ tel que $d_ix = x_i$ pour $i \not = k$ et $f(x) = y$.

\medskip\noindent
Puisque $f$ est surjectif en chaque degré, on peut trouver $x \in X_n$ tel que
$f(x) = y$. Remplaçons $y$ par $0 = y - f(x)$ et $x_i$ par $x_i - d_ix$
pour $i \not = k$. On voit alors que l'on peut supposer $y = 0$. En
particulier, $f(x_i) = 0$. Autrement dit, on peut remplacer $X$ par
$\Ker(f) \subset X$ et $Y$ par $0$. Dans ce cas, l'assertion devient le
lemme \ref{lemma-simplicial-group-kan}.
\end{proof}

\begin{lemma}
\label{lemma-qis-simplicial-abelian-groups}
Soit $f : X \to Y$ un homomorphisme de groupes abéliens simpliciaux
qui est surjectif en chaque degré et induit un quasi-isomorphisme entre les
complexes de chaînes associés. Alors $f$ est une fibration de Kan triviale
d'ensembles simpliciaux.
\end{lemma}

\begin{proof}
Considérons un diagramme commutatif dont les flèches pleines sont données
$$
\xymatrix{
\partial \Delta[n] \ar[r]_a \ar[d] & X \ar[d] \\
\Delta[n] \ar[r]^b \ar@{-->}[ru] & Y
}
$$
comme dans la définition \ref{definition-trivial-kan}. L'application $a$
correspond à $x_0, \ldots, x_n \in X_{n - 1}$ satisfaisant
$d_i x_j = d_{j - 1} x_i$ pour $i < j$. L'application $b$ correspond à un
élément $y \in Y_n$ tel que $d_iy = f(x_i)$. Il s'agit de produire un
$x \in X_n$ tel que $d_ix = x_i$ et $f(x) = y$.

\medskip\noindent
Puisque $f$ est surjectif en chaque degré, on peut trouver $x \in X_n$ tel que
$f(x) = y$. Remplaçons $y$ par $0 = y - f(x)$ et $x_i$ par $x_i - d_ix$.
On voit alors que l'on peut supposer $y = 0$. En particulier,
$f(x_i) = 0$. Autrement dit, on peut remplacer $X$ par
$\Ker(f) \subset X$ et $Y$ par $0$. Cela convient, car, d'après
Homologie, lemme \ref{homology-lemma-long-exact-sequence-chain},
l'homologie du complexe de chaînes associé à $\Ker(f)$
est nulle et donc $\Ker(f) \to 0$ induit un quasi-isomorphisme
entre les complexes de chaînes associés.

\medskip\noindent
Puisque $X$ est un complexe de Kan (lemme
\ref{lemma-simplicial-group-kan}), on peut trouver $x \in X_n$ tel que
$d_i x = x_i$ pour $i = 0, \ldots, n - 1$.
Après avoir remplacé $x_i$ par $x_i - d_ix$ pour $i = 0, \ldots, n$, on
peut supposer que $x_0 = x_1 = \ldots = x_{n - 1} = 0$.
Dans ce cas, on voit que $d_i x_n = 0$ pour $i = 0, \ldots, n - 1$.
Ainsi, $x_n \in N(X)_{n - 1}$ et appartient au noyau de la différentielle
$N(X)_{n - 1} \to N(X)_{n - 2}$. Ici, $N(X)$ est le complexe de chaînes
normalisé associé à $X$, voir la section \ref{section-complexes}.
Puisque $N(X)$ est quasi-isomorphe à $s(X)$
(lemme \ref{lemma-quasi-isomorphism})
et est donc acyclique, on trouve $x \in N(X_n)$ dont la différentielle est
$x_n$. Cet $x$ répond à la question posée par le lemme, et la démonstration
est achevée.
\end{proof}

\begin{lemma}
\label{lemma-homotopy-equivalence}
Soit $f : X \to Y$ un morphisme de groupes abéliens simpliciaux. Si $f$
est une équivalence d'homotopie d'ensembles simpliciaux, alors $f$ induit
un quasi-isomorphisme entre les complexes de chaînes associés.
\end{lemma}

\begin{proof}
Dans cette démonstration, on écrira $H_n(Z) = H_n(s(Z)) = H_n(N(Z))$
lorsque $Z$ est un groupe abélien simplicial, avec $s$ et $N$ comme dans la
section \ref{section-complexes}.
Notons $\mathbf{Z}[X]$ le groupe abélien libre sur $X$ considéré comme
ensemble simplicial, et de même pour $\mathbf{Z}[Y]$. Considérons le
diagramme commutatif
$$
\xymatrix{
\mathbf{Z}[X] \ar[r]_g \ar[d] &
\mathbf{Z}[Y] \ar[d] \\
X \ar[r]^f & Y
}
$$
de groupes abéliens simpliciaux. Puisque la formation du groupe abélien
libre sur un ensemble est fonctorielle, on voit que la flèche horizontale
est une équivalence d'homotopie de groupes abéliens simpliciaux, voir le
lemme \ref{lemma-functorial-homotopy}. D'après le lemme
\ref{lemma-homotopy-equivalence-s-N}, on voit que
$H_n(g) : H_n(\mathbf{Z}[X]) \to H_n(\mathbf{Z}[Y])$ est bijective pour tout
$n \geq 0$.

\medskip\noindent
Soit $\xi \in H_n(Y)$. Par définition de $N(Y)$, on peut représenter
$\xi$ par un élément $y \in N(Y_n)$ dont le bord est nul.
Cela signifie que $y \in Y_n$ avec
$d^n_0(y) = \ldots = d^n_{n - 1}(y) = 0$, puisque $y \in N(Y_n)$, et
$d^n_n(y) = 0$, puisque le bord de $y$ est nul. Notons $0_n \in Y_n$
l'élément nul. On voit alors que
$$
\tilde y = [y] - [0_n] \in (\mathbf{Z}[Y])_n
$$
est un élément tel que
$d^n_0(\tilde y) = \ldots = d^n_{n - 1}(\tilde y) = 0$
et $d^n_n(\tilde y) = 0$. Ainsi, $\tilde y$ appartient à
$N(\mathbf{Z}[Y])_n$ et est de bord $0$, c'est-à-dire que $\tilde y$
détermine une classe $\tilde \xi \in H_n(\mathbf{Z}[Y])$ dont l'image est
$\xi$. Puisque $H_n(\mathbf{Z}[X]) \to H_n(\mathbf{Z}[Y])$ est bijective,
on peut relever $\tilde \xi$ en une classe de $H_n(\mathbf{Z}[X])$.
En considérant le diagramme commutatif ci-dessus, on voit que $\xi$
appartient à l'image de $H_n(X) \to H_n(Y)$.

\medskip\noindent
Soit $\xi \in H_n(X)$ un élément dont l'image dans $H_n(Y)$ est nulle.
Exactement comme dans le paragraphe précédent, on peut représenter
$\xi$ par un élément $x \in N(X_n)$ dont le bord est nul, c'est-à-dire
$d^n_0(x) = \ldots = d^n_{n - 1}(x) = d^n_n(x) = 0$.
En particulier, on voit que $[x] - [0_n]$ est un élément de
$N(\mathbf{Z}[X])_n$ dont le bord est nul, et définit donc un relèvement
$\tilde \xi \in H_n(\mathbf{Z}[X])$ de $\xi$.
Le fait que l'image de $\xi$ dans $H_n(Y)$ soit nulle signifie qu'il existe
un $y \in N(Y_{n + 1})$ dont le bord est $f_n(x)$.
Cela signifie que $d^{n + 1}_0(y) = \ldots = d^{n + 1}_n(y) = 0$ et
$d^{n + 1}_{n + 1}(y) = f(x)$.
Mais cela signifie exactement que
$z = [y] - [0_{n + 1}]$ appartient à $N(\mathbf{Z}[Y])_{n + 1}$ et que
$$
g([x] - [0_n]) = [f(x)] - [0_n] = \text{bord de }z
$$
Ceci prouve que l'image de $\tilde \xi$ dans $H_n(\mathbf{Z}[Y])$ est nulle.
Comme $H_n(\mathbf{Z}[X]) \to H_n(\mathbf{Z}[Y])$ est bijective,
on conclut que $\tilde \xi = 0$, et donc que $\xi = 0$.
\end{proof}






\section{Une équivalence d'homotopie}
\label{section-homotopy-equivalence}

\noindent
Supposons que $A$, $B$ soient des ensembles et que $f : A \to B$
soit une application. Considérons l'application associée d'ensembles
simpliciaux
$$
\xymatrix{
\text{cosk}_0(A) \ar@{=}[r] &
\Big(\ldots A \times A \times A
\ar[d] \ar@<2ex>[r] \ar@<0ex>[r] \ar@<-2ex>[r] &
A \times A \ar[d] \ar@<1ex>[r] \ar@<-1ex>[r] \ar@<1ex>[l] \ar@<-1ex>[l] &
A \Big) \ar[d] \ar@<0ex>[l] \\
\text{cosk}_0(B) \ar@{=}[r] &
\Big( \ldots B \times B \times B \ar@<2ex>[r] \ar@<0ex>[r] \ar@<-2ex>[r] &
B \times B \ar@<1ex>[r] \ar@<-1ex>[r] \ar@<1ex>[l] \ar@<-1ex>[l] &
B \Big) \ar@<0ex>[l]
}
$$
Voir l'exemple \ref{example-cosk0}. Le cas $n = 0$ du lemme suivant affirme
que cette application d'ensembles simpliciaux est une fibration de Kan
triviale si $f$ est surjective.

\begin{lemma}
\label{lemma-section}
Soit $f : V \to U$ un morphisme d'ensembles simpliciaux. Soit $n \geq 0$
un entier. Supposons que
\begin{enumerate}
\item L'application $f_i : V_i \to U_i$ est une bijection pour $i < n$.
\item L'application $f_n : V_n \to U_n$ est une surjection.
\item Le morphisme canonique $U \to \text{cosk}_n \text{sk}_n U$
est un isomorphisme.
\item Le morphisme canonique $V \to \text{cosk}_n \text{sk}_n V$
est un isomorphisme.
\end{enumerate}
Alors $f$ est une fibration de Kan triviale.
\end{lemma}

\begin{proof}
Considérons un diagramme dont les flèches pleines sont données
$$
\xymatrix{
\partial \Delta[k] \ar[r] \ar[d] & V \ar[d] \\
\Delta[k] \ar[r] \ar@{-->}[ru] & U
}
$$
comme dans la définition \ref{definition-trivial-kan}. Soit $x \in U_k$
le $k$-simplexe correspondant à la flèche horizontale inférieure.
Si $k \leq n$, alors la flèche pointillée est celle qui correspond à un
relèvement $y \in V_k$ de $x$ ; le diagramme commute, car les autres
simplexes non dégénérés de $\Delta[k]$ sont de degrés $< k$, où $f$ est un
isomorphisme. Si $k > n$, alors, d'après les conditions (3) et (4), on a
(en utilisant l'adjonction entre les foncteurs squelette et cosquelette)
$$
\Mor(\Delta[k], U) =
\Mor(\text{sk}_n\Delta[k], \text{sk}_nU) =
\Mor(\text{sk}_n\partial\Delta[k], \text{sk}_nU) =
\Mor(\partial \Delta[k], U)
$$
et de même pour $V$, car
$\text{sk}_n\Delta[k] = \text{sk}_n\partial\Delta[k]$ pour $k > n$.
On obtient donc, dans ce cas également, une unique flèche pointillée qui
complète le diagramme.
\end{proof}

\noindent
Soient $A, B$ des ensembles. Soient $f^0, f^1 : A \to B$ des applications
d'ensembles. Considérons les applications induites
$f^0, f^1 : \text{cosk}_0(A) \to \text{cosk}_0(B)$, désignées abusivement
par les mêmes symboles. Le lemme suivant, pour $n = 0$, affirme que $f^0$
est homotope à $f^1$. En fait, il existe une homotopie
$h : \text{cosk}_0(A) \times \Delta[1] \to \text{cosk}_0(B)$
de $f^0$ à $f^1$, de composantes
\begin{eqnarray*}
h_m : A \times \ldots \times A \times \Mor_{\Delta}([m], [1])
& \longrightarrow &
B \times \ldots \times B, \\
(a_0, \ldots, a_m, \alpha) & \longmapsto &
(f^{\alpha(0)}(a_0), \ldots, f^{\alpha(m)}(a_m))
\end{eqnarray*}
Pour vérifier que cela convient, remarquons que, pour une application
$\varphi : [k] \to [m]$, les applications induites sont
$(a_0, \ldots, a_m) \mapsto (a_{\varphi(0)}, \ldots, a_{\varphi(k)})$
et $\alpha \mapsto \alpha \circ \varphi$. Ainsi, $h = (h_m)_{m \geq 0}$
est manifestement une application d'ensembles simpliciaux, comme souhaité.

\begin{lemma}
\label{lemma-homotopy}
Soient $f^0, f^1 : V \to U$ des applications d'ensembles simpliciaux.
Soit $n \geq 0$ un entier.
Supposons que
\begin{enumerate}
\item Les applications $f^j_i : V_i \to U_i$, $j = 0, 1$, sont égales pour $i < n$.
\item Le morphisme canonique $U \to \text{cosk}_n \text{sk}_n U$
est un isomorphisme.
\item Le morphisme canonique $V \to \text{cosk}_n \text{sk}_n V$
est un isomorphisme.
\end{enumerate}
Alors $f^0$ est homotope à $f^1$.
\end{lemma}

\begin{proof}[Première démonstration]
Soit $W$ l'ensemble simplicial $n$-tronqué tel que $W_i = U_i$ pour
$i < n$ et $W_n = U_n / \sim$, où $\sim$ est la relation d'équivalence
engendrée par $f^0(y) \sim f^1(y)$ pour $y \in V_n$. Ceci a un sens, car les
morphismes $U(\varphi) : U_n \to U_i$ correspondant à
$\varphi : [i] \to [n]$ pour $i < n$ se factorisent par l'application
quotient $U_n \to W_n$, puisque $f^0$ et $f^1$ sont des morphismes
d'ensembles simpliciaux et sont égaux aux degrés $< n$. Ensuite, on étend
$W$ en un ensemble simplicial en prenant $\text{cosk}_n W$. D'après le
lemme \ref{lemma-section}, le morphisme $g : U \to W$ est une fibration de
Kan triviale. Remarquons que $g \circ f^0 = g \circ f^1$ par construction,
et notons $f : V \to W$ ce morphisme. Considérons le diagramme
$$
\xymatrix{
\partial \Delta[1] \times V \ar[rr]_{f^0, f^1} \ar[d] & & U \ar[d] \\
\Delta[1] \times V \ar[rr]^f \ar@{-->}[rru] & & W
}
$$
D'après le lemme \ref{lemma-trivial-kan}, la flèche pointillée existe, ce
qui achève la démonstration.
\end{proof}

\begin{proof}[Deuxième démonstration]
Il faut construire un morphisme d'ensembles simpliciaux
$h : V \times \Delta[1] \to U$ qui redonne $f^i$ par composition avec $e_i$.
Le cas $n = 0$ a été traité avant le lemme.
On peut donc supposer que $n \geq 1$.
L'application $\Delta[1] \to \text{cosk}_1 \text{sk}_1 \Delta[1]$
est un isomorphisme, voir le lemme \ref{lemma-simplex-cosk}.
On voit donc que $\Delta[1] \to \text{cosk}_n \text{sk}_n \Delta[1]$
est un isomorphisme puisque $n \geq 1$, voir le
lemme \ref{lemma-cosk-up}. Par conséquent, $V \times \Delta[1] \to
\text{cosk}_n \text{sk}_n (V \times \Delta[1])$
est également un isomorphisme, voir le lemme \ref{lemma-cosk-product}.
Autrement dit, pour construire l'homotopie, il suffit de construire un
morphisme convenable d'ensembles simpliciaux $n$-tronqués
$h : \text{sk}_n V \times \text{sk}_n \Delta[1] \to \text{sk}_n U$.

\medskip\noindent
Pour $k = 0, \ldots, n - 1$, on définit $h_k$ par la
formule $h_k(v, \alpha) = f^0(v) = f^1(v)$.
L'application $h_n : V_n \times \Mor_{\Delta}([n], [1]) \to U_n$
est définie comme suit. Soient $v \in V_n$ et $\alpha : [n] \to [1]$ :
\begin{enumerate}
\item Si $\Im(\alpha) = \{0\}$, alors on pose $h_n(v, \alpha) = f^0(v)$.
\item Si $\Im(\alpha) = \{0, 1\}$, alors on pose $h_n(v, \alpha) = f^0(v)$.
\item Si $\Im(\alpha) = \{1\}$, alors on pose $h_n(v, \alpha) = f^1(v)$.
\end{enumerate}
Soit $\varphi : [k] \to [l]$ un morphisme de $\Delta_{\leq n}$.
Montrons que le diagramme
$$
\xymatrix{
V_{l} \times \Mor([l], [1]) \ar[r] \ar[d] & U_{l} \ar[d] \\
V_{k} \times \Mor([k], [1]) \ar[r] & U_{k}
}
$$
est commutatif.
Soient $v \in V_{l}$ et $\alpha : [l] \to [1]$.
La commutativité signifie que
$$
h_k(V(\varphi)(v), \alpha \circ \varphi)
=
U(\varphi)(h_l(v, \alpha)).
$$
Dans presque tous les cas, cette égalité est vraie, car
$h_k(V(\varphi)(v), \alpha \circ \varphi) = f^0(V(\varphi)(v))$
et $U(\varphi)(h_l(v, \alpha)) = U(\varphi)(f^0(v))$, compte tenu du fait
que $f^0$ est un morphisme d'ensembles simpliciaux.
Les seuls cas où ce raisonnement ne s'applique pas sont ceux où
soit (A) $\Im(\alpha) = \{1\}$ et $l = n$,
soit (B) $\Im(\alpha \circ \varphi) = \{1\}$ et $k = n$.
Remarquons en outre que l'on a nécessairement $f^0(v) = f^1(v)$
pour tout $n$-simplexe dégénéré de $V$.
On peut donc restreindre davantage les cas ci-dessus aux cas
(A) $\Im(\alpha) = \{1\}$, $l = n$ et $v$ non dégénéré, et
(B) $\Im(\alpha \circ \varphi) = \{1\}$, $k = n$ et
$V(\varphi)(v)$ non dégénéré.

\medskip\noindent
Dans le cas (A), on voit que $\Im(\alpha \circ \varphi) = \{1\}$ également.
On voit donc que non seulement $h_l(v, \alpha) = f^1(v)$, mais aussi que
$h_k(V(\varphi)(v), \alpha \circ \varphi) = f^1(V(\varphi)(v))$.
Ainsi, la relation est vraie parce que $f^1$ est un morphisme d'ensembles
simpliciaux.

\medskip\noindent
Dans le cas (B), on conclut que $l = k = n$ et que
$\varphi$ est bijective, car sinon $V(\varphi)(v)$
est dégénéré. Ainsi, $\varphi = \text{id}_{[n]}$, ce qui est un cas trivial.
\end{proof}

\begin{lemma}
\label{lemma-cosk-minus-one-equivalence}
Soient $A$, $B$ des ensembles et $f : A \to B$ une application. Considérons
l'ensemble simplicial $U$ dont les $n$-simplexes sont
$$
A \times_B A \times_B \ldots \times_B A\ (n + 1 \text{ facteurs)}.
$$
Voir l'exemple \ref{example-fibre-products-simplicial-object}.
Si $f$ est surjective, le morphisme $U \to B$,
où $B$ désigne l'ensemble simplicial constant de valeur $B$,
est une fibration de Kan triviale.
\end{lemma}

\begin{proof}
Remarquons que $U$ s'insère dans un carré cartésien
$$
\xymatrix{
U \ar[d] \ar[r] & \text{cosk}_0(A) \ar[d] \\
B \ar[r] & \text{cosk}_0(B)
}
$$
Puisque la flèche verticale de droite est une fibration de Kan triviale
d'après le lemme \ref{lemma-section}, il en est de même de celle de gauche
d'après le lemme \ref{lemma-trivial-kan-base-change}.
\end{proof}




\section{Préparation aux résolutions standard}
\label{section-standard-pre}

\noindent
Le contenu de cette section se trouve dans \cite[Appendice 1]{Godement}

\begin{example}
\label{example-godement}
Soit $Y : \mathcal{C} \to \mathcal{C}$ un foncteur d'une catégorie dans
elle-même et supposons données des transformations de foncteurs
$$
d : Y \longrightarrow \text{id}_\mathcal{C}
\quad\text{et}\quad
s : Y \longrightarrow Y \circ Y
$$
À l'aide de ces transformations, nous pouvons construire quelque chose qui
ressemble à un objet simplicial. Plus précisément, pour $n \geq 0$, posons
$$
X_n = Y \circ \ldots \circ Y \quad (n + 1\text{ compositions})
$$
Remarquons que $X_{n + m + 1} = X_n \circ X_m$ pour $n, m \geq 0$.
Ensuite, pour $n \geq 0$ et $0 \leq j \leq n$, nous définissons, avec les
notations de Catégories, section \ref{categories-section-formal-cat-cat},
$$
d^n_j = 1_{X_{j - 1}} \star d \star 1_{X_{n - j - 1}} :
X_n
\to
X_{n - 1}
\quad\text{et}\quad
s^n_j = 1_{X_{j - 1}} \star s \star 1_{X_{n - j - 1}} :
X_n
\to
X_{n + 1}
$$
Ainsi, $d^n_j$, resp.\ $s^n_j$, est la transformation naturelle qui applique
$d$, resp.\ $s$, au $j$-ième $Y$ (compté à partir de la gauche) dans la
composition définissant $X_n$.
\end{example}

\begin{lemma}
\label{lemma-godement}
Dans l'exemple \ref{example-godement}, si
$$
1_Y = (d \star 1_Y) \circ s = (1_Y \star d) \circ s
\quad\text{et}\quad
(s \star 1) \circ s = (1 \star s) \circ s
$$
alors $X = (X_n, d^n_j, s^n_j)$ est un objet simplicial de la catégorie
des endofoncteurs de $\mathcal{C}$ et $d : X_0 = Y \to \text{id}_\mathcal{C}$
définit une augmentation.
\end{lemma}

\begin{proof}
Pour voir que nous obtenons un objet simplicial, nous devons vérifier que les
relations (1)(a) -- (e) du lemme \ref{lemma-characterize-simplicial-object}
sont satisfaites. Nous utiliserons la notation abrégée
$$
1_a = 1_{X_{a - 1}} = 1_Y \star \ldots \star 1_Y \quad (a\text{ facteurs})
$$
pour $a \geq 0$. Avec cette notation, nous avons
$$
d^n_j = 1_j \star d \star 1_{n - j}
\quad\text{et}\quad
s^n_j = 1_j \star s \star 1_{n - j}
$$
Nous utiliserons à plusieurs reprises la règle selon laquelle, pour des
transformations de foncteurs $a, a', b, b'$, nous avons
$(a' \circ a) \star (b' \circ b) = (a' \star b') \circ (a \star b)$
pourvu que les compositions $\star$ et $\circ$ de cette formule soient
définies ; voir
Catégories, lemme \ref{categories-lemma-properties-2-cat-cats}.

\medskip\noindent
La condition (1)(a) est toujours satisfaite (aucune condition n'est requise
sur $d$ et $s$). En effet, soit $0 \leq i < j \leq n + 1$. Nous devons
montrer que $d^n_i \circ d^{n + 1}_j = d^n_{j - 1} \circ d^{n + 1}_i$,
c'est-à-dire
$$
(1_i \star d \star 1_{n - i}) \circ
(1_j \star d \star 1_{n + 1 - j}) =
(1_{j - 1} \star d \star 1_{n + 1 - j}) \circ
(1_i \star d \star 1_{n + 1 - i})
$$
Nous pouvons réécrire le membre de gauche sous la forme
\begin{align*}
& (1_i \star d \star 1_{j - i - 1} \star 1_{n + 1 - j}) \circ
(1_i \star 1_1 \star 1_{j - i - 1} \star d \star 1_{n + 1 - j}) \\
& =
1_i \star
\left((d \star 1_{j - i - 1}) \circ
(1_1 \star 1_{j - i - 1} \star d)\right) \star 1_{n + 1 - j} \\
& =
1_i \star d \star 1_{j - i - 1} \star d \star 1_{n + 1 - j} 
\end{align*}
La seconde égalité est vraie parce que $d \circ 1_1 = d$ et
$1_{j - i} \circ (1_{j - i - 1} \star d) = 1_{j - i - 1} \star d$.
Un calcul analogue donne le même résultat pour le membre de droite.

\medskip\noindent
Vérifions la condition (1)(b). Soit $0 \leq i < j \leq n - 1$. Nous devons
montrer que $d^n_i \circ s^{n - 1}_j = s^{n - 2}_{j - 1} \circ d^{n - 1}_i$,
c'est-à-dire
$$
(1_i \star d \star 1_{n - i}) \circ
(1_j \star s \star 1_{n - 1 - j}) =
(1_{j - 1} \star s \star 1_{n - 1 - j}) \circ
(1_i \star d \star 1_{n - 1 - i})
$$
Par le même type de calcul que dans le cas (1)(a), les deux membres se
simplifient en $1_i \star d \star 1_{j - i - 1} \star s \star 1_{n - j - 1}$.

\medskip\noindent
Vérifions la condition (1)(c). Soit $0 \leq j \leq n - 1$. Nous devons montrer
$\text{id} = d^n_j \circ s^{n - 1}_j = d^n_{j + 1} \circ s^{n - 1}_j$,
c'est-à-dire
$$
1_n =
(1_j \star d \star 1_{n - j}) \circ
(1_j \star s \star 1_{n - 1 - j}) =
(1_{j + 1} \star d \star 1_{n - j - 1}) \circ
(1_j \star s \star 1_{n - 1 - j})
$$
Cela résulte immédiatement de la première hypothèse du lemme.

\medskip\noindent
Vérifions la condition (1)(d). Soit $0 < j + 1 < i \leq n$. Nous devons
montrer $d^n_i \circ s^{n - 1}_j = s^{n - 2}_j \circ d^{n - 1}_{i - 1}$,
c'est-à-dire
$$
(1_i \star d \star 1_{n - i}) \circ
(1_j \star s \star 1_{n - 1 - j}) =
(1_j \star s \star 1_{n - 2 - j}) \circ
(1_{i - 1} \star d \star 1_{n - i})
$$
Par le même type de calcul que dans le cas (1)(a), les deux membres se
simplifient en $1_j \star s \star 1_{i - j - 2} \star d \star 1_{n - i}$.

\medskip\noindent
Vérifions la condition (1)(e). Soit $0 \leq i \leq j \leq n - 1$. Nous
devons montrer que
$s^n_i \circ s^{n - 1}_j = s^n_{j + 1} \circ s^{n - 1}_i$, c'est-à-dire
$$
(1_i \star s \star 1_{n - i}) \circ
(1_j \star s \star 1_{n - 1 - j}) =
(1_{j + 1} \star s \star 1_{n - 1 - j}) \circ
(1_i \star s \star 1_{n - 1 - i})
$$
Par le même type de calcul que dans le cas (1)(a), cela se réduit à
$$
(s \star 1_{j - i + 1}) \circ (1_{j - i} \star s) =
(1_{j - i + 1} \star s) \circ (s \star 1_{j - i})
$$
Si $j = i$, c'est exactement l'une des deux hypothèses du lemme.
Si $j > i$, les membres de gauche et de droite se réduisent tous deux à
$s \star 1_{j - i - 1} \star s$ par des calculs semblables à ceux que nous
avons effectués dans le cas (1)(a).

\medskip\noindent
Enfin, pour montrer que $d$ définit une augmentation, nous devons montrer que
$d \circ (1_1 \star d) = d \circ (d \star 1_1)$, ce qui est vrai puisque les
deux membres sont égaux à $d \star d$.
\end{proof}

\begin{example}
\label{example-godement-functorial}
Soient $\mathcal{C}$, $Y$, $d$, $s$ comme dans l'exemple
\ref{example-godement}, satisfaisant les équations du lemme
\ref{lemma-godement}. Étant donnés des foncteurs
$F : \mathcal{A} \to \mathcal{C}$ et $G : \mathcal{C} \to \mathcal{B}$,
nous obtenons un objet simplicial $G \circ X \circ F$ dans la catégorie des
foncteurs de $\mathcal{A}$ dans $\mathcal{B}$, muni d'une augmentation vers
$G \circ F$.
\end{example}

\begin{lemma}
\label{lemma-godement-section}
Soient $\mathcal{A}$, $\mathcal{B}$, $\mathcal{C}$, $Y$, $d$, $s$, $F$, $G$
comme dans l'exemple \ref{example-godement-functorial}. Étant donnée une
transformation de foncteurs $h_0 : G \circ F \to G \circ Y \circ F$ telle que
$$
1_{G \circ F} = (1_G \star d \star 1_F) \circ h_0
$$
alors il existe un morphisme $h : G \circ F \to G \circ X \circ F$
d'objets simpliciaux tel que $\epsilon \circ h = \text{id}$, où
$\epsilon : G \circ X \circ F \to G \circ F$ est l'augmentation.
\end{lemma}

\begin{proof}
Notons $u_n : Y = X_0 \to X_n$ le morphisme de l'objet simplicial $X$
correspondant à l'unique morphisme $[n] \to [0]$ de $\Delta$.
Posons $h_n : G \circ F \to G \circ X_n \circ F$ égal à
$(1_G \star u_n \star 1_F) \circ h_0$.

\medskip\noindent
Pour tout objet simplicial $X = (X_n)$ d'une catégorie quelconque,
$u =(u_n) : X_0 \to X$ est un morphisme de l'objet simplicial constant de
valeur $X_0$ vers $X$. Par conséquent, $h$ est un morphisme d'objets
simpliciaux, car il est la composée de $1_G \star u \star 1_F$ et de $h_0$.

\medskip\noindent
Vérifions que $\epsilon \circ h = \text{id}$. Nous calculons
$$
\epsilon_n \circ (1_G \star u_n \star 1_F) \circ h_0 =
\epsilon_0 \circ h_0 = \text{id}
$$
La première égalité vient de ce que $\epsilon$ est un morphisme d'objets
simpliciaux, et la seconde de ce que
$\epsilon_0 = (1_G \star d \star 1_F)$ et que nous pouvons appliquer
l'hypothèse de l'énoncé du lemme.
\end{proof}

\begin{lemma}
\label{lemma-godement-two-maps}
Soient $\mathcal{A}$, $\mathcal{B}$, $\mathcal{C}$, $Y$, $d$, $s$, $F$, $G$
comme dans l'exemple \ref{example-godement-functorial}. Soient
$F' : \mathcal{A} \to \mathcal{C}$ et
$G' : \mathcal{C} \to \mathcal{B}$ deux foncteurs.
Soit $(a_n) : G \circ X \to G' \circ X$ un morphisme d'objets simpliciaux
compatible avec $a : G \to G'$ via les augmentations.
Soit $(b_n) : X \circ F \to X \circ F'$ un morphisme d'objets simpliciaux
compatible avec $b : F \to F'$ via les augmentations. Alors les deux morphismes
$$
a \star (b_n), (a_n) \star b : G \circ X \circ F \to G' \circ X \circ F'
$$
sont homotopes.
\end{lemma}

\begin{proof}
Pour montrer que les morphismes sont homotopes, nous construisons des morphismes
$$
h_{n, i} : G \circ X_n \circ F \to G' \circ X_n \circ F'
$$
pour $n \geq 0$ et $0 \leq i \leq n + 1$, satisfaisant les relations décrites
dans le lemme \ref{lemma-relations-homotopy}. Voir aussi la remarque
\ref{remark-homotopy-better}. Pour satisfaire la condition (1) du lemme
\ref{lemma-relations-homotopy}, nous sommes contraints de poser
$h_{n, 0} = a \star b_n$ et $h_{n , n + 1} = a_n \star b$.
Un choix naturel est donc
$$
h_{n , i} = a_{i - 1} \star b_{n - i}
$$
pour $1 \leq i \leq n$. En posant $a = a_{-1}$ et $b = b_{-1}$, nous voyons
que la formule affichée est valable pour $0 \leq i \leq n + 1$.

\medskip\noindent
Rappelons que
$$
d^n_j = 1_G \star 1_j \star d \star 1_{n - j} \star 1_F
$$
sur $G \circ X \circ F$, où nous utilisons la notation
$1_a = 1_{Y \circ \ldots \circ Y}$ introduite dans la démonstration du lemme
\ref{lemma-godement}. Nous utiliserons ci-dessous les réécritures suivantes :
\begin{align*}
d^n_j
& =
d^j_j \star 1_{n - j} =
d^{j + 1}_j \star 1_{n - j - 1} = \ldots = d^{n - 1}_j \star 1_1 \\
& =
1_j \star d^{n - j}_0 = 1_{j - 1} \star d^{n - j + 1}_1 = \ldots =
1_1 \star d^{n - 1}_{j - 1}
\end{align*}
Bien entendu, nous avons les formules analogues pour $d^n_j$ sur
$G' \circ X \circ F'$.

\medskip\noindent
Vérifions la condition (2) du lemme \ref{lemma-relations-homotopy}.
Soit $i > j$. Nous devons montrer
$$
d^n_j \circ (a_{i - 1} \star b_{n - i}) =
(a_{i - 2} \star b_{n - i}) \circ d^n_j
$$
Puisque $i - 1 \geq j$, nous pouvons utiliser l'une des descriptions possibles
de $d^n_j$ pour réécrire le membre de gauche sous la forme
$$
(d^{i - 1}_j \star 1_{n - i + 1}) \circ (a_{i - 1} \star b_{n - i}) =
(d^{i - 1}_j \circ a_{i - 1}) \star b_{n - i} =
(a_{i - 2} \circ d^{i - 1}_j) \star b_{n - i}
$$
De même, le membre de droite devient
$$
(a_{i - 2} \star b_{n - i}) \circ (d^{i - 1}_j \star 1_{n - i + 1}) =
(a_{i - 2} \circ d^{i - 1}_j) \star b_{n - i}
$$
Nous obtenons donc le même résultat, ce qui vérifie (2).

\medskip\noindent
Vérifions la condition (3) du lemme \ref{lemma-relations-homotopy}.
Soit $i \leq j$. Nous devons montrer
$$
d^n_j \circ (a_{i - 1} \star b_{n - i}) =
(a_{i - 1} \star b_{n - 1 - i}) \circ d^n_j
$$
Puisque $j \geq i$, nous pouvons réécrire le membre de gauche sous la forme
$$
(1_i \star d^{n - i}_{j - i}) \circ (a_{i - 1} \star b_{n - i}) =
a_{i - 1} \star (b_{n - 1 - i} \circ d^{n - i}_{j - i})
$$
Une manipulation analogue montre que cela coïncide avec le membre de droite.

\medskip\noindent
Rappelons que
$$
s^n_j = 1_G \star 1_j \star s \star 1_{n - j} \star 1_F
$$
sur $G \circ X \circ F$. Nous utiliserons ci-dessous les réécritures suivantes :
\begin{align*}
s^n_j
& =
s^j_j \star 1_{n - j} = s^{j + 1}_j \star 1_{n - j - 1} = \ldots =
s^{n - 1}_j \star 1_1 \\
& =
1_j \star s^{n - j}_0 = 1_{j - 1} \star s^{n - j + 1}_1 = \ldots =
1_1 \star s^{n - 1}_{j - 1}
\end{align*}
Bien entendu, nous avons les formules analogues pour $s^n_j$ sur
$G' \circ X \circ F'$.

\medskip\noindent
Vérifions la condition (4) du lemme \ref{lemma-relations-homotopy}.
Soit $i > j$. Nous devons montrer
$$
s^n_j \circ (a_{i - 1} \star b_{n - i}) =
(a_i \star b_{n - i}) \circ s^n_j
$$
Puisque $i - 1 \geq j$, nous pouvons réécrire le membre de gauche sous la forme
$$
(s^{i - 1}_j \star 1_{n - i + 1}) \circ (a_{i - 1} \star b_{n - i}) =
(s^{i - 1}_j \circ a_{i - 1}) \star b_{n - i} =
(a_i \circ s^{i - 1}_j) \star b_{n - i}
$$
De même, le membre de droite devient
$$
(a_i \star b_{n - i}) \circ (s^{i - 1}_j \star 1_{n - i + 1}) =
(a_i \circ s^{i - 1}_j) \star b_{n - i}
$$
comme voulu.

\medskip\noindent
Vérifions la condition (5) du lemme \ref{lemma-relations-homotopy}.
Soit $i \leq j$. Nous devons montrer
$$
s^n_j \circ (a_{i - 1} \star b_{n - i}) =
(a_{i - 1} \star b_{n + 1 - i}) \circ s^n_j
$$
Cette égalité est vraie, car les deux membres sont égaux à
$a_{i - 1} \star (s^{n - i}_{j - i} \circ b_{n - i}) =
a_{i - 1} \star (b_{n + 1 - i} \circ s^{n - i}_{j - i})$
par les mêmes arguments que ci-dessus.
\end{proof}

\begin{lemma}
\label{lemma-godement-before-after}
Soient $\mathcal{C}$, $Y$, $d$, $s$ comme dans l'exemple
\ref{example-godement}, satisfaisant les équations du lemme
\ref{lemma-godement}. Soit
$f : \text{id}_\mathcal{C} \to \text{id}_\mathcal{C}$ un endomorphisme du
foncteur identité. Alors $f \star 1_X, 1_X \star f : X \to X$ sont des
morphismes d'objets simpliciaux compatibles avec $f$ via l'augmentation
$\epsilon : X \to \text{id}_\mathcal{C}$. De plus, $f \star 1_X$ et
$1_X \star f$ sont homotopes.
\end{lemma}

\begin{proof}
Le morphisme $f \star 1_X$ est celui dont les composantes sont
$$
X_n = \text{id}_\mathcal{C} \circ X_n
\xrightarrow{f \star 1_{X_n}}
\text{id}_\mathcal{C} \circ X_n = X_n
$$
Pour une transformation $a : F \to G$ d'endofoncteurs de $\mathcal{C}$, nous
avons $a \circ (f \star 1_F) = f \star a = (f \star 1_G) \circ a$. Ainsi,
$f \star 1_X$ est bien un morphisme d'objets simpliciaux. Il en va de même
pour $1_X \star f$.

\medskip\noindent
Pour montrer que les morphismes sont homotopes, nous construisons des
morphismes $h_{n, i} : X_n \to X_n$ pour $n \geq 0$ et $0 \leq i \leq n + 1$
satisfaisant les relations décrites dans le lemme
\ref{lemma-relations-homotopy}. Voir aussi la remarque
\ref{remark-homotopy-better}. On constate que nous pouvons prendre
$$
h_{n, i} = 1_i \star f \star 1_{n + 1 - i}
$$
où $1_i$ est la transformation identité de $Y \circ \ldots \circ Y$, comme
dans la démonstration du lemme \ref{lemma-godement}. Nous avons
$h_{n, 0} = f \star 1_{X_n}$ et $h_{n, n + 1} = 1_{X_n} \star f$
ce qui vérifie la première condition. Pour vérifier les autres conditions,
nous utilisons les observations faites dans la démonstration du lemme
\ref{lemma-godement-two-maps} au sujet des morphismes $d^n_j$ et $s^n_j$.

\medskip\noindent
Vérifions la condition (2) du lemme \ref{lemma-relations-homotopy}.
Soit $i > j$. Nous devons montrer
$$
d^n_j \circ (1_i \star f \star 1_{n + 1 - i}) =
(1_{i - 1} \star f \star 1_{n + 1 - i}) \circ d^n_j
$$
Puisque $i - 1 \geq j$, nous pouvons utiliser l'une des descriptions possibles
de $d^n_j$ pour réécrire le membre de gauche sous la forme
$$
(d^{i - 1}_j \star 1_{n - i + 1}) \circ (1_i \star f \star 1_{n + 1 - i}) =
d^{i - 1}_j \star f \star 1_{n + 1 - i}
$$
De même, le membre de droite devient
$$
(1_{i - 1} \star f \star 1_{n + 1 - i}) \circ
(d^{i - 1}_j \star 1_{n - i + 1}) =
d^{i - 1}_j \star f \star 1_{n + 1 - i}
$$
Nous obtenons donc le même résultat, ce qui vérifie (2).

\medskip\noindent
Les conditions (3), (4) et (5) du lemme \ref{lemma-relations-homotopy}
se vérifient exactement de la même manière, en suivant la stratégie de la
démonstration du lemme \ref{lemma-godement-two-maps}. Nous omettons les
détails\footnote{Lorsque $f$ est inversible, il suffit de démontrer que
$(a_n) = 1_X$ et $(b_n) = f^{-1} \star 1_X \star f$ sont homotopes.
Mais cela résulte du lemme \ref{lemma-godement-two-maps}, car dans ce cas
$a = b = 1_{\text{id}_\mathcal{C}}$.}.
\end{proof}
\section{Résolutions standard}
\label{section-standard}

\noindent
Une partie du contenu de cette section se trouve dans
\cite[Appendice 1]{Godement} et \cite[I 1.5]{cotangent}.

\begin{situation}
\label{situation-adjoint-functors}
Soient $\mathcal{A}$, $\mathcal{S}$ des catégories et soit
$V : \mathcal{A} \to \mathcal{S}$ un foncteur admettant pour adjoint à gauche
$U : \mathcal{S} \to \mathcal{A}$.
\end{situation}

\noindent
Dans cette situation très générale, nous allons construire un objet
simplicial $X$ dans la catégorie des foncteurs de $\mathcal{A}$ vers $\mathcal{A}$.
Nous conseillons de consulter les exemples présentés plus loin avant de lire
le texte de cette section.

\medskip\noindent
Pour cette construction, nous utiliserons la composition horizontale telle
qu'elle est définie dans Catégories, section \ref{categories-section-formal-cat-cat}.
La définition des morphismes d'adjonction\footnote{Nous ne pouvons pas utiliser
$\epsilon$ pour la co-unité de l'adjonction
car nous voulons utiliser $\epsilon$ pour
l'augmentation de notre objet simplicial.}
$$
d : U \circ V \to \text{id}_\mathcal{A} \quad (\text{co-unité})
\quad\text{et}\quad
\eta : \text{id}_\mathcal{S} \to V \circ U \quad (\text{unité})
$$
dans Catégories, section \ref{categories-section-adjoint}
montre que les composés
\begin{equation}
\label{equation-composition}
V \xrightarrow{\eta \star 1_V} V \circ U \circ V \xrightarrow{1_V \star d} V
\quad\text{et}\quad
U \xrightarrow{1_U \star \eta} U \circ V \circ U \xrightarrow{d \star 1_U} U
\end{equation}
sont les morphismes identité. Ici, pour définir le morphisme $\eta \star 1_V$,
nous identifions tacitement $V$ à $\text{id}_\mathcal{S} \circ V$ et
$1_V$ désigne $\text{id}_V : V \to V$. Nous utiliserons cette notation et
ces relations à plusieurs reprises dans la suite.
Pour $n \geq 0$, nous posons
$$
X_n = (U \circ V)^{\circ (n + 1)} =
U \circ V \circ U \circ \ldots \circ U \circ V
$$
Autrement dit, $X_n$ est le composé de $(n + 1)$ copies de $U \circ V$.
Nous posons aussi $X_{-1} = \text{id}_\mathcal{A}$.
Nous avons $X_{n + m + 1} = X_n \circ X_m$ pour tous $n, m \geq -1$.
Nous allons munir cette suite de foncteurs d'une structure d'objet simplicial
de $\text{Fun}(\mathcal{A}, \mathcal{A})$ en construisant les morphismes de
foncteurs
$$
d^n_j : X_n \to X_{n - 1},\quad s^n_j : X_n \to X_{n + 1}
$$
qui satisfont aux relations du
lemme \ref{lemma-relations-face-degeneracy}.
Plus précisément, nous posons
$$
d^n_j = 1_{X_{j - 1}} \star d \star 1_{X_{n - j - 1}}
\quad\text{et}\quad
s^n_j = 1_{X_{j - 1} \circ U} \star \eta \star 1_{V \circ X_{n - j - 1}}
$$
Enfin, posons $\epsilon_0 = d : X_0 \to X_{-1}$.

\begin{lemma}
\label{lemma-standard-simplicial}
Dans la situation \ref{situation-adjoint-functors},
le système $X = (X_n, d^n_j, s^n_j)$
est un objet simplicial de $\text{Fun}(\mathcal{A}, \mathcal{A})$
et $\epsilon_0$ définit une augmentation $\epsilon$ de $X$
vers l'objet simplicial constant de valeur $X_{-1} = \text{id}_\mathcal{A}$.
\end{lemma}

\begin{proof}
Considérons $Y = U \circ V : \mathcal{A} \to \mathcal{A}$.
Nous disposons déjà de la transformation
$d : Y = U \circ V \to \text{id}_\mathcal{A}$. Notons
$$
s = 1_U \star \eta \star 1_V :
Y = U \circ \text{id}_\mathcal{S} \circ V
\longrightarrow
U \circ V \circ U \circ V = Y \circ Y
$$
Nous sommes alors dans la situation de l'exemple \ref{example-godement}.
Il résulte immédiatement des formules que
les $X, d^n_i, s^n_i$ construits ci-dessus et les
$X, d^n_i, s^n_i$ construits à partir de $Y, d, s$
dans l'exemple \ref{example-godement} coïncident. Ainsi, d'après le
lemme \ref{lemma-godement}, il suffit de montrer que
$$
1_Y = (d \star 1_Y) \circ s = (1_Y \star d) \circ s
\quad\text{et}\quad
(s \star 1) \circ s = (1 \star s) \circ s
$$
La première égalité équivaut à l'égalité
$$
1_U \star 1_V = (d \star 1_U \star 1_V) \circ (1_U \star \eta \star 1_V)
$$
qui est vérifiée si $1_U = (d \star 1_U) \circ (1_U \star \eta)$,
ce qui résulte de (\ref{equation-composition}). Il en va de même pour la
seconde égalité. Pour la dernière équation, il faut montrer
$$
(1_U \star \eta \star 1_V \star 1_U \star 1_V) \circ
(1_U \star \eta \star 1_V) =
(1_U \star 1_V \star 1_U \star \eta \star 1_V) \circ
(1_U \star \eta \star 1_V)
$$
Pour cela, il suffit de montrer
$(\eta \star 1_V \star 1_U) \circ \eta = (1_V \star 1_U \star \eta) \circ \eta$
ce qui est vrai car les deux membres sont égaux à $\eta \star \eta$.
\end{proof}

\noindent
Avant de lire la démonstration du lemme suivant, nous conseillons au
lecteur de consulter l'exemple \ref{example-polynomial-algebra-homotopy}
présenté plus loin
afin de comprendre le but du lemme.

\begin{lemma}
\label{lemma-standard-simplicial-homotopy}
Dans la situation \ref{situation-adjoint-functors}, les applications
$$
1_V \star \epsilon : V \circ X \to V,
\quad\text{et}\quad
\epsilon \star 1_U : X \circ U \to U
$$
sont des équivalences d'homotopie.
\end{lemma}

\begin{proof}
Comme dans la démonstration du lemme \ref{lemma-standard-simplicial},
nous posons $Y = U \circ V$, de sorte que nous sommes
dans la situation de l'exemple \ref{example-godement}.

\medskip\noindent
Démonstration de la première équivalence d'homotopie. D'après le lemme
\ref{lemma-godement-section}, pour construire une application
$h : V \to V \circ X$ inverse à droite de $1_V \star \epsilon$, il
suffit de construire une application $h_0 : V \to V \circ Y = V \circ U \circ V$
telle que $1_V = (1_V \star d) \circ h_0$. Bien entendu, nous prenons
$h_0 = \eta \star 1_V$ et l'égalité résulte de (\ref{equation-composition}).
Pour achever la démonstration, il faut montrer que les deux applications
$$
(1_V \star \epsilon) \circ h, 1_V \star \text{id}_X :
V \circ X \longrightarrow V \circ X
$$
sont homotopes. Cela résulte immédiatement du
lemme \ref{lemma-godement-two-maps} (avec
$G = G' = V$ et $F = F' = \text{id}_\mathcal{S}$).

\medskip\noindent
Démonstration de la seconde équivalence d'homotopie. D'après le
lemme \ref{lemma-godement-section}, pour construire une application
$h : U \to X \circ U$ inverse à droite de $\epsilon \star 1_U$, il
suffit de construire une application $h_0 : U \to Y \circ U = U \circ V \circ U$
telle que $1_U = (d \star 1_U) \circ h_0$. Bien entendu, nous prenons
$h_0 = 1_U \star \eta$ et l'égalité résulte de (\ref{equation-composition}).
Pour achever la démonstration, il faut montrer que les deux applications
$$
(\epsilon \star 1_U) \circ h, \text{id}_X \star 1_U :
X \circ U \longrightarrow X \circ U
$$
sont homotopes. Cela résulte immédiatement du
lemme \ref{lemma-godement-two-maps} (avec
$G = G' = \text{id}_\mathcal{A}$ et $F = F' = U$).
\end{proof}


\begin{example}
\label{example-module}
Soit $R$ un anneau. Comme exemple de ce qui précède, nous pouvons prendre
pour $i : \text{Mod}_R \to \textit{Sets}$ le foncteur d'oubli
et pour $F : \textit{Sets} \to \text{Mod}_R$ le foncteur qui associe
à un ensemble $E$ le $R$-module libre $R[E]$ de base $E$. Pour un $R$-module $M$,
le $R$-module simplicial $X(M)$ est de la forme suivante
$$
X(M) =
\left( \ldots
\xymatrix{
R[R[R[M]]]
\ar@<2ex>[r]
\ar@<0ex>[r]
\ar@<-2ex>[r]
&
R[R[M]]
\ar@<1ex>[r]
\ar@<-1ex>[r]
\ar@<1ex>[l]
\ar@<-1ex>[l]
&
R[M]
\ar@<0ex>[l]
}
\right)
$$
et est muni d'une augmentation vers $M$. Nous montrerons aussi que cette
augmentation est une équivalence d'homotopie sur les ensembles simpliciaux sous-jacents. D'après les lemmes
\ref{lemma-trivial-kan-homotopy}, \ref{lemma-homotopy-equivalence} et
\ref{lemma-qis-simplicial-abelian-groups}
cela revient à demander que $M$ soit le seul groupe d'homologie non nul
du complexe de chaînes associé au module simplicial $X(M)$.
\end{example}

\begin{example}
\label{example-polynomial-algebra}
Soit $A$ un anneau. Soit $\textit{Alg}_A$ la catégorie des
$A$-algèbres commutatives. Comme exemple de ce qui précède, nous pouvons prendre
pour $i : \textit{Alg}_A \to \textit{Sets}$ le foncteur d'oubli
et pour $F : \textit{Sets} \to \textit{Alg}_A$ le foncteur qui associe
à un ensemble $E$ l'algèbre de polynômes $A[E]$ en les indéterminées $E$,
à coefficients dans $A$.
(Nous nous excusons du conflit de notation entre cet exemple et
l'exemple \ref{example-module}.) Pour une $A$-algèbre $B$,
l'$A$-algèbre simpliciale $X(B)$ est de la forme suivante
$$
X(B) = \left( \ldots
\xymatrix{
A[A[A[B]]]
\ar@<2ex>[r]
\ar@<0ex>[r]
\ar@<-2ex>[r]
&
A[A[B]]
\ar@<1ex>[r]
\ar@<-1ex>[r]
\ar@<1ex>[l]
\ar@<-1ex>[l]
&
A[B]
\ar@<0ex>[l]
}
\right)
$$
et est munie d'une augmentation vers $B$. Nous montrerons aussi que cette
augmentation est une équivalence d'homotopie sur les ensembles simpliciaux sous-jacents. D'après les lemmes
\ref{lemma-trivial-kan-homotopy}, \ref{lemma-homotopy-equivalence} et
\ref{lemma-qis-simplicial-abelian-groups}
cela revient à demander que $B$ soit le seul groupe d'homologie non nul
du complexe de chaînes de $A$-modules
associé à $X(B)$ considéré comme un $A$-module simplicial.
\end{example}

\begin{example}
\label{example-module-maps}
Dans l'exemple \ref{example-module}, nous avons $X_n(M) = R[R[\ldots [M]\ldots]]$
avec $n + 1$ paires de crochets. Nous décrivons les applications construites ci-dessus à l'aide d'un
élément type
$$
\xi = \sum\nolimits_i r_i\left[\sum\nolimits_j r_{ij}[m_{ij}]\right]
$$
de $X_1(M)$. Les applications $d_0, d_1 : R[R[M]] \to R[M]$ sont données par
$$
d_0(\xi) = \sum\nolimits_{i, j} r_ir_{ij}[m_{ij}]
\quad\text{et}\quad
d_1(\xi) = \sum\nolimits_i r_i\left[\sum\nolimits_j r_{ij}m_{ij}\right].
$$
Les applications $s_0, s_1 : R[R[M]] \to R[R[R[M]]]$
sont données par
$$
s_0(\xi) =
\sum\nolimits_i r_i\left[\left[\sum\nolimits_j r_{ij}[m_{ij}]\right]\right]
\quad\text{et}\quad
s_1(\xi) = \sum\nolimits_i r_i\left[\sum\nolimits_j r_{ij}[[m_{ij}]]\right].
$$
\end{example}

\begin{example}
\label{example-polynomial-algebra-maps}
Dans l'exemple \ref{example-polynomial-algebra}, nous avons
$X_n(B) = A[A[\ldots [B]\ldots]]$ avec $n + 1$ paires de crochets.
Nous décrivons les applications construites ci-dessus à l'aide d'un élément type
$$
\xi = \sum\nolimits_i a_i [x_{i, 1}] \ldots [x_{i, m_i}] \in A[A[B]] = X_1(B)
$$
où, pour chaque $i, j$, nous pouvons écrire
$$
x_{i, j} = \sum a_{i, j, k} [b_{i, j, k, 1}] \ldots [b_{i, j, k, n_{i, j, k}}]
\in A[B]
$$
De toute évidence, cette notation est affreuse ! Pour l'alléger et déterminer
l'action des morphismes de $A$-algèbres $d_0, d_1 : A[A[B]] \to A[B]$, il suffit
d'examiner l'image des variables $[x]$, où
$$
x = \sum a_k [b_{k, 1}] \ldots [b_{k, n_k}] \in A[B]
$$
est un élément quelconque. On obtient alors
$$
d_0([x]) = x = \sum a_k [b_{k, 1}] \ldots [b_{k, n_k}]
\quad\text{et}\quad
d_1([x]) = \left[\sum a_k b_{k, 1} \ldots b_{k, n_k}\right]
$$
Les applications $s_0, s_1 : A[A[B]] \to A[A[A[B]]]$ sont données par
$$
s_0([x]) = \left[\left[\sum a_k [b_{k, 1}] \ldots [b_{k, n_k}]\right]\right]
\quad\text{et}\quad
s_1([x]) = \left[\sum a_k [[b_{k, 1}]] \ldots [[b_{k, n_k}]]\right]
$$
\end{example}

\begin{example}
\label{example-polynomial-algebra-homotopy}
Revenons à l'exemple \ref{example-polynomial-algebra}
étudié ci-dessus.
Notre lemme \ref{lemma-standard-simplicial-homotopy}
signifie que, pour tout morphisme d'anneaux $A \to B$, l'application d'anneaux simpliciaux
$$
\xymatrix{
A[A[A[B]]] \ar[d]
\ar@<2ex>[r]
\ar@<0ex>[r]
\ar@<-2ex>[r]
&
A[A[B]] \ar[d]
\ar@<1ex>[r]
\ar@<-1ex>[r]
\ar@<1ex>[l]
\ar@<-1ex>[l]
&
A[B] \ar[d]
\ar@<0ex>[l] \\
B
\ar@<2ex>[r]
\ar@<0ex>[r]
\ar@<-2ex>[r]
&
B
\ar@<1ex>[r]
\ar@<-1ex>[r]
\ar@<1ex>[l]
\ar@<-1ex>[l]
&
B
\ar@<0ex>[l]
}
$$
est une équivalence d'homotopie sur les ensembles simpliciaux sous-jacents. De plus, l'application
inverse construite dans le lemme \ref{lemma-standard-simplicial-homotopy}
est donnée en degré $n$ par
$$
b \longmapsto [\ldots[b]\ldots]
$$
avec les notations évidentes. Dans l'autre sens, le lemme nous dit que,
pour tout ensemble $E$, il existe une équivalence d'homotopie
$$
\xymatrix{
A[A[A[A[E]]]] \ar[d]
\ar@<2ex>[r]
\ar@<0ex>[r]
\ar@<-2ex>[r]
&
A[A[A[E]]] \ar[d]
\ar@<1ex>[r]
\ar@<-1ex>[r]
\ar@<1ex>[l]
\ar@<-1ex>[l]
&
A[A[E]] \ar[d]
\ar@<0ex>[l] \\
A[E]
\ar@<2ex>[r]
\ar@<0ex>[r]
\ar@<-2ex>[r]
&
A[E]
\ar@<1ex>[r]
\ar@<-1ex>[r]
\ar@<1ex>[l]
\ar@<-1ex>[l]
&
A[E]
\ar@<0ex>[l]
}
$$
d'anneaux simpliciaux. L'application inverse construite dans le lemme est donnée en degré $n$
par le morphisme d'anneaux
$$
\sum a_{e_1,\ldots,e_p}[e_1][e_2] \ldots [e_p]
\longmapsto
\sum a_{e_1,\ldots,e_p}[\ldots[e_1]\ldots][\ldots[e_2]\ldots]
\ldots [\ldots[e_p]\ldots]
$$
(avec les notations évidentes).
\end{example}










\begin{multicols}{2}[\section{Autres chapitres}]
\noindent
Préliminaires
\begin{enumerate}
\item \hyperref[introduction-section-phantom]{Introduction}
\item \hyperref[conventions-section-phantom]{Conventions}
\item \hyperref[sets-section-phantom]{Théorie des ensembles}
\item \hyperref[categories-section-phantom]{Catégories}
\item \hyperref[topology-section-phantom]{Topologie}
\item \hyperref[sheaves-section-phantom]{Faisceaux sur les espaces}
\item \hyperref[sites-section-phantom]{Sites et faisceaux}
\item \hyperref[stacks-section-phantom]{Champs}
\item \hyperref[fields-section-phantom]{Corps}
\item \hyperref[algebra-section-phantom]{Algèbre commutative}
\item \hyperref[brauer-section-phantom]{Groupes de Brauer}
\item \hyperref[homology-section-phantom]{Algèbre homologique}
\item \hyperref[derived-section-phantom]{Catégories dérivées}
\item \hyperref[simplicial-section-phantom]{Méthodes simpliciales}
\item \hyperref[more-algebra-section-phantom]{Compléments d'algèbre}
\item \hyperref[smoothing-section-phantom]{Lissage des morphismes d'anneaux}
\item \hyperref[modules-section-phantom]{Faisceaux de modules}
\item \hyperref[sites-modules-section-phantom]{Modules sur les sites}
\item \hyperref[injectives-section-phantom]{Objets injectifs}
\item \hyperref[cohomology-section-phantom]{Cohomologie des faisceaux}
\item \hyperref[sites-cohomology-section-phantom]{Cohomologie sur les sites}
\item \hyperref[dga-section-phantom]{Algèbre différentielle graduée}
\item \hyperref[dpa-section-phantom]{Algèbre à puissances divisées}
\item \hyperref[sdga-section-phantom]{Faisceaux différentiels gradués}
\item \hyperref[hypercovering-section-phantom]{Hyperrecouvrements}
\end{enumerate}
Schémas
\begin{enumerate}
\setcounter{enumi}{25}
\item \hyperref[schemes-section-phantom]{Schémas}
\item \hyperref[constructions-section-phantom]{Constructions de schémas}
\item \hyperref[properties-section-phantom]{Propriétés des schémas}
\item \hyperref[morphisms-section-phantom]{Morphismes de schémas}
\item \hyperref[coherent-section-phantom]{Cohomologie des schémas}
\item \hyperref[divisors-section-phantom]{Diviseurs}
\item \hyperref[limits-section-phantom]{Limites de schémas}
\item \hyperref[varieties-section-phantom]{Variétés}
\item \hyperref[topologies-section-phantom]{Topologies sur les schémas}
\item \hyperref[descent-section-phantom]{Descente}
\item \hyperref[perfect-section-phantom]{Catégories dérivées de schémas}
\item \hyperref[more-morphisms-section-phantom]{Compléments sur les morphismes}
\item \hyperref[flat-section-phantom]{Compléments sur la platitude}
\item \hyperref[groupoids-section-phantom]{Schémas en groupoïdes}
\item \hyperref[more-groupoids-section-phantom]{Compléments sur les schémas en groupoïdes}
\item \hyperref[etale-section-phantom]{Morphismes étales de schémas}
\end{enumerate}
Thèmes de théorie des schémas
\begin{enumerate}
\setcounter{enumi}{41}
\item \hyperref[chow-section-phantom]{Homologie de Chow}
\item \hyperref[intersection-section-phantom]{Théorie de l'intersection}
\item \hyperref[pic-section-phantom]{Schémas de Picard des courbes}
\item \hyperref[weil-section-phantom]{Théories cohomologiques de Weil}
\item \hyperref[adequate-section-phantom]{Modules adéquats}
\item \hyperref[dualizing-section-phantom]{Complexes dualisants}
\item \hyperref[duality-section-phantom]{Dualité pour les schémas}
\item \hyperref[discriminant-section-phantom]{Discriminants et différentes}
\item \hyperref[derham-section-phantom]{Cohomologie de de Rham}
\item \hyperref[local-cohomology-section-phantom]{Cohomologie locale}
\item \hyperref[algebraization-section-phantom]{Géométrie algébrique et formelle}
\item \hyperref[curves-section-phantom]{Courbes algébriques}
\item \hyperref[resolve-section-phantom]{Résolution des surfaces}
\item \hyperref[models-section-phantom]{Réduction semi-stable}
\item \hyperref[functors-section-phantom]{Foncteurs et morphismes}
\item \hyperref[equiv-section-phantom]{Catégories dérivées de variétés}
\item \hyperref[pione-section-phantom]{Groupes fondamentaux de schémas}
\item \hyperref[etale-cohomology-section-phantom]{Cohomologie étale}
\item \hyperref[crystalline-section-phantom]{Cohomologie cristalline}
\item \hyperref[proetale-section-phantom]{Cohomologie pro-étale}
\item \hyperref[relative-cycles-section-phantom]{Cycles relatifs}
\item \hyperref[more-etale-section-phantom]{Compléments de cohomologie étale}
\item \hyperref[trace-section-phantom]{La formule des traces}
\end{enumerate}
Espaces algébriques
\begin{enumerate}
\setcounter{enumi}{64}
\item \hyperref[spaces-section-phantom]{Espaces algébriques}
\item \hyperref[spaces-properties-section-phantom]{Propriétés des espaces algébriques}
\item \hyperref[spaces-morphisms-section-phantom]{Morphismes d'espaces algébriques}
\item \hyperref[decent-spaces-section-phantom]{Espaces algébriques décents}
\item \hyperref[spaces-cohomology-section-phantom]{Cohomologie des espaces algébriques}
\item \hyperref[spaces-limits-section-phantom]{Limites d'espaces algébriques}
\item \hyperref[spaces-divisors-section-phantom]{Diviseurs sur les espaces algébriques}
\item \hyperref[spaces-over-fields-section-phantom]{Espaces algébriques sur des corps}
\item \hyperref[spaces-topologies-section-phantom]{Topologies sur les espaces algébriques}
\item \hyperref[spaces-descent-section-phantom]{Descente et espaces algébriques}
\item \hyperref[spaces-perfect-section-phantom]{Catégories dérivées d'espaces}
\item \hyperref[spaces-more-morphisms-section-phantom]{Compléments sur les morphismes d'espaces}
\item \hyperref[spaces-flat-section-phantom]{Platitude sur les espaces algébriques}
\item \hyperref[spaces-groupoids-section-phantom]{Groupoïdes dans les espaces algébriques}
\item \hyperref[spaces-more-groupoids-section-phantom]{Compléments sur les groupoïdes dans les espaces}
\item \hyperref[bootstrap-section-phantom]{Amorçage}
\item \hyperref[spaces-pushouts-section-phantom]{Sommes amalgamées d'espaces algébriques}
\end{enumerate}
Thèmes de géométrie
\begin{enumerate}
\setcounter{enumi}{81}
\item \hyperref[spaces-chow-section-phantom]{Groupes de Chow des espaces}
\item \hyperref[groupoids-quotients-section-phantom]{Quotients de groupoïdes}
\item \hyperref[spaces-more-cohomology-section-phantom]{Compléments sur la cohomologie des espaces}
\item \hyperref[spaces-simplicial-section-phantom]{Espaces simpliciaux}
\item \hyperref[spaces-duality-section-phantom]{Dualité pour les espaces}
\item \hyperref[formal-spaces-section-phantom]{Espaces algébriques formels}
\item \hyperref[restricted-section-phantom]{Algébrisation des espaces formels}
\item \hyperref[spaces-resolve-section-phantom]{Résolution des surfaces revisitée}
\end{enumerate}
Théorie des déformations
\begin{enumerate}
\setcounter{enumi}{89}
\item \hyperref[formal-defos-section-phantom]{Théorie formelle des déformations}
\item \hyperref[defos-section-phantom]{Théorie des déformations}
\item \hyperref[cotangent-section-phantom]{Le complexe cotangent}
\item \hyperref[examples-defos-section-phantom]{Problèmes de déformation}
\end{enumerate}
Champs algébriques
\begin{enumerate}
\setcounter{enumi}{93}
\item \hyperref[algebraic-section-phantom]{Champs algébriques}
\item \hyperref[examples-stacks-section-phantom]{Exemples de champs}
\item \hyperref[stacks-sheaves-section-phantom]{Faisceaux sur les champs algébriques}
\item \hyperref[criteria-section-phantom]{Critères de représentabilité}
\item \hyperref[artin-section-phantom]{Axiomes d'Artin}
\item \hyperref[quot-section-phantom]{Espaces de Quot et de Hilbert}
\item \hyperref[stacks-properties-section-phantom]{Propriétés des champs algébriques}
\item \hyperref[stacks-morphisms-section-phantom]{Morphismes de champs algébriques}
\item \hyperref[stacks-limits-section-phantom]{Limites de champs algébriques}
\item \hyperref[stacks-cohomology-section-phantom]{Cohomologie des champs algébriques}
\item \hyperref[stacks-perfect-section-phantom]{Catégories dérivées de champs}
\item \hyperref[stacks-introduction-section-phantom]{Introduction aux champs algébriques}
\item \hyperref[stacks-more-morphisms-section-phantom]{Compléments sur les morphismes de champs}
\item \hyperref[stacks-geometry-section-phantom]{La géométrie des champs}
\end{enumerate}
Thèmes de théorie des espaces de modules
\begin{enumerate}
\setcounter{enumi}{107}
\item \hyperref[moduli-section-phantom]{Champs de modules}
\item \hyperref[moduli-curves-section-phantom]{Espaces de modules de courbes}
\end{enumerate}
Divers
\begin{enumerate}
\setcounter{enumi}{109}
\item \hyperref[examples-section-phantom]{Exemples}
\item \hyperref[exercises-section-phantom]{Exercices}
\item \hyperref[guide-section-phantom]{Guide bibliographique}
\item \hyperref[desirables-section-phantom]{Desiderata}
\item \hyperref[coding-section-phantom]{Style de programmation}
\item \hyperref[obsolete-section-phantom]{Obsolète}
\item \hyperref[fdl-section-phantom]{Licence de documentation libre GNU}
\item \hyperref[index-section-phantom]{Index généré automatiquement}
\end{enumerate}
\end{multicols}


\bibliography{my}
\bibliographystyle{amsalpha}

\end{document}
